\documentclass[11pt, a4paper]{article}
\usepackage[a4paper, margin=2.2cm]{geometry}
\usepackage{amsmath, amssymb}
\usepackage{enumitem}
\usepackage{booktabs}
\usepackage{hyperref}

\hypersetup{
  colorlinks=true,
  linkcolor=blue,
  pdftitle={TIFR GS Mathematics Question Papers (2010--2025)}
}

\title{\textbf{TIFR Graduate School Admissions (GS)}\\Mathematics Question Papers (2010--2025)}
\author{Tata Institute of Fundamental Research}
\date{}

\begin{document}

\maketitle
\tableofcontents
\newpage

% ==============================================================================
\section{GS-2010 (Mathematics)}
% ==============================================================================

\begin{center}
  \textbf{TATA INSTITUTE OF FUNDAMENTAL RESEARCH}\\
  \textbf{Written Test in MATHEMATICS --- December 13, 2009}\\
  \textbf{Duration:} Two hours (2 hours)
\end{center}

\subsection*{Instructions}
\begin{enumerate}
  \item Write your FULL NAME and REFERENCE CODE in block letters on this page and also fill-in all details on the ANSWER SHEET.
  \item The Answer Sheet is machine-readable. Please read the instructions on the reverse of the answer sheet before you start filling it up. Only use HB pencils to fill-in the answer sheet.
  \item There are thirty (30) questions divided into TWO parts (Part A \& B) of 15 questions each. If you wish to be considered ONLY for the Integrated Ph.D.\ programme at Bengaluru, you need to answer only Part A. All other candidates need to answer both Parts.
  \item Indicate your ANSWER ON THE ANSWER SHEET by blackening the appropriate circle for each question. Each correct answer will get $1$ mark; each wrong answer will get a $-1$ mark, and a question not answered will not get you any mark. Do not mark more than one circle for any question; this will be treated as a wrong answer.
  \item We advise you to first mark the correct answers in the QUESTION SHEET and then TRANSFER these to the ANSWER SHEET only when you are sure of your choice.
  \item Rough work may be done on blank pages of the question paper. If needed, you may ask for extra rough sheets from an Invigilator. Please do not scribble or do rough work on the reverse of your hall ticket. If found, the hall ticket will be retained.
  \item Use of calculators is NOT permitted.
  \item Do NOT ask for clarifications from the invigilators regarding the questions. They have been instructed not to respond to any such inquiries from candidates. In case a correction/clarification is deemed necessary, the invigilator(s) will announce it publicly.
  \item This set of question paper must be returned along with your answer sheet and all extra rough sheets used.
  \item See below for Notation and Conventions used in this test.
\end{enumerate}

\subsection*{Notation and Conventions}
\begin{itemize}
  \item $\mathbb{Z} =$ set of integers; $\mathbb{N} =$ set of natural numbers; $\mathbb{Q} =$ set of rational numbers; $\mathbb{R} =$ set of real numbers; $\mathbb{C} =$ set of complex numbers.
  \item $\mathbb{R}^n =$ Euclidean space of dimension $n$.
  \item $[a, b] = \{x \in \mathbb{R} : a \le x \le b\}$ for real numbers $a$ and $b$ with $a < b$.
  \item For a differentiable function $f: \mathbb{R} \to \mathbb{R}$, $f'$ denotes its derivative.
  \item For any natural number $n$, $\mathbb{Z}/n\mathbb{Z}$ denotes the ring of integers modulo $n$.
  \item Subsets of $\mathbb{R}^n$ are assumed to carry the induced topology and metric.
\end{itemize}

\subsection*{Part A}
\begin{enumerate}
  \item A cyclic group of order $60$ has
  \begin{enumerate}[label=(\alph*)]
    \item $12$ generators.
    \item $15$ generators.
    \item $16$ generators.
    \item $20$ generators.
  \end{enumerate}

  \item Which of the following is false?
  \begin{enumerate}[label=(\alph*)]
    \item Any abelian group of order $27$ is cyclic.
    \item Any abelian group of order $14$ is cyclic.
    \item Any abelian group of order $21$ is cyclic.
    \item Any abelian group of order $30$ is cyclic.
  \end{enumerate}

  \item The last digit of $2^{80}$ is
  \begin{enumerate}[label=(\alph*)]
    \item $2$
    \item $4$
    \item $6$
    \item $8$
  \end{enumerate}

  \item The sum of the series $\displaystyle \frac{1}{1 \cdot 2} + \frac{1}{2 \cdot 3} + \frac{1}{3 \cdot 4} + \dots + \frac{1}{100 \cdot 101}$ is
  \begin{enumerate}[label=(\alph*)]
    \item $\frac{99}{101}$
    \item $\frac{98}{101}$
    \item $\frac{99}{100}$
    \item None of the above.
  \end{enumerate}

  \item Let $f$ be a one to one function from the closed interval $[0, 1]$ to the set of real numbers $\mathbb{R}$, then
  \begin{enumerate}[label=(\alph*)]
    \item $f$ must be onto.
    \item range of $f$ must contain a rational number.
    \item range of $f$ must contain an irrational number.
    \item range of $f$ must contain both rational and irrational numbers.
  \end{enumerate}

  \item The maximum value of $f(x) = x^n(1 - x)^n$ for a natural number $n \ge 1$ and $0 \le x \le 1$ is
  \begin{enumerate}[label=(\alph*)]
    \item $\frac{1}{2^n}$
    \item $\frac{1}{3^n}$
    \item $\frac{1}{5^n}$
    \item $\frac{1}{4^n}$
  \end{enumerate}

  \item The sequence $\sqrt{7}, \sqrt{7 + \sqrt{7}}, \sqrt{7 + \sqrt{7 + \sqrt{7}}}, \dots$ converges to
  \begin{enumerate}[label=(\alph*)]
    \item $\frac{1 + \sqrt{33}}{2}$
    \item $\frac{1 + \sqrt{32}}{2}$
    \item $\frac{1 + \sqrt{30}}{2}$
    \item $\frac{1 + \sqrt{29}}{2}$
  \end{enumerate}

  \item Let $f(x) = |x|^{3/2}$, $x \in \mathbb{R}$. Then
  \begin{enumerate}[label=(\alph*)]
    \item $f$ is uniformly continuous.
    \item $f$ is continuous, but not differentiable at $x = 0$.
    \item $f$ is differentiable and $f'$ is continuous.
    \item $f$ is differentiable, but $f'$ is discontinuous at $x = 0$.
  \end{enumerate}

  \item The total number of subsets of a set of $6$ elements is
  \begin{enumerate}[label=(\alph*)]
    \item $720$
    \item $6^6$
    \item $21$
    \item None of the above.
  \end{enumerate}

  \item Let $M_n(\mathbb{R})$ be the set of $n \times n$-matrices with real entries. Which of the following statements is true?
  \begin{enumerate}[label=(\alph*)]
    \item Any matrix $A \in M_4(\mathbb{R})$ has a real eigenvalue.
    \item Any matrix $A \in M_5(\mathbb{R})$ has a real eigenvalue.
    \item Any matrix $A \in M_2(\mathbb{R})$ has a real eigenvalue.
    \item None of the above.
  \end{enumerate}

  \item The series $\displaystyle \sum_{n=1}^\infty \frac{(-1)^{n+1}}{\sqrt{n}}$
  \begin{enumerate}[label=(\alph*)]
    \item converges but not absolutely.
    \item converges absolutely.
    \item diverges.
    \item none of the above.
  \end{enumerate}

  \item The sum of the roots of the equation $x^5 + 3x^2 + 7 = 0$ is
  \begin{enumerate}[label=(\alph*)]
    \item $-3$
    \item $\frac{3}{7}$
    \item $\frac{-1}{7}$
    \item $0$
  \end{enumerate}

  \item $\displaystyle \lim_{x \to 0} x \sin\left(\frac{1}{x}\right)$ is
  \begin{enumerate}[label=(\alph*)]
    \item $1$
    \item $0$
    \item $\frac{1}{2}$
    \item does not exist.
  \end{enumerate}

  \item The solution of the ordinary differential equation $\displaystyle \frac{dy}{dx} = y$, $y(0) = 0$
  \begin{enumerate}[label=(\alph*)]
    \item {[obscured in source]}
    \item {[obscured in source]}
    \item {[obscured in source]}
    \item is zero.
  \end{enumerate}

  \item Let $G$ be the set of all $2 \times 2$ symmetric, invertible matrices with real entries. Then with matrix multiplication, $G$ is
  \begin{enumerate}[label=(\alph*)]
    \item an infinite group.
    \item a finite group.
    \item not a group.
    \item an abelian group.
  \end{enumerate}
\end{enumerate}

\subsection*{Part B}
\begin{enumerate}
  \item Let $u_n = \sin\left(\frac{\pi}{n}\right)$ and consider the series $\sum u_n$. Which of the following statements is false?
  \begin{enumerate}[label=(\alph*)]
    \item $\sum u_n$ is convergent.
    \item $u_n \to 0$ as $n \to \infty$.
    \item $\sum u_n$ is divergent.
    \item $\sum u_n$ is absolutely convergent.
  \end{enumerate}

  \item If $V$ is a vector space over the field $\mathbb{Z}/5\mathbb{Z}$ and $\dim_{\mathbb{Z}/5\mathbb{Z}}(V) = 3$, then $V$ has
  \begin{enumerate}[label=(\alph*)]
    \item $125$ elements.
    \item $15$ elements.
    \item $243$ elements.
    \item None of the above.
  \end{enumerate}

  \item If $f, g: \mathbb{R} \to \mathbb{R}$ are uniformly continuous functions, then their composition $g \circ f$ is
  \begin{enumerate}[label=(\alph*)]
    \item uniformly continuous.
    \item continuous but not uniformly continuous.
    \item continuous and bounded.
    \item None of the above.
  \end{enumerate}

  \item Which of the following statements is false?
  \begin{enumerate}[label=(\alph*)]
    \item There exists a natural number which when divided by $3$ leaves remainder $1$ and which when divided by $4$ leaves remainder $0$.
    \item There exists a natural number which when divided by $6$ leaves remainder $2$ and when divided by $9$ leaves remainder $1$.
    \item There exists a natural number which when divided by $7$ leaves remainder $1$ and when divided by $11$ leaves remainder $3$.
    \item There exists a natural number which when divided by $12$ leaves remainder $7$ and when divided by $8$ leaves remainder $3$.
  \end{enumerate}

  \item If $f_n(x)$ are continuous functions from $[0, 1]$ to $[0, 1]$, and $f_n(x) \to f(x)$ as $n \to \infty$, then which of the following statements is true?
  \begin{enumerate}[label=(\alph*)]
    \item $f_n(x)$ converges to $f(x)$ uniformly on $[0, 1]$.
    \item $f_n(x)$ converges to $f(x)$ uniformly on $(0, 1)$.
    \item $f(x)$ is continuous on $[0, 1]$.
    \item None of the above.
  \end{enumerate}

  \item Let $A, B$ be subsets of $\mathbb{R}$. Define $A + B$ to be the set of all sums $x + y$ with $x \in A$ and $y \in B$. Which of the following statements is false?
  \begin{enumerate}[label=(\alph*)]
    \item If $A$ and $B$ are bounded, then $A + B$ is bounded.
    \item If $A$ and $B$ are open, then $A + B$ is open.
    \item If $A$ and $B$ are closed, then $A + B$ is closed.
    \item If $A$ and $B$ are connected, then $A + B$ is connected.
  \end{enumerate}

  \item Number of solutions of the ordinary differential equation $\displaystyle \frac{d^2 y}{dx^2} - y = 0$; $y(0) = 0$, $y(\pi) = 1$
  \begin{enumerate}[label=(\alph*)]
    \item is $0$.
    \item is $1$.
    \item is $2$.
    \item None of the above.
  \end{enumerate}

  \item The function $f(x)$ defined by
  \[
    f(x) = \begin{cases} 0 & \text{if } x \text{ is rational,} \\ x & \text{if } x \text{ is irrational} \end{cases}
  \]
  \begin{enumerate}[label=(\alph*)]
    \item is not continuous at any point.
    \item is continuous at every point.
    \item is continuous at every rational number.
    \item is continuous at $x = 0$.
  \end{enumerate}

  \item Let $G = \{z \in \mathbb{C} \mid z^n = 1 \text{ for some positive integer } n\}$. Then under multiplication of complex numbers,
  \begin{enumerate}[label=(\alph*)]
    \item $G$ is a group of finite order.
    \item $G$ is a group of infinite order, but every element of $G$ has finite order.
    \item $G$ is a cyclic group.
    \item None of the above.
  \end{enumerate}

  \item Let $x$ and $y$ in $\mathbb{R}^n$ be non-zero column vectors. Form the matrix $A = x y^t$ where $y^t$ is the transpose of $y$. Then the rank of $A$ is
  \begin{enumerate}[label=(\alph*)]
    \item $2$
    \item $0$
    \item at least $n/2$.
    \item None of the above.
  \end{enumerate}

  \item Which of the following is true?
  \begin{enumerate}[label=(\alph*)]
    \item The matrix $\begin{pmatrix} 1 & 0 \\ 1 & 2 \end{pmatrix}$ is not diagonalisable.
    \item The matrix $\begin{pmatrix} 1 & 5 \\ 0 & 2 \end{pmatrix}$ is diagonalisable.
    \item The matrix $\begin{pmatrix} 1 & 1 \\ 0 & 1 \end{pmatrix}$ is diagonalisable.
    \item None of the above.
  \end{enumerate}

  \item If $n$ and $m$ are positive integers and $n^9 = 19m + r$, then the possible values for $r$ modulo $19$ are
  \begin{enumerate}[label=(\alph*)]
    \item only $0$.
    \item only $0, \pm 1$.
    \item only $\pm 1$.
    \item None of the above.
  \end{enumerate}

  \item Define $\{x_n\}$ as $x_1 = 0.1$, $x_2 = 0.101$, $x_3 = 0.101001, \dots$. Then the sequence $\{x_n\}$
  \begin{enumerate}[label=(\alph*)]
    \item converges to a rational number.
    \item converges to an irrational number.
    \item does not converge.
    \item oscillates.
  \end{enumerate}

  \item The equations
  \begin{align*}
    x_1 + 2x_2 + 3x_3 &= 1 \\
    x_1 + 4x_2 + 9x_3 &= 1 \\
    x_1 + 8x_2 + 27x_3 &= 1
  \end{align*}
  have
  \begin{enumerate}[label=(\alph*)]
    \item only one solution.
    \item two solutions.
    \item infinitely many solutions.
    \item no solution.
  \end{enumerate}

  \item Which of the following statements is false?
  \begin{enumerate}[label=(\alph*)]
    \item The polynomial $x^2 + x + 1$ is irreducible in $\mathbb{Z}/2\mathbb{Z}[x]$.
    \item The polynomial $x^2 - 2$ is irreducible in $\mathbb{Q}[x]$.
    \item The polynomial $x^2 + 1$ is reducible in $\mathbb{Z}/5\mathbb{Z}[x]$.
    \item The polynomial $x^2 + 1$ is reducible in $\mathbb{Z}/7\mathbb{Z}[x]$.
  \end{enumerate}
\end{enumerate}

\newpage

% ==============================================================================
\section{GS-2011 (Mathematics)}
% ==============================================================================

\subsection*{Notation and Conventions}
\begin{itemize}
  \item $\mathbb{Z} =$ set of integers; $\mathbb{N} =$ set of natural numbers; $\mathbb{Q} =$ set of rational numbers; $\mathbb{R} =$ set of real numbers; $\mathbb{C} =$ set of complex numbers.
  \item $\mathbb{R}^n =$ Euclidean space of dimension $n$.
  \item For a natural number $n$, the product of all the natural numbers from $1$ up to $n$ is denoted by $n!$.
  \item $[a, b] = \{x \in \mathbb{R} : a \le x \le b\}$ for real numbers $a$ and $b$ with $a < b$.
  \item $(a, b) = \{x \in \mathbb{R} : a < x < b\}$ for real numbers $a$ and $b$ with $a < b$.
  \item For a differentiable function $f: \mathbb{R} \to \mathbb{R}$, $f'$ denotes its derivative.
  \item For any natural number $n$, $\mathbb{Z}/n\mathbb{Z}$ denotes the ring of integers modulo $n$.
  \item Subsets of $\mathbb{R}^n$ are assumed to carry the induced topology and metric.
\end{itemize}

\subsection*{Part A}
\begin{enumerate}
  \item Consider the sequence $\{x_n\}$ defined by $x_n = \frac{[nx]}{n}$ for $x \in \mathbb{R}$ where $[\,]$ denotes the integer part. Then $\{x_n\}$
  \begin{enumerate}[label=(\alph*)]
    \item converges to $x$.
    \item converges but not to $x$.
    \item does not converge.
    \item oscillates.
  \end{enumerate}

  \item $\displaystyle \lim_{x \to 0} x \sin(1/x^2)$ equals
  \begin{enumerate}[label=(\alph*)]
    \item $1$.
    \item $0$.
    \item $\infty$.
    \item oscillates.
  \end{enumerate}

  \item Let $A$ be a $5 \times 5$ matrix with real entries, then $A$ has
  \begin{enumerate}[label=(\alph*)]
    \item an eigenvalue which is purely imaginary.
    \item at least one real eigenvalue.
    \item at least two eigenvalues which are not real.
    \item at least $2$ distinct real eigenvalues.
  \end{enumerate}

  \item The groups $\mathbb{Z}_9$ and $\mathbb{Z}_3 \times \mathbb{Z}_3$ are
  \begin{enumerate}[label=(\alph*)]
    \item isomorphic
    \item abelian
    \item non abelian
    \item cyclic
  \end{enumerate}

  \item The differential equation $\displaystyle \frac{dy}{dx} = y^{1/3}$, $y(0) = 0$ has
  \begin{enumerate}[label=(\alph*)]
    \item a unique solution.
    \item no nontrivial solution.
    \item finite number of solutions.
    \item infinite number of solutions.
  \end{enumerate}

  \item The function $f_n(x) = n \sin(x/n)$
  \begin{enumerate}[label=(\alph*)]
    \item does not converge for any $x$ as $n \to \infty$.
    \item converges to the constant function $1$ as $n \to \infty$.
    \item converges to the function $x$ as $n \to \infty$.
    \item does not converge for all $x$ as $n \to \infty$.
  \end{enumerate}

  \item The equation $x^{22} \equiv 2 \pmod{23}$ has
  \begin{enumerate}[label=(\alph*)]
    \item no solutions.
    \item $23$ solutions.
    \item exactly one solution.
    \item $22$ solutions.
  \end{enumerate}

  \item The sum of the squares of the roots of the cubic equation $x^3 - 4x^2 + 6x + 1 = 0$ is
  \begin{enumerate}[label=(\alph*)]
    \item $0$.
    \item $4$.
    \item $16$.
    \item none of the above.
  \end{enumerate}

  \item The function $f(x)$ defined by
  \[
    f(x) = \begin{cases} ax + b, & x \ge 1, \\ x^2 + 3x + 3, & x \le 1 \end{cases}
  \]
  is {[obscured in source]}
  \begin{enumerate}[label=(\alph*)]
    \item for a unique value of $a$ and infinitely many values of $b$.
    \item for a unique value of $b$ and infinitely many values of $a$.
    \item for infinitely many values of $a$ and $b$.
    \item none of the above.
  \end{enumerate}

  \item Let $m \le n$ be natural numbers. The number of injective maps from a set of cardinality $m$ to a set of cardinality $n$ is
  \begin{enumerate}[label=(\alph*)]
    \item $m!$
    \item $n!$
    \item $(n - m)!$
    \item none of the above.
  \end{enumerate}

  \item \emph{(True/False)} For any real number $c$, the polynomial $x^3 + x + c$ has exactly one real root.
  \item \emph{(True/False)} $e^{\sqrt{2}} > 3$.
  \item \emph{(True/False)} $A$ is a $3 \times 4$-matrix of rank $3$. Then the system of equations $Ax = b$ has exactly one solution.
  \item \emph{(True/False)} $\log x$ is uniformly continuous on $(\text{[obscured]}, \infty)$.
  \item \emph{(True/False)} If $A, B$ are closed subsets of $[0, \infty)$, then $A + B = \{x + y \mid x \in A, \, y \in B\}$ is closed in $[0, \infty)$.
  \item \emph{(True/False)} The polynomial $x^4 + 7x^3 - 13x^2 + 11x$ has exactly one real root.
  \item \emph{(True/False)} The value of the infinite product $\displaystyle \prod_{n=2}^\infty \left(1 - \frac{1}{n^2}\right)$ is $1$.
  \item \emph{(True/False)} Consider the map $T$ from the vector space of polynomials of degree at most $5$ over the reals to $\mathbb{R} \times \mathbb{R}$, given by sending a polynomial $P$ to the pair $(P(3), P'(3))$ where $P'$ is the derivative of $P$. Then the dimension of the kernel is $3$.
  \item \emph{(True/False)} The derivative of the function $\displaystyle \int_0^{\sqrt{x}} e^{-t^2}\,dt$ at $x = 1$ is $e^{-1}$.
  \item \emph{(True/False)} The equation $63x + 70y + 15z = 2010$ has an integral solution.
  \item \emph{(True/False)} Any continuous function from the open unit interval $(0, 1)$ to itself has a fixed point.
  \item \emph{(True/False)} There exists a group with a proper subgroup isomorphic to itself.
  \item \emph{(True/False)} The space of solutions of infinitely differentiable functions satisfying the equation $y'' + y = 0$ is infinite dimensional.
  \item \emph{(True/False)} The series $\displaystyle \sum_{n=1}^\infty \frac{\sqrt{n+1} - \sqrt{n}}{n}$ diverges.
  \item \emph{(True/False)} The function
  \[
    f(x) = \begin{cases} 0 & \text{if } x \text{ is rational,} \\ x & \text{if } x \text{ is irrational} \end{cases}
  \]
  is not continuous anywhere on the real line.
\end{enumerate}

\subsection*{Part B (True/False)}
\begin{enumerate}
  \item Let $A$ be a $2 \times 2$-matrix with complex entries. The number of $2 \times 2$-matrices $A$ with complex entries satisfying the equation $A^3 = A$ is infinite.
  \item In the ring $\mathbb{Z}/8\mathbb{Z}$ the equation $x^2 = 1$ has exactly $2$ solutions.
  \item There are $n$ homomorphisms from the group $\mathbb{Z}/n\mathbb{Z}$ to the additive group of rationals.
  \item A bounded continuous function on $\mathbb{R}$ is uniformly continuous.
  \item The symmetric group $S_5$ consisting of permutations on $5$ symbols has an element of order $6$.
  \item Suppose $f_n(x)$ is a sequence of {[obscured in source]} on the closed interval $[0, 1]$ converging to $0$ pointwise. {[obscured in source]} converges to $0$.
  \item {[obscured in source]} from $S_3$ to $\mathbb{Z}/3\mathbb{Z}$.
  \item If $A$ and $B$ are $3 \times 3$ matrices and $A$ is invertible, then there exists an integer $n$ such that $A + nB$ is invertible.
  \item Let $P$ be a degree $3$ polynomial with complex coefficients such that the constant term is $2010$. Then $P$ has a root $\alpha$ with $|\alpha| > 10$.
  \item Suppose a box contains three cards, one with both sides white, one with both sides black, and one with one side white and the other side black. If you pick a card at random, and the side facing you is white, then the probability that the other side is white is $1/2$.
  \item There exists a set $A \subset \{1, 2, \dots, 100\}$ with $65$ elements, such that $65$ cannot be expressed as a sum of two elements in $A$.
  \item Let $S$ be a finite subset of $\mathbb{R}^3$ such that any three elements in $S$ span a two dimensional subspace. Then $S$ spans a two dimensional space.
  \item Any non-singular $k \times k$-matrix with real entries can be made singular by changing exactly one entry.
  \item Let $f$ be a continuous integrable function of $\mathbb{R}$ such that either $f(x) > 0$ or $f(x) + f(x + 1) > 0$ for all $x \in \mathbb{R}$. Then $\int_{-\infty}^\infty f(x)\,dx > 0$.
  \item A gardener throws $18$ seeds onto an {[obscured in source]} two seeds are within a distance of $25$ centimetres.
\end{enumerate}

\newpage

% ==============================================================================
\section{GS-2012 (Mathematics)}
% ==============================================================================

\begin{center}
  \textbf{TATA INSTITUTE OF FUNDAMENTAL RESEARCH}\\
  \textbf{Written Test in MATHEMATICS --- December 11, 2011}\\
  \textbf{Duration:} Two hours (2 hours)
\end{center}

\subsection*{Instructions}
There are 4 parts and 40 True/False questions in total, consisting of 10 questions in each part (Part A, Part B, Part C, and Part D). Every correct answer carries $+1$ mark and every wrong answer carries $-1$ mark.

\subsection*{Notation and Conventions}
\begin{itemize}
  \item $\mathbb{N} :=$ Set of natural numbers; $\mathbb{Z} :=$ Set of integers; $\mathbb{Q} :=$ Set of rational numbers; $\mathbb{R} :=$ Set of real numbers; $\mathbb{C} :=$ Set of complex numbers.
  \item $\mathbb{R}^n :=$ $n$-dimensional vector space over $\mathbb{R}$.
  \item $(a, b) := \{x \in \mathbb{R} \mid a < x < b\}$.
  \item For a differentiable real valued function $f: \mathbb{R} \to \mathbb{R}$, $f'$ denotes its derivative and $f^{(k)}$ means the $k^{\text{th}}$ derivative.
  \item Subsets of $\mathbb{R}^n$ are assumed to carry the induced topology and the metric.
\end{itemize}

\subsection*{Part A (True/False)}
\begin{enumerate}
  \item If $H_1$ \& $H_2$ are subgroups of a group $G$ then $H_1 \cdot H_2 = \{h_1 h_2 \in G \mid h_1 \in H_1, \, h_2 \in H_2\}$ is a subgroup of $G$.
  \item There exist polynomials $f(x)$ and $g(x)$, with complex coefficients, such that $\left(\frac{f(x)}{g(x)}\right)^2 = x$.
  \item Let $f$ be real valued, differentiable on $(a, b)$ and $f'(x) \ne 0$ for all $x \in (a, b)$. Then $f$ is $1$--$1$.
  \item The inequality $\displaystyle \sum_{n=0}^\infty \frac{(\log \log 2)^n}{n!} > \frac{3}{5}$ holds.
  \item Every subgroup of order $74$ in a group of order $148$ is normal.
  \item Let $u_1, u_2, u_3, u_4$ be vectors in $\mathbb{R}^2$ and $u = \sum_{j=1}^4 t_j u_j$; $t_j > 0$ and $\sum_{j=1}^4 t_j = 1$. Then three vectors $v_1, v_2, v_3 \in \mathbb{R}^2$ may be chosen from $\{u_1, u_2, u_3, u_4\}$ such that $u = \sum_{j=1}^3 s_j v_j$, $s_j \ge 0$, $\sum_{j=1}^3 s_j = 1$.
  \item The inequality $\sqrt{1 + x} < 1 + x/2$ for $x \in (-1, 10)$ is true.
  \item If $n$ is not a multiple of $23$ then the remainder when $n^{11}$ is divided by $23$ is $\pm 1 \pmod{23}$.
  \item Suppose $A$ is a nilpotent matrix and $I$ is the identity matrix. Then $(I + A)$ is invertible.
  \item The equations
  \begin{align*}
    x_1 + \frac{1}{2} x_2 + \frac{1}{3} x_3 &= 1 \\
    x_1 + \frac{1}{4} x_2 + \frac{1}{9} x_3 &= 1 \\
    x_1 + \frac{1}{8} x_2 + \frac{1}{27} x_3 &= 1
  \end{align*}
  have no solution.
\end{enumerate}

\subsection*{Part B (True/False)}
\begin{enumerate}[resume]
  \item The automorphism group $\operatorname{Aut}(\mathbb{Z}/2 \times \mathbb{Z}/2)$ is abelian.
  \item Let $V$ be the vector space consisting of polynomials of $\mathbb{R}[t]$ of $\deg \le 2$. The map $T: V \to V$ sending $f(t)$ to $f(t) + f'(t)$ is invertible.
  \item The polynomials $(t - 1)(t - 2)$, $(t - 2)(t - 3)$, $(t - 3)(t - 4)$, $(t - 4)(t - 6) \in \mathbb{R}[t]$ are linearly independent.
  \item $A \in M_2(\mathbb{C})$ and $A$ is nilpotent then $A^2 = 0$.
  \item Let $P$ be an $n \times n$ matrix whose row sums equal $1$. Then for any positive integer $m$ the row sums of the matrix $P^m$ equal $1$.
  \item There is a non trivial group homomorphism from $\mathbb{C}$ to $\mathbb{R}$.
  \item If the equation {[obscured in source]} holds in a {[obscured in source]} $yzx = 1$.
  \item Any $3 \times 3$ and $5 \times 5$ skew-symmetric matrices have always zero determinants.
  \item The rank of the matrix
  \[
    \begin{pmatrix}
      11 & 12 & 13 & 14 \\
      21 & 22 & 23 & 24 \\
      31 & 32 & 33 & 34 \\
      41 & 42 & 43 & 44
    \end{pmatrix}
  \]
  is $2$.
  \item The number $2$ is a prime in $\mathbb{Z}[i]$.
\end{enumerate}

\subsection*{Part C (True/False)}
\begin{enumerate}[resume]
  \item Let $f: \mathbb{R}^2 \to \mathbb{R}$ be a continuous function. Then the derivative $\frac{\partial^2 f}{\partial x \partial y}$ can exist without $\frac{\partial f}{\partial x}$ existing.
  \item If $f$ is continuous on $[0, 1]$ and if $\int_0^1 f(x) x^n\,dx = 0$ for $n = 0, 1, 2, 3, \dots$, then $\int_0^1 f^2(x)\,dx = 0$.
  \item Suppose that $f \in \mathfrak{L}^2(\mathbb{R})$. Then $f \in \mathfrak{L}^1(\mathbb{R})$.
  \item The integral $\displaystyle \int_{-\infty}^{+\infty} \frac{e^{-x}}{1 + x^2}\,dx$ is convergent.
  \item If $A \subset \mathbb{R}$ and open then the interior of the closure $\overline{A}^\circ$ is $A$.
  \item If $f \in C^\infty$ and $f^{(k)}(0) = 0$ for all integer $k \ge 0$, then $f \equiv 0$.
  \item Let $f: [0, 1] \to [0, 1]$ be continuous then $f$ assumes the value $\int_0^1 f^2(t)\,dt$ somewhere in $[0, 1]$.
  \item Let $f: \mathbb{R} \to \mathbb{R}$ {[obscured in source]} that $\lim_{h \to 0} \frac{f(x + h) - f(x - h)}{h}$ exists for all $x \in \mathbb{R}$. Then $f$ is differentiable in $\mathbb{R}$.
  \item The functions $f(x) = x|x|$ and $x|\sin x|$ are not differentiable at $x = 0$.
  \item The composition of two uniformly continuous functions need not always be uniformly continuous.
\end{enumerate}

\subsection*{Part D (True/False)}
\begin{enumerate}[resume]
  \item $f: [0, \infty] \to [0, \infty]$ is continuous and bounded then $f$ has a fixed point.
  \item The polynomial $X^8 + 1$ is irreducible in $\mathbb{R}[X]$.
  \item The matrix $\begin{pmatrix} 1 & \pi & 3 \\ 0 & 2 & 4 \\ 0 & 0 & 3 \end{pmatrix}$ is diagonalisable.
  \item If a rectangle $R := \{(x, y) \in \mathbb{R}^2 \mid A \le x \le B, \, C \le y \le D\}$ can be covered (allowing overlaps) by $25$ discs of radius $1$ then it can also be covered by $101$ discs of radius $\frac{1}{2}$.
  \item Given any integer $n \ge 2$, we can always find an integer $m$ such that each of the $n - 1$ consecutive integers $m + 2, m + 3, \dots, m + n$ are composite.
  \item The $10 \times 10$ matrix {[obscured in source; last column $v_1 w_{10}, v_2 w_{10}, \dots, v_{10} w_{10}$]} has rank $2$, where $v_i, w_i \in \mathbb{C}$.
  \item If every continuous function on $X \subset \mathbb{R}^2$ is bounded, then $X$ is compact.
  \item The graph of $xy = 1$ in $\mathbb{C}^2$ is connected.
  \item If $z_1, z_2, z_3, z_4 \in \mathbb{C}$ satisfy $z_1 + z_2 + z_3 + z_4 = 0$ and $|z_1|^2 + |z_2|^2 + |z_3|^2 + |z_4|^2 = 1$, then the least value of $|z_1 - z_2|^2 + |z_1 - z_4|^2 + |z_2 - z_3|^2 + |z_3 - z_4|^2$ is $2$.
  \item Consider the differential equations (with $y$ a function of $x$):
  \begin{enumerate}[label=(\arabic*)]
    \item $\dfrac{dy}{dx} = y$, $y(0) = 0$,
    \item $\dfrac{dy}{dx} = |y|^{1/3}$, $y(0) = 0$.
  \end{enumerate}
  Then (1) has infinitely many solutions but (2) has finite number of solutions.
\end{enumerate}

\newpage

% ==============================================================================
\section{GS-2013 (Mathematics)}
% ==============================================================================

\begin{center}
  \textbf{TATA INSTITUTE OF FUNDAMENTAL RESEARCH}\\
  \textbf{Written Test in MATHEMATICS --- December 9, 2012}\\
  \textbf{Duration:} Two hours (2 hours)
\end{center}

\subsection*{Instructions}
There are 4 parts and 40 True/False questions in total, consisting of 10 questions in each part (Part A, Part B, Part C, and Part D). Every correct answer carries $+1$ mark and every wrong answer carries $-1$ mark.

\subsection*{Notation and Conventions}
\begin{itemize}
  \item $\mathbb{N} :=$ Set of natural numbers; $\mathbb{Z} :=$ Set of integers; $\mathbb{Q} :=$ Set of rational numbers; $\mathbb{R} :=$ Set of real numbers; $\mathbb{C} :=$ Set of complex numbers.
  \item $\mathbb{R}^n :=$ $n$-dimensional vector space over $\mathbb{R}$.
  \item $(a, b) := \{x \in \mathbb{R} \mid a < x < b\}$, the open interval.
  \item A sequence is always indexed by natural numbers.
  \item Subsets of $\mathbb{R}^n$ are assumed to carry the induced topology.
\end{itemize}

\subsection*{Part A (True/False)}
\begin{enumerate}
  \item Every countable group $G$ has only countably many distinct subgroups.
  \item Any automorphism of the group $(\mathbb{Q}, +)$ under addition is of the form $x \mapsto qx$ for some $q \in \mathbb{Q}$.
  \item The equation $x^3 + 3x - 4 = 0$ has exactly one real root.
  \item The equation $x^3 + 10x^2 - 100x + 1729 = 0$ has at least one complex root $\alpha$ such that $|\alpha| > 12$.
  \item All non-trivial proper subgroups of $(\mathbb{R}, +)$ are cyclic.
  \item Every infinite abelian group has at least one element of infinite order.
  \item If $A$ and $B$ are similar matrices then every eigenvector of $A$ is an eigenvector of $B$.
  \item If a real square matrix $A$ is similar to a diagonal matrix and satisfies $A^n = 0$ for some $n$, then $A$ must be the zero matrix.
  \item There is an element of order $51$ in the multiplicative group $(\mathbb{Z}/103\mathbb{Z})^*$.
  \item Any normal subgroup of order $2$ is contained in the center of the group.
\end{enumerate}

\subsection*{Part B (True/False)}
\begin{enumerate}[resume]
  \item Consider the sequences $x_n = \sum_{j=1}^n \frac{1}{j}$ and $y_n = \sum_{j=1}^n \frac{1}{j^2}$. Then $\{x_n\}$ is Cauchy but $\{y_n\}$ is not.
  \item $\displaystyle \lim_{x \to 0} \frac{\sin(x^2)}{\text{[obscured]}} = 1$.
  \item Let $f: [a, b] \to [c, d]$ and $g: [c, d] \to \mathbb{R}$ be Riemann integrable functions defined on the closed intervals $[a, b]$ and $[c, d]$ respectively. Then the composite $g \circ f$ is also Riemann integrable.
  \item Let $f: \mathbb{R} \to \mathbb{R}$ be defined by $f(x) = \sin x^3$. Then $f$ is continuous but not uniformly continuous.
  \item Let $x_1 \in (0, 1)$ be a real number between $0$ and $1$. For $n \ge 1$, define $x_{n+1} = x_n - x_n^{n+1}$. Then $\lim_{n \to \infty} x_n$ exists.
  \item Suppose $\{a_i\}$ is a sequence in $\mathbb{R}$ such that $\sum |a_i||x_i| < \infty$ whenever $\sum |x_i| < \infty$. Then $\{a_i\}$ is a bounded sequence.
  \item The integral $\displaystyle \int_0^\infty e^{-x^5}\,dx$ is convergent.
  \item Let $P(x) = 1 + x + \frac{x^2}{2!} + \frac{x^3}{3!} + \dots + \frac{x^n}{n!}$ where $n$ is a large positive integer. Then $\lim_{x \to \infty} \frac{e^x}{P(x)} = 1$.
  \item Every differentiable function $f: (0, 1) \to [0, 1]$ is uniformly continuous.
  \item Consider the function $f(x) = ax + b$ with $a, b \in \mathbb{R}$. Then the iteration $x_{n+1} = f(x_n)$, $n \ge 0$, for a given $x_0$ converges to $b/(1 - a)$ whenever $0 < a < 1$.
\end{enumerate}

\subsection*{Part C (True/False)}
\begin{enumerate}[resume]
  \item Every homeomorphism of the $2$-sphere to itself has a fixed point.
  \item The intervals $[0, 1)$ and $(0, 1)$ are homeomorphic.
  \item Let $X$ be a complete metric space such that distance between any two points is less than $1$. Then $X$ is compact.
  \item There exists a continuous surjective function from $S^1$ onto $\mathbb{R}$.
  \item There exists a complete metric on the open interval $(0, 1)$ inducing the usual topology.
  \item There exists a continuous surjective map from the complex plane onto the non-zero reals.
  \item If every differentiable function on a subset $X \subset \mathbb{R}^n$ (i.e., restriction of a differentiable function on a neighbourhood of $X$) is bounded, then $X$ is compact.
  \item Let $f: X \to Y$ be a continuous map between metric spaces. If $f$ is a bijection, then its inverse is also continuous.
  \item Let $f$ be a function on the closed interval $[0, 1]$ defined by
  \[
    f(x) = \begin{cases} x & \text{if } x \text{ is rational,} \\ x^2 & \text{if } x \text{ is irrational.} \end{cases}
  \]
  Then $f$ is continuous at $0$ and $1$.
  \item There exists an infinite subset $S \subset \mathbb{R}^3$ such that any three vectors in $S$ are linearly independent.
\end{enumerate}

\subsection*{Part D (True/False)}
\begin{enumerate}[resume]
  \item The inequality $\sqrt{n+1} - \sqrt{n} < \frac{1}{\sqrt{n}}$ is false for all $n$ such that $101 \le n \le 2000$.
  \item $\displaystyle \lim_{n \to \infty} \left((n + 1)^{1/3} - n^{1/3}\right) = \infty$.
  \item There exists a bijection between $\mathbb{R}^2$ and the open interval $(0, 1)$.
  \item Let $S$ be the set of all sequences $\{a_1, a_2, \dots, a_n, \dots\}$ where each entry $a_i$ is either $0$ or $1$. Then $S$ is countable.
  \item Let $\{a_n\}$ be any non-constant sequence in $\mathbb{R}$ such that $a_{n+1} = \frac{a_n + a_{n+2}}{2}$ for all $n \ge 1$. Then $\{a_n\}$ is unbounded.
  \item The function $f: \mathbb{Z} \to \mathbb{R}$ defined by $f(n) = n^3 - 3n$ is injective.
  \item The polynomial $x^3 + 3x - 2\pi$ is irreducible over $\mathbb{R}$.
  \item Let $V$ be the vector space consisting of polynomials with real coefficients in variable $t$ of degree $< 9$. Let $D: V \to V$ be the linear operator defined by $D(f) := \frac{df}{dt}$. Then $0$ is an eigenvalue of $D$.
  \item If $A$ is a complex $n \times n$ matrix with $A^2 = A$ then $\operatorname{rank} A = \operatorname{trace} A$.
  \item The series $\displaystyle 1 - \frac{1}{\sqrt{2}} + \frac{1}{\sqrt{3}} - \frac{1}{\sqrt{4}} + \dots$ is divergent.
\end{enumerate}

\newpage

% ==============================================================================
\section{GS-2014 (Mathematics)}
% ==============================================================================

\begin{center}
  \textbf{TATA INSTITUTE OF FUNDAMENTAL RESEARCH}\\
  \textbf{Written Test in MATHEMATICS --- December 8, 2013}\\
  \textbf{Duration:} Two hours (2 hours)
\end{center}

\subsection*{Instructions}
There are thirty (30) multiple choice questions divided into two parts: Part I consists of 20 questions and Part II consists of 10 questions. Each correct answer gets $1$ mark, with no negative marking for wrong answers.

\subsection*{Notation and Conventions}
\begin{itemize}
  \item $\mathbb{N} := \{1, 2, 3, \dots\}$; $\mathbb{Z} :=$ Set of integers; $\mathbb{Q} :=$ Set of rational numbers; $\mathbb{R} :=$ Set of real numbers; $\mathbb{C} :=$ Set of complex numbers; $\mathbb{R}^* :=$ Set of non-zero real numbers; $\mathbb{C}^* :=$ Set of non-zero complex numbers.
  \item $\mathbb{R}^n :=$ $n$-dimensional vector space over $\mathbb{R}$.
  \item $(a, b) := \{x \in \mathbb{R} \mid a < x < b\}$, $[a, b) := \{x \in \mathbb{R} \mid a \le x < b\}$, $[a, b] := \{x \in \mathbb{R} \mid a \le x \le b\}$.
  \item A sequence is always indexed by the set of natural numbers.
  \item The cyclic group with $n$ elements is denoted by $\mathbb{Z}/n$.
  \item Subsets of $\mathbb{R}^n$ are assumed to carry the induced topology.
  \item For any set $S$, the cardinality of the set is denoted by $|S|$.
\end{itemize}

\subsection*{Part I}
\begin{enumerate}
  \item Let $A, B, C$ be three subsets of $\mathbb{R}$. The negation of the following statement:
  \begin{quote}
    For every $\epsilon > 1$, there exists $a \in A$ and $b \in B$ such that for all $c \in C$, $|a - c| < \epsilon$ and $|b - c| > \epsilon$
  \end{quote}
  is
  \begin{enumerate}[label=\Alph*.]
    \item there exists $\epsilon \le 1$, such that for all $a \in A$ and $b \in B$ there exists $c \in C$ such that $|a - c| \ge \epsilon$ or $|b - c| \le \epsilon$.
    \item there exists $\epsilon \le 1$, such that for all $a \in A$ and $b \in B$ there exists $c \in C$ such that $|a - c| \ge \epsilon$ and $|b - c| \le \epsilon$.
    \item there exists $\epsilon > 1$, such that for all $a \in A$ and $b \in B$ there exists $c \in C$ such that $|a - c| \ge \epsilon$ and $|b - c| \le \epsilon$.
    \item there exists $\epsilon > 1$, such that for all $a \in A$ and $b \in B$ there exists $c \in C$ such that $|a - c| \ge \epsilon$ or $|b - c| \le \epsilon$.
  \end{enumerate}

  \item Let $f: \mathbb{R} \to \mathbb{R}$ be a continuous bounded function, then:
  \begin{enumerate}[label=\Alph*.]
    \item $f$ has to be uniformly continuous.
    \item there exists an $x \in \mathbb{R}$ such that $f(x) = x$.
    \item $f$ cannot be increasing.
    \item $\lim_{x \to \infty} f(x)$ exists.
  \end{enumerate}

  \item Let $f: \mathbb{R} \to \mathbb{R}$ be a differentiable function such that $\lim_{x \to +\infty} f'(x) = 1$; then
  \begin{enumerate}[label=\Alph*.]
    \item $f$ is bounded.
    \item $f$ is increasing.
    \item $f$ is unbounded.
    \item $f'$ is bounded.
  \end{enumerate}

  \item Let $f$ be the real valued function on $[0, \infty)$ defined by
  \[
    f(x) = \begin{cases} x^{2/3} \log x, & \text{for } x > 0, \\ 0, & \text{if } x = 0. \end{cases}
  \]
  Then
  \begin{enumerate}[label=\Alph*.]
    \item $f$ is discontinuous at $x = 0$.
    \item $f$ is continuous on $[0, \infty)$, but not uniformly continuous on $[0, \infty)$.
    \item $f$ is uniformly continuous on $[0, \infty)$.
    \item $f$ is not uniformly continuous on $[0, \infty)$, but uniformly continuous on $(0, \infty)$.
  \end{enumerate}

  \item Let $a_n = (n + 1)^{100} e^{-\sqrt{n}}$ for $n \ge 1$. Then the sequence $(a_n)_n$ is
  \begin{enumerate}[label=\Alph*.]
    \item unbounded.
    \item bounded but does not converge.
    \item bounded and converges to $1$.
    \item bounded and converges to $0$.
  \end{enumerate}

  \item Let $f: [0, 1] \to \mathbb{R}$ be a continuous function. Which of the following statements is always true?
  \begin{enumerate}[label=\Alph*.]
    \item $\int_0^1 f^2(x)\,dx = \left(\int_0^1 f(x)\,dx\right)^2$.
    \item $\int_0^1 f^2(x)\,dx \le \left(\int_0^1 |f(x)|\,dx\right)^2$.
    \item $\int_0^1 f^2(x)\,dx \ge \left(\int_0^1 |f(x)|\,dx\right)^2$.
    \item $\int_0^1 f^2(x)\,dx \le \left(\int_0^1 f(x)\,dx\right)^2$.
  \end{enumerate}

  \item Let $f_n(x)$, for $n \ge 1$, be a sequence of continuous nonnegative functions on $[0, 1]$ such that $\lim_{n \to \infty} \int_0^1 f_n(x)\,dx = 0$. Which of the following statements is always correct?
  \begin{enumerate}[label=\Alph*.]
    \item $f_n \to 0$ uniformly on $[0, 1]$.
    \item $f_n$ may not converge uniformly but converges to $0$ point-wise.
    \item $f_n$ will converge point-wise and the limit may be non-zero.
    \item $f_n$ is not guaranteed to have a point-wise limit.
  \end{enumerate}

  \item Let $f: \mathbb{R} \to \mathbb{R}$ be a continuous function such that $|f(x) - f(y)| \ge \frac{1}{2}|x - y|$ for all $x, y \in \mathbb{R}$. Then
  \begin{enumerate}[label=\Alph*.]
    \item $f$ is both one-to-one and onto.
    \item $f$ is one-to-one but may not be onto.
    \item $f$ is onto but may not be one-to-one.
    \item $f$ is neither one-to-one nor onto.
  \end{enumerate}

  \item Let $A(\theta) = \begin{pmatrix} \cos\theta & \sin\theta \\ -\sin\theta & \cos\theta \end{pmatrix}$, where $\theta \in (0, 2\pi)$. Mark the correct statement below.
  \begin{enumerate}[label=\Alph*.]
    \item $A(\theta)$ has eigenvectors in $\mathbb{R}^2$ for all $\theta \in (0, 2\pi)$.
    \item $A(\theta)$ does not have an eigenvector in $\mathbb{R}^2$ for any $\theta \in (0, 2\pi)$.
    \item $A(\theta)$ has eigenvectors in $\mathbb{R}^2$ for exactly one value of $\theta \in (0, 2\pi)$.
    \item $A(\theta)$ has eigenvectors in $\mathbb{R}^2$ for exactly $2$ values of $\theta \in (0, 2\pi)$.
  \end{enumerate}

  \item Let $\mathcal{C} \subset \mathbb{Z} \times \mathbb{Z}$ be the set of integer pairs $(a, b)$ for which the three complex roots $r_1, r_2$ and $r_3$ of the polynomial $p(x) = x^3 - 2x^2 + ax - b$ satisfy $r_1^3 + r_2^3 + r_3^3 = 0$. Then the cardinality of $\mathcal{C}$ is
  \begin{enumerate}[label=\Alph*.]
    \item $|\mathcal{C}| = \infty$
    \item $|\mathcal{C}| = 0$
    \item $|\mathcal{C}| = 1$
    \item $1 < |\mathcal{C}| < \infty$
  \end{enumerate}

  \item Let $A$ be an $n \times n$ matrix with real entries such that $A^k = 0$ ($0$-matrix) for some $k \in \mathbb{N}$. Then
  \begin{enumerate}[label=\Alph*.]
    \item $A$ has to be the $0$ matrix.
    \item $\operatorname{trace}(A)$ could be non-zero.
    \item $A$ is diagonalizable.
    \item $0$ is the only eigenvalue of $A$.
  \end{enumerate}

  \item There exists a map $f: \mathbb{Z} \to \mathbb{Q}$ such that $f$
  \begin{enumerate}[label=\Alph*.]
    \item is bijective and increasing.
    \item is onto and decreasing.
    \item is bijective and satisfies $f(n) \ge 0$ if $n \le 0$.
    \item has uncountable image.
  \end{enumerate}

  \item Let $S$ be the set of all tuples $(x, y)$ with $x, y$ non-negative real numbers satisfying $x + y = 2n$ for a fixed $n \in \mathbb{N}$. Then the supremum value of $x^2 y^2 (x^2 + y^2)$ on the set $S$ is
  \begin{enumerate}[label=\Alph*.]
    \item $3n^6$
    \item $2n^6$
    \item $4n^6$
    \item $n^6$.
  \end{enumerate}

  \item Let $G$ be a group and let $H$ and $K$ be two subgroups of $G$. If both $H$ and $K$ have $12$ elements, which of the following numbers cannot be the cardinality of the set $HK = \{hk : h \in H, \, k \in K\}$?
  \begin{enumerate}[label=\Alph*.]
    \item $72$
    \item $60$
    \item $48$
    \item $36$.
  \end{enumerate}

  \item How many proper subgroups does the group $\mathbb{Z} \oplus \mathbb{Z}$ have?
  \begin{enumerate}[label=\Alph*.]
    \item $1$
    \item $2$
    \item $3$
    \item infinitely many.
  \end{enumerate}

  \item $X$ is a metric space. $Y$ is a closed subset of $X$ such that the distance between any two points in $Y$ is at most $1$. Then
  \begin{enumerate}[label=\Alph*.]
    \item $Y$ is compact.
    \item any continuous function from $Y \to \mathbb{R}$ is bounded.
    \item $Y$ is not an open subset of $X$.
    \item none of the above.
  \end{enumerate}

  \item Let $f: \mathbb{R} \to \mathbb{R}$ be a continuous function and let $S$ be a non-empty proper subset of $\mathbb{R}$. Which one of the following statements is always true? (Here $\overline{A}$ denotes the closure of $A$ and $A^\circ$ denotes the interior of $A$.)
  \begin{enumerate}[label=\Alph*.]
    \item $f(S)^\circ \subseteq f(S^\circ)$
    \item $f(\overline{S}) \subseteq \overline{f(S)}$
    \item $f(\overline{S}) \supseteq \overline{f(S)}$
    \item $f(S)^\circ \supseteq f(S^\circ)$.
  \end{enumerate}

  \item What is the last digit of $97^{2013}$?
  \begin{enumerate}[label=\Alph*.]
    \item $1$
    \item $3$
    \item $7$
    \item $9$.
  \end{enumerate}

  \item For $n \in \mathbb{N}$ we define $s_n = 1^3 + 2^3 + 3^3 + \dots + n^3$. Which of the following holds for all $n \in \mathbb{N}$?
  \begin{enumerate}[label=\Alph*.]
    \item $s_n$ is an odd integer.
    \item $s_n = n^2(n + 1)^2 / 4$.
    \item $s_n = n(n + 1)(2n + 1) / 6$.
    \item none of the above.
  \end{enumerate}

  \item Let $C$ denote the cube $[-1, 1]^3 \subset \mathbb{R}^3$. How many rotations are there in $\mathbb{R}^3$ {[obscured in source]}?
  \begin{enumerate}[label=\Alph*.]
    \item $6$
    \item {[obscured in source]}
    \item {[obscured in source]}
    \item $24$.
  \end{enumerate}
\end{enumerate}

\subsection*{Part II}
\begin{enumerate}[resume]
  \item Let $f: [0, 1] \to [0, \infty)$ be continuous. Suppose $\int_0^x f(t)\,dt \ge f(x)$ for all $x \in [0, 1]$. Then
  \begin{enumerate}[label=\Alph*.]
    \item no such function exists.
    \item there are infinitely many such functions.
    \item there is only one such function.
    \item there are exactly two such functions.
  \end{enumerate}

  \item Let $f: \mathbb{R}^2 \to \mathbb{R}$ be a continuous map such that $f(x) = 0$ for only finitely many values of $x$. Which of the following is true?
  \begin{enumerate}[label=\Alph*.]
    \item either $f(x) \le 0$ for all $x$, or $f(x) \ge 0$ for all $x$.
    \item the map $f$ is onto.
    \item the map $f$ is one-to-one.
    \item none of the above.
  \end{enumerate}

  \item Let $S_n$ be the symmetric group of $n$ letters. There exists an onto group homomorphism
  \begin{enumerate}[label=\Alph*.]
    \item from $S_5$ to $S_4$
    \item from $S_4$ to $S_2$
    \item from $S_5$ to $\mathbb{Z}/5$
    \item from $S_4$ to $\mathbb{Z}/4$.
  \end{enumerate}

  \item Let $H_1, H_2$ be two distinct subgroups of a finite group $G$, each of order $2$. Let $H$ be the smallest subgroup containing $H_1$ and $H_2$. Then the order of $H$ is
  \begin{enumerate}[label=\Alph*.]
    \item always $2$
    \item always $4$
    \item always $8$
    \item none of the above.
  \end{enumerate}

  \item Which of the following groups are isomorphic?
  \begin{enumerate}[label=\Alph*.]
    \item $\mathbb{R}$ and $\mathbb{C}$
    \item $\mathbb{R}^*$ and $\mathbb{C}^*$
    \item $S_3 \times \mathbb{Z}/4$ and $S_4$
    \item $\mathbb{Z}/2 \times \mathbb{Z}/2$ and $\mathbb{Z}/4$.
  \end{enumerate}

  \item The number of irreducible polynomials of the form $x^2 + ax + b$ with $a, b$ in the field $\mathbb{F}_7$ of $7$ elements is:
  \begin{enumerate}[label=\Alph*.]
    \item $7$
    \item $21$
    \item $35$
    \item $49$.
  \end{enumerate}

  \item $X$ is a topological space of infinite cardinality which is homeomorphic to $X \times X$. Then
  \begin{enumerate}[label=\Alph*.]
    \item $X$ is not connected.
    \item $X$ is not compact.
    \item $X$ is not homeomorphic to a subset of $\mathbb{R}$.
    \item none of the above.
  \end{enumerate}

  \item Let $X$ be a non-empty topological space such that every function $f: X \to \mathbb{R}$ is continuous. Then
  \begin{enumerate}[label=\Alph*.]
    \item $X$ has the discrete topology.
    \item $X$ has the indiscrete topology.
    \item $X$ is compact.
    \item $X$ is not connected.
  \end{enumerate}

  \item Let $f: X \to Y$ be a continuous map between metric spaces. Then $f(X)$ is a complete subset of $Y$ if
  \begin{enumerate}[label=\Alph*.]
    \item the space $X$ is compact.
    \item the space $Y$ is compact.
    \item the space $X$ is complete.
    \item the space $Y$ is complete.
  \end{enumerate}

  \item How many maps $\phi: \mathbb{N} \cup \{0\} \to \mathbb{N} \cup \{0\}$ are there with the property that $\phi(ab) = \phi(a) + \phi(b)$ for all $a, b \in \mathbb{N} \cup \{0\}$?
  \begin{enumerate}[label=\Alph*.]
    \item none
    \item {[obscured in source]}
    \item {[obscured in source]}
    \item {[obscured in source]}
  \end{enumerate}
\end{enumerate}

\newpage

% ==============================================================================
\section{GS-2015 (Mathematics)}
% ==============================================================================

\begin{center}
  \textbf{TATA INSTITUTE OF FUNDAMENTAL RESEARCH}\\
  \textbf{Written Test in MATHEMATICS --- December 14, 2014}\\
  \textbf{Duration:} Two hours (2 hours)
\end{center}

\subsection*{Instructions}
There are thirty (30) multiple choice questions divided into two parts: Part I consists of 15 questions and Part II consists of 15 questions. Each correct answer gets $1$ mark, with no negative marking for wrong answers.

\subsection*{Notation and Conventions}
\begin{itemize}
  \item $\mathbb{N} := \{1, 2, 3, \dots\}$; $\mathbb{Z} :=$ Set of integers; $\mathbb{Q} :=$ Set of rational numbers; $\mathbb{R} :=$ Set of real numbers; $\mathbb{C} :=$ Set of complex numbers.
  \item $\mathbb{R}^n :=$ $n$-dimensional vector space over $\mathbb{R}$.
  \item $(a, b) := \{x \in \mathbb{R} \mid a < x < b\}$, $(a, b] := \{x \in \mathbb{R} \mid a < x \le b\}$, $[a, b) := \{x \in \mathbb{R} \mid a \le x < b\}$, $[a, b] := \{x \in \mathbb{R} \mid a \le x \le b\}$.
  \item A sequence is always indexed by the set of natural numbers.
  \item The cyclic group with $n$ elements is denoted by $\mathbb{Z}_n$.
  \item Unless stated otherwise, subsets of $\mathbb{R}^n$ carry the induced topology.
  \item For any set $S$, the cardinality of $S$ is denoted by $|S|$.
\end{itemize}

\subsection*{Part I}
\begin{enumerate}
  \item Let $A$ be an invertible $10 \times 10$ matrix with real entries such that the sum of each row is $1$. Then
  \begin{enumerate}[label=\Alph*.]
    \item The sum of the entries of each row of the inverse of $A$ is $1$.
    \item The sum of the entries of each column of the inverse of $A$ is $1$.
    \item The trace of the inverse of $A$ is non-zero.
    \item None of the above.
  \end{enumerate}

  \item Let $f: \mathbb{R} \to \mathbb{R}$ be a continuous function. Which one of the following sets cannot be the image of $(0, 1]$ under $f$?
  \begin{enumerate}[label=\Alph*.]
    \item $\{0\}$
    \item $(0, 1)$
    \item $[0, 1)$
    \item $[0, 1]$.
  \end{enumerate}

  \item Let $A$ be a $10 \times 10$ matrix with complex entries such that all its eigenvalues are non-negative real numbers, and at least one eigenvalue is positive. Which of the following statements is always false?
  \begin{enumerate}[label=\Alph*.]
    \item There exists a matrix $B$ such that $AB - BA = B$.
    \item There exists a matrix $B$ such that $AB - BA = A$.
    \item There exists a matrix $B$ such that $AB + BA = A$.
    \item There exists a matrix $B$ such that $AB + BA = B$.
  \end{enumerate}

  \item Let $S$ be the collection of (isomorphism classes of) groups $G$ which have the property that every element of $G$ commutes only with the identity element and itself. Then
  \begin{enumerate}[label=\Alph*.]
    \item $|S| = 1$
    \item $|S| = 2$
    \item $|S| \ge 3$ and is finite
    \item $|S| = \infty$.
  \end{enumerate}

  \item Let $f: \mathbb{R} \to \mathbb{R}$ denote the function defined by $f(x) = (1 - x^2)^{3/2}$ if $|x| < 1$, and $f(x) = 0$ if $|x| \ge 1$. Which of the following statements is correct?
  \begin{enumerate}[label=\Alph*.]
    \item $f$ is not continuous.
    \item $f$ is continuous but not differentiable.
    \item $f$ is differentiable but $f'$ is not continuous.
    \item $f$ is differentiable and $f'$ is continuous.
  \end{enumerate}

  \item Let $A$ be the $2 \times 2$ matrix $\begin{pmatrix} \sin\frac{\pi}{18} & -\sin\frac{4\pi}{9} \\ \sin\frac{4\pi}{9} & \sin\frac{\pi}{18} \end{pmatrix}$. Then the smallest number $n \in \mathbb{N}$ such that $A^n = I$ is
  \begin{enumerate}[label=\Alph*.]
    \item {[obscured in source]}
    \item {[obscured in source]}
    \item {[obscured in source]}
    \item $27$.
  \end{enumerate}

  \item Let $f$ and $g$ be two functions from $[0, 1]$ to $[0, 1]$ with $f$ strictly increasing. Which of the following statements is always correct?
  \begin{enumerate}[label=\Alph*.]
    \item If $g$ is continuous, then $f \circ g$ is continuous.
    \item If $f$ is continuous, then $f \circ g$ is continuous.
    \item If $f$ and $f \circ g$ are continuous, then $g$ is continuous.
    \item If $g$ and $f \circ g$ are continuous, then $f$ is continuous.
  \end{enumerate}

  \item Let $f(x) = \frac{e^{-1/x}}{x}$, where $x \in (0, 1)$. Then, on $(0, 1)$
  \begin{enumerate}[label=\Alph*.]
    \item $f$ is uniformly continuous.
    \item $f$ is continuous but not uniformly continuous.
    \item $f$ is unbounded.
    \item $f$ is not continuous.
  \end{enumerate}

  \item Let $\{a_n\}$ be a sequence of real numbers such that $|a_{n+1} - a_n| \le \frac{n^2}{2^n}$ for all $n \in \mathbb{N}$. Then
  \begin{enumerate}[label=\Alph*.]
    \item The sequence $\{a_n\}$ may be unbounded.
    \item The sequence $\{a_n\}$ is bounded but may not converge.
    \item The sequence $\{a_n\}$ has exactly two limit points.
    \item The sequence $\{a_n\}$ is convergent.
  \end{enumerate}

  \item For a group $G$, let $\operatorname{Aut}(G)$ denote the group of automorphisms of $G$. Which of the following statements is true?
  \begin{enumerate}[label=\Alph*.]
    \item $\operatorname{Aut}(\mathbb{Z})$ is isomorphic to $\mathbb{Z}_2$.
    \item If $G$ is cyclic, then $\operatorname{Aut}(G)$ is cyclic.
    \item If $\operatorname{Aut}(G)$ is trivial, then $G$ is trivial.
    \item $\operatorname{Aut}(\mathbb{Z})$ is isomorphic to $\mathbb{Z}$.
  \end{enumerate}

  \item Let $\{a_n\}$ be a sequence of real numbers. Which of the following is true?
  \begin{enumerate}[label=\Alph*.]
    \item If $\sum a_n$ converges, then so does $\sum a_n^4$.
    \item If $\sum |a_n|$ converges, then so does $\sum a_n^2$.
    \item If $\sum a_n$ diverges, then so does $\sum a_n^3$.
    \item If $\sum |a_n|$ diverges, then so does $\sum a_n^2$.
  \end{enumerate}

  \item Let $f: \mathbb{R} \to \mathbb{R}$ be an infinitely differentiable function that vanishes at $10$ distinct points in $\mathbb{R}$. Suppose $f^{(n)}$ denotes the $n$-th derivative of $f$, for $n \ge 1$. Which of the following statements is always true?
  \begin{enumerate}[label=\Alph*.]
    \item $f^{(n)}$ has at least $10$ zeros, for $1 \le n \le 8$.
    \item $f^{(n)}$ has at least one zero, for $1 \le n \le 9$.
    \item $f^{(n)}$ has at least $10$ zeros, for $n \ge 10$.
    \item $f^{(n)}$ has at least one zero, for $n \ge 9$.
  \end{enumerate}

  \item For a real number $t > 0$, let $\sqrt{t}$ denote the positive square root of $t$. For a real number $x > 0$, let $F(x) = \int_{x^2}^{4x^2} \sin\sqrt{t}\,dt$. If $F'$ is the derivative of $F$, then
  \begin{enumerate}[label=\Alph*.]
    \item $F'\left(\frac{\pi}{2}\right) = 0$
    \item $F'\left(\frac{\pi}{2}\right) = \pi$
    \item $F'\left(\frac{\pi}{2}\right) = -\pi$
    \item $F'\left(\frac{\pi}{2}\right) = 2\pi$.
  \end{enumerate}

  \item Let $n \in \mathbb{N}$ be a six digit number whose base $10$ expansion is of the form $abcabc$, where $a, b, c$ are digits between $0$ and $9$ and $a$ is non-zero. Then
  \begin{enumerate}[label=\Alph*.]
    \item $n$ is divisible by $5$.
    \item $n$ is divisible by $8$.
    \item $n$ is divisible by $13$.
    \item $n$ is divisible by $17$.
  \end{enumerate}

  \item The series $\displaystyle \sum_{n=1}^\infty \frac{\cos(3^n x)}{2^n}$
  \begin{enumerate}[label=\Alph*.]
    \item diverges, for all rational $x \in \mathbb{R}$.
    \item diverges, for some irrational $x \in \mathbb{R}$.
    \item converges, for some but not all $x \in \mathbb{R}$.
    \item converges, for all $x \in \mathbb{R}$.
  \end{enumerate}
\end{enumerate}

\subsection*{Part II}
\begin{enumerate}[resume]
  \item Let $X$ be a proper closed subset of $[0, 1]$. Which of the following statements is always true?
  \begin{enumerate}[label=\Alph*.]
    \item The set $X$ is countable.
    \item There exists $x \in X$ such that $X \setminus \{x\}$ is closed.
    \item The set $X$ contains an open interval.
    \item None of the above.
  \end{enumerate}

  \item In how many ways can the group $\mathbb{Z}_5$ act on the set $\{1, 2, 3, 4, 5\}$?
  \begin{enumerate}[label=\Alph*.]
    \item $5$
    \item $24$
    \item $25$
    \item $120$.
  \end{enumerate}

  \item Let $f$ be a function from $\{1, 2, \dots, 10\}$ to $\mathbb{R}$ such that
  \[
    \left(\sum_{i=1}^{10} \frac{|f(i)|}{2^i}\right)^2 = \left(\sum_{i=1}^{10} |f(i)|^2\right) \left(\sum_{i=1}^{10} \frac{1}{4^i}\right).
  \]
  Mark the correct statement.
  \begin{enumerate}[label=\Alph*.]
    \item There are uncountably many $f$ with this property.
    \item There are only countably infinitely many $f$ with this property.
    \item There is exactly one such $f$.
    \item There is no such $f$.
  \end{enumerate}

  \item Let $U_1 \supset U_2 \supset \cdots$ be a decreasing sequence of open sets in Euclidean $3$-space $\mathbb{R}^3$. What can we say about the set $\bigcap U_i$?
  \begin{enumerate}[label=\Alph*.]
    \item It is infinite.
    \item It is open.
    \item It is non-empty.
    \item None of the above.
  \end{enumerate}

  \item Let $n \ge 1$ and let $A$ be an $n \times n$ matrix with real entries such that $A^k = 0$ for some $k \ge 1$. Let $I$ be the identity $n \times n$ matrix. Then
  \begin{enumerate}[label=\Alph*.]
    \item $I + A$ need not be invertible.
    \item $\det(I + A)$ can be any non-zero real number.
    \item $\det(I + A) = 1$.
    \item $A^n$ is a non-zero matrix.
  \end{enumerate}

  \item Let $f: [0, 1] \to \mathbb{R}$ be a fixed continuous function such that $f$ is differentiable on $(0, 1)$ and $f(0) = f(1) = 0$. Then the equation $f(x) = f'(x)$ admits
  \begin{enumerate}[label=\Alph*.]
    \item no solution $x \in (0, 1)$.
    \item more than one solution $x \in (0, 1)$.
    \item exactly one solution $x \in (0, 1)$.
    \item at least one solution $x \in (0, 1)$.
  \end{enumerate}

  \item A complex number $\alpha \in \mathbb{C}$ is called algebraic if there is a non-zero polynomial $P(x) \in \mathbb{Q}[x]$ with rational coefficients such that $P(\alpha) = 0$. Which of the following statements is true?
  \begin{enumerate}[label=\Alph*.]
    \item There are only finitely many algebraic numbers.
    \item All complex numbers are algebraic.
    \item $\sin\left(\frac{\pi}{3}\right) + \cos\left(\frac{\pi}{4}\right)$ is algebraic.
    \item None of the above.
  \end{enumerate}

  \item For a group $G$, let $F(G)$ denote the collection of all subgroups of $G$. Which one of the following situations can occur?
  \begin{enumerate}[label=\Alph*.]
    \item $G$ is finite but $F(G)$ is infinite.
    \item $G$ is infinite but $F(G)$ is finite.
    \item $G$ is countable but $F(G)$ is uncountable.
    \item $G$ is uncountable but $F(G)$ is countable.
  \end{enumerate}

  \item Let $f: \mathbb{R} \to \mathbb{R}$ be a continuous function and $A \subset \mathbb{R}$ be defined by
  \[
    A = \{y \in \mathbb{R} : y = \lim_{n \to \infty} f(x_n) \text{ for some sequence } x_n \to +\infty\}.
  \]
  Then the set $A$ is necessarily
  \begin{enumerate}[label=\Alph*.]
    \item a connected set.
    \item a compact set.
    \item a singleton set.
    \item none of the above.
  \end{enumerate}

  \item How many finite sequences $x_1, x_2, \dots, x_m$ are there such that each $x_i = 1$ or $2$, and $\sum_{i=1}^m x_i = 10$?
  \begin{enumerate}[label=\Alph*.]
    \item $89$
    \item $91$
    \item $92$
    \item $120$.
  \end{enumerate}

  \item Let $(X, d)$ be a path connected metric space with at least two elements, and let $S = \{d(x, y) : x, y \in X\}$. Which of the following statements is not necessarily true?
  \begin{enumerate}[label=\Alph*.]
    \item $S$ is infinite.
    \item $S$ contains a non-zero rational number.
    \item $S$ is connected.
    \item $S$ is a closed subset of $\mathbb{R}$.
  \end{enumerate}

  \item Let $X \subset \mathbb{R}$ and let $f, g: X \to X$ be continuous functions such that $f(X) \cap g(X) = \emptyset$ and $f(X) \cup g(X) = X$. Which one of the following sets cannot be equal to $X$?
  \begin{enumerate}[label=\Alph*.]
    \item $[0, 1]$
    \item $(0, 1)$
    \item $[0, 1)$
    \item $\mathbb{R}$.
  \end{enumerate}

  \item Let $X = \{(x, y) \in \mathbb{R}^2 : 2x^2 + 3y^2 = 1\}$. Endow $\mathbb{R}^2$ with the discrete topology, and $X$ with the subspace topology. Then
  \begin{enumerate}[label=\Alph*.]
    \item $X$ is a compact subset of $\mathbb{R}^2$ in this topology.
    \item $X$ is a connected subset of $\mathbb{R}^2$ in this topology.
    \item $X$ is an open subset of $\mathbb{R}^2$ in this topology.
    \item None of the above.
  \end{enumerate}

  \item Let $G$ be a group. Suppose $|G| = p^2 q$, where $p$ and $q$ are distinct prime numbers satisfying $q \not\equiv 1 \pmod p$. Which of the following is always true?
  \begin{enumerate}[label=\Alph*.]
    \item $G$ has more than one $p$-Sylow subgroup.
    \item $G$ has a normal $p$-Sylow subgroup.
    \item The number of $q$-Sylow subgroups of $G$ is divisible by $p$.
    \item $G$ has a unique $q$-Sylow subgroup.
  \end{enumerate}

  \item Let $d(x, y)$ be the usual Euclidean metric on $\mathbb{R}^2$. Which of the following metric spaces is complete?
  \begin{enumerate}[label=\Alph*.]
    \item $\mathbb{Q}^2 \subset \mathbb{R}^2$ with the metric $d(x, y)$.
    \item $[0, 1] \times [0, \infty) \subset \mathbb{R}^2$ with the metric $d'(x, y) = \frac{d(x, y)}{1 + d(x, y)}$.
    \item $(0, \infty) \times [0, \infty) \subset \mathbb{R}^2$ with the metric $d(x, y)$.
    \item $[0, 1] \times [0, 1) \subset \mathbb{R}^2$ with the metric $d''(x, y) = \min\{1, d(x, y)\}$.
  \end{enumerate}
\end{enumerate}

\newpage

% ==============================================================================
\section{GS-2016 (Mathematics)}
% ==============================================================================

\begin{center}
  \textbf{TATA INSTITUTE OF FUNDAMENTAL RESEARCH}\\
  \textbf{Written Test in MATHEMATICS --- December 13, 2015}\\
  For the Ph.D.\ Programs at TIFR (Mumbai, and CAM and ICTS, Bangalore) and for the Int.\ Ph.D.\ Programs at TIFR (CAM, Bangalore and Mumbai).\\
  \textbf{Duration:} Two hours (2 hours)
\end{center}

\subsection*{Instructions}
There are thirty (30) multiple choice questions divided into two parts: Part I consists of 20 questions and Part II consists of 10 questions. Each correct answer gets $1$ mark, with no negative marking for wrong answers.

\subsection*{Notation and Conventions}
\begin{itemize}
  \item $\mathbb{N} := \{1, 2, 3, \dots\}$; $\mathbb{Z} :=$ Set of integers; $\mathbb{Q} :=$ Set of rational numbers; $\mathbb{R} :=$ Set of real numbers; $\mathbb{C} :=$ Set of complex numbers.
  \item $\mathbb{R}^n :=$ $n$-dimensional vector space over $\mathbb{R}$.
  \item $(a, b) := \{x \in \mathbb{R} \mid a < x < b\}$, $(a, b] := \{x \in \mathbb{R} \mid a < x \le b\}$, $[a, b) := \{x \in \mathbb{R} \mid a \le x < b\}$, $[a, b] := \{x \in \mathbb{R} \mid a \le x \le b\}$.
  \item A sequence is always indexed by the set of natural numbers.
  \item The cyclic group with $n$ elements is denoted by $\mathbb{Z}/n\mathbb{Z}$.
  \item Unless stated otherwise, subsets of $\mathbb{R}^n$ carry the induced topology.
  \item For any set $S$, the cardinality of $S$ is denoted by $|S|$.
\end{itemize}

\subsection*{Part I}
\begin{enumerate}
  \item The value of the product $\left(1 + \frac{1}{1!} + \frac{1}{2!} + \dots\right)\left(1 - \frac{1}{1!} + \frac{1}{2!} - \frac{1}{3!} + \dots\right)$ is
  \begin{enumerate}[label=\Alph*.]
    \item $1$
    \item $e^2$
    \item $0$
    \item $\log_e 2$.
  \end{enumerate}

  \item Which of the following is false?
  \begin{enumerate}[label=\Alph*.]
    \item $\sum_{n=1}^\infty \sin\frac{1}{n}$ diverges
    \item $\sum_{n=1}^\infty \sin\frac{1}{n^2}$ converges
    \item $\sum_{n=1}^\infty \cos\frac{1}{n}$ diverges
    \item $\sum_{n=1}^\infty \cos\frac{1}{n^2}$ converges.
  \end{enumerate}

  \item The value of the series $\displaystyle \sum_{n=1}^\infty \frac{n}{2^n}$ is
  \begin{enumerate}[label=\Alph*.]
    \item $1$
    \item $2$
    \item $3$
    \item $4$.
  \end{enumerate}

  \item Let $f: \mathbb{R} \to \mathbb{R}$ be the function defined by $f(x) = \frac{\sin x}{|x| + \cos x}$. Then
  \begin{enumerate}[label=\Alph*.]
    \item $f$ is differentiable at all $x \in \mathbb{R}$.
    \item $f$ is not differentiable at $x = 0$.
    \item $f$ is differentiable at $x = 0$ but $f'$ is not continuous at $x = 0$.
    \item $f$ is not differentiable at $x = \frac{\pi}{2}$.
  \end{enumerate}

  \item Which of the following continuous functions $f: (0, \infty) \to \mathbb{R}$ can be extended to a continuous function on $[0, \infty)$?
  \begin{enumerate}[label=\Alph*.]
    \item $f(x) = \sin\frac{1}{x}$
    \item $f(x) = \frac{1 - \cos x}{x^2}$
    \item $f(x) = \cos\frac{1}{x}$
    \item $f(x) = \frac{1}{x}$.
  \end{enumerate}

  \item Let $V$ be the vector space over $\mathbb{R}$ consisting of polynomials $p(t)$ over $\mathbb{R}$ of degree less than or equal to $4$. Let $D: V \to V$ be the linear operator that takes any polynomial $p(t)$ to its derivative $p'(t)$. Then the characteristic polynomial $f(x)$ of $D$ is
  \begin{enumerate}[label=\Alph*.]
    \item $x^4$
    \item $x^5$
    \item $x^3(x - 1)$
    \item $x^4(x - 1)$.
  \end{enumerate}

  \item Let $A = \{\text{[obscured]} : a_i = 0, 1, 2, 3 \text{ or } 4\} \subset \mathbb{R}$. Then
  \begin{enumerate}[label=\Alph*.]
    \item $A$ is {[obscured in source]}
    \item {[obscured in source]}
    \item $A$ is uncountable but does not contain an open interval.
    \item $A$ contains an open interval.
  \end{enumerate}

  \item The number of group homomorphisms from $\mathbb{Z}/20\mathbb{Z}$ to $\mathbb{Z}/29\mathbb{Z}$ is
  \begin{enumerate}[label=\Alph*.]
    \item $1$
    \item $20$
    \item $29$
    \item $580$.
  \end{enumerate}

  \item Let $p(x)$ be a polynomial of degree $3$ with real coefficients. Which of the following is possible?
  \begin{enumerate}[label=\Alph*.]
    \item $p(x)$ has no real roots.
    \item $p(x)$ has exactly $2$ real roots.
    \item $p(1) = -1$, $p(2) = 1$, $p(3) = 11$ and $p(4) = 35$.
    \item $i - 1$ and $i + 1$ are roots of $p(x)$, where $i$ is the square root of $-1$.
  \end{enumerate}

  \item Let $\{a_n\}_{n=1}^\infty$ and $\{b_n\}_{n=1}^\infty$ be two sequences of real numbers such that the series $\sum_{n=1}^\infty a_n^2$ and $\sum_{n=1}^\infty b_n^2$ converge. Then the series $\sum_{n=1}^\infty a_n b_n$
  \begin{enumerate}[label=\Alph*.]
    \item is absolutely convergent.
    \item may not converge.
    \item is always convergent, but may not converge absolutely.
    \item converges to $0$.
  \end{enumerate}

  \item Let $v_i = (v_i^{(1)}, v_i^{(2)}, v_i^{(3)}, v_i^{(4)})$, for $i = 1, 2, 3, 4$, be four vectors in $\mathbb{R}^4$ such that $\sum_{i=1}^4 v_i^{(j)} = 0$, for each $j = 1, 2, 3, 4$. Let $W$ be the subspace of $\mathbb{R}^4$ spanned by $\{v_1, v_2, v_3, v_4\}$. Then the dimension of $W$ over $\mathbb{R}$ is always
  \begin{enumerate}[label=\Alph*.]
    \item either equal to $1$ or equal to $4$.
    \item less than or equal to $3$.
    \item greater than or equal to $2$.
    \item either equal to $0$ or equal to $4$.
  \end{enumerate}

  \item Let $A$ be a subset of $[0, 1]$ with non-empty interior, and let $\mathbb{Q} + A = \{q + a : q \in \mathbb{Q}, \, a \in A\}$. Which of the following is true?
  \begin{enumerate}[label=\Alph*.]
    \item $\mathbb{Q} + A = \mathbb{R}$.
    \item $\mathbb{Q} + A$ can be a proper subset of $\mathbb{R}$.
    \item $\mathbb{Q} + A$ need not be closed in $\mathbb{R}$.
    \item $\mathbb{Q} + A$ need not be open in $\mathbb{R}$.
  \end{enumerate}

  \item Let $f: \mathbb{R} \to \mathbb{R}$ be a continuously differentiable function such that $|f(x) - f(y)| \ge |x - y|$ for all $x, y \in \mathbb{R}$. Then the equation $f'(x) = \frac{1}{2}$
  \begin{enumerate}[label=\Alph*.]
    \item has exactly one solution.
    \item has no solution.
    \item has a countably infinite number of solutions.
    \item has uncountably many solutions.
  \end{enumerate}

  \item Let $f: \mathbb{R} \to [0, \infty)$ be a continuous function such that $g(x) = (f(x))^2$ is uniformly continuous. Which of the following statements is always true?
  \begin{enumerate}[label=\Alph*.]
    \item $f$ is bounded.
    \item $f$ may not be uniformly continuous.
    \item $f$ is uniformly continuous.
    \item $f$ is unbounded.
  \end{enumerate}

  \item Which of the following sequences of functions $\{f_n\}_{n=1}^\infty$ converges uniformly?
  \begin{enumerate}[label=\Alph*.]
    \item $f_n(x) = x^n$ on $[0, 1]$
    \item $f_n(x) = 1 - x^n$ on $[0, 1]$
    \item $f_n(x) = \frac{1}{1 + n x^2}$ on $[0, 1]$
    \item $f_n(x) = \frac{x}{1 + n x^2}$ on $[0, 1]$.
  \end{enumerate}

  \item Let $S$ be a collection of subsets of $\{1, 2, \dots, 100\}$ such that the intersection of any two sets in $S$ is non-empty. What is the maximum possible cardinality $|S|$ of $S$?
  \begin{enumerate}[label=\Alph*.]
    \item $100$
    \item $2^{100}$
    \item $2^{99}$
    \item $2^{98}$.
  \end{enumerate}

  \item Let $S$ be the set of all $3 \times 3$ matrices $A$ with integer entries such that the product $AA^t$ is the identity matrix. Here $A^t$ denotes the transpose of $A$. Then $|S| =$
  \begin{enumerate}[label=\Alph*.]
    \item $12$
    \item $24$
    \item $48$
    \item $60$.
  \end{enumerate}

  \item Let $A$ be a $3 \times 3$ matrix with integer entries such that $\det(A) = 1$. What is the maximum possible number of entries of $A$ that are even?
  \begin{enumerate}[label=\Alph*.]
    \item $2$
    \item $3$
    \item $6$
    \item $8$.
  \end{enumerate}

  \item The limit
  \[
    \lim_{n \to \infty} \left(\frac{1}{n} + \frac{1}{n + 1} + \dots + \frac{1}{2n}\right) =
  \]
  \begin{enumerate}[label=\Alph*.]
    \item $e$
    \item $2$
    \item $\log_e 2$
    \item $e^2$.
  \end{enumerate}

  \item Let $G = \mathbb{Z}/100\mathbb{Z}$ and let $S = \{h \in G : \operatorname{Order}(h) = 50\}$. Then $|S|$ equals
  \begin{enumerate}[label=\Alph*.]
    \item $20$
    \item $25$
    \item $30$
    \item $50$.
  \end{enumerate}
\end{enumerate}

\subsection*{Part II}
\begin{enumerate}[resume]
  \item Let $A_1 \supset A_2 \supset \cdots \supset A_n \supset A_{n+1} \supset \cdots$ be an infinite sequence of non-empty subsets of $\mathbb{R}^3$. Which of the following conditions ensures that their intersection is non-empty?
  \begin{enumerate}[label=\Alph*.]
    \item Each $A_i$ is uncountable.
    \item Each $A_i$ is open.
    \item Each $A_i$ is connected.
    \item Each $A_i$ is compact.
  \end{enumerate}

  \item Let $(X, d)$ be a metric space. Which of the following is possible?
  \begin{enumerate}[label=\Alph*.]
    \item $X$ has exactly $3$ dense subsets.
    \item $X$ has exactly $4$ dense subsets.
    \item $X$ has exactly $5$ dense subsets.
    \item $X$ has exactly $6$ dense subsets.
  \end{enumerate}

  \item Let $\{f_n\}_{n=1}^\infty$ be the sequence of functions on $\mathbb{R}$ defined by $f_n(x) = n^2 x^n$. Let $A$ be the set of all points $a$ in $\mathbb{R}$ such that the sequence $\{f_n(a)\}_{n=1}^\infty$ converges. Then
  \begin{enumerate}[label=\Alph*.]
    \item $A = \{0\}$
    \item $A = [0, 1)$
    \item $A = \mathbb{R} \setminus \{-1, 1\}$
    \item $A = (-1, 1)$.
  \end{enumerate}

  \item Let $f: \mathbb{R} \to \mathbb{R}$ be a continuous function such that $f(i) = 0$ for all $i \in \mathbb{Z}$. Which of the following statements is always true?
  \begin{enumerate}[label=\Alph*.]
    \item $\operatorname{Image}(f)$ is closed in $\mathbb{R}$.
    \item $\operatorname{Image}(f)$ is open in $\mathbb{R}$.
    \item $f$ is uniformly continuous.
    \item None of the above.
  \end{enumerate}

  \item Let $S^1 = \{z \in \mathbb{C} : |z| = 1\}$ be the unit circle. Which of the following is false? Any continuous function from $S^1$ to $\mathbb{R}$
  \begin{enumerate}[label=\Alph*.]
    \item is bounded.
    \item is uniformly continuous.
    \item has image containing a non-empty open subset of $\mathbb{R}$.
    \item has a point $z \in S^1$ such that $f(z) = f(-z)$.
  \end{enumerate}

  \item Which of the following is false?
  \begin{enumerate}[label=\Alph*.]
    \item Any continuous function from $[0, 1]$ to $[0, 1]$ has a fixed point.
    \item Any homeomorphism from $[0, 1)$ to $[0, 1)$ has a fixed point.
    \item Any bounded continuous function from $[0, \infty)$ to $[0, \infty)$ has a fixed point.
    \item Any continuous function from $(0, 1)$ to $(0, 1)$ has a fixed point.
  \end{enumerate}

  \item For $n \ge 1$, let $S_n$ denote the group of all permutations on $n$ symbols. Which of the following statements is true?
  \begin{enumerate}[label=\Alph*.]
    \item $S_3$ has an element of order $4$.
    \item $S_4$ has an element of order $6$.
    \item $S_4$ has an element of order $5$.
    \item $S_5$ has an element of order $6$.
  \end{enumerate}

  \item Which of the following rings is an integral domain?
  \begin{enumerate}[label=\Alph*.]
    \item $\mathbb{R}[x]/(x^2 + x + 1)$
    \item $\mathbb{R}[x]/(x^2 + 5x + 6)$
    \item $\mathbb{R}[x]/(x^3 - 2)$
    \item $\mathbb{R}[x]/(x^7 + 1)$.
  \end{enumerate}

  \item Let $f: \mathbb{R} \to (0, \infty)$ be a twice differentiable function such that $f(0) = 1$ and $\int_a^b f(x)\,dx = \int_a^b f'(x)\,dx$ for all $a, b \in \mathbb{R}$ with $a \le b$. Which of the following statements is false?
  \begin{enumerate}[label=\Alph*.]
    \item $f$ is one to one.
    \item The image of $f$ is compact.
    \item $f$ is unbounded.
    \item There is only one such function.
  \end{enumerate}

  \item For $X \subset \mathbb{R}^n$, consider $X$ as a metric space with metric induced by the usual Euclidean metric on $\mathbb{R}^n$. Which of the following metric spaces $X$ is complete?
  \begin{enumerate}[label=\Alph*.]
    \item $X = \mathbb{Z} \times \mathbb{Z} \subset \mathbb{R} \times \mathbb{R}$
    \item $X = \mathbb{Q} \times \mathbb{R} \subset \mathbb{R} \times \mathbb{R}$
    \item $X = (-\pi, \pi) \cap \mathbb{Q} \subset \mathbb{R}$
    \item $X = [-\pi, \pi] \cap (\mathbb{R} \setminus \mathbb{Q}) \subset \mathbb{R}$.
  \end{enumerate}
\end{enumerate}

\newpage

% ==============================================================================
\section{GS-2017 (Mathematics)}
% ==============================================================================

\begin{center}
  \textbf{TATA INSTITUTE OF FUNDAMENTAL RESEARCH}\\
  \textbf{Written Test in MATHEMATICS --- December 11, 2016}\\
  For the Ph.D.\ Programs at TIFR, Mumbai and CAM \& ICTS, Bangalore and for the Int.\ Ph.D.\ Programs at TIFR, Mumbai and CAM, Bangalore.\\
  \textbf{Duration:} Three hours (3 hours)
\end{center}

\subsection*{Instructions}
\begin{enumerate}
  \item Please read all instructions carefully before you attempt the questions.
  \item Please fill in details about name, reference code etc.\ on the answer sheet for Part I as well on the answer booklet of Part II. The Answer Sheet for Part I is machine-readable. Use only Black/Blue ball point pen to fill in the answer sheet.

  \textbf{PART I:} There are thirty (30) True/False type questions in Part I of the question paper. Allotted time for Part I is 90 minutes. The answer sheet for Part I will be collected at the end of 90 minutes. Part I questions carry $+2$ for a correct answer, $-1$ (negative marks) for a wrong answer and $0$ for not answering.

  Indicate your answer ON THE ANSWER SHEET by blackening the appropriate circle for each question. Do not mark more than one circle for any question: this will be treated as a wrong answer.

  We advise you to first mark the correct answers on the QUESTION PAPER and then to TRANSFER these to the ANSWER SHEET only when you are sure of your choice.
  \item \textbf{PART II:} 10 problems to be solved. The solutions should be written in the Answer Booklet for Part II that is provided. Extra blank sheets will be provided if needed. All Part II questions carry equal marks, and there are no negative marks. Partial credit will be given for partial solutions. Candidate can begin answering questions on Part II anytime. The answer booklet for Part II will be collected at the end of the exam.
  \item \textbf{Selection Procedure:} The answers for Part I will be machine-graded. Part I score will be used to decide a cut-off. Answer papers for Part II will be graded only for those candidates whose score is above the cut-off. List of candidates to be called for interview for the final selection for admission in the various programs will be decided based on the combined performance in Part I and II, weighted appropriately for each program.
  \item Rough work may be done on blank pages of the question paper. If needed, you may ask for extra rough sheets from an invigilator.
  \item Use of calculators, mobile phones, laptops, tablets (or other electronic devices, including those connecting to the internet) is NOT permitted.
  \item Do NOT ask for clarifications from the invigilators regarding the questions. They have been instructed not to respond to any such inquiries from candidates. In case a correction/clarification is deemed necessary, the invigilators will announce it publicly.
  \item Notation and Conventions used in this test are given below.
\end{enumerate}

\subsection*{Notation and Conventions}
\begin{itemize}
  \item $\mathbb{N}$ denotes the set of natural numbers $\{0, 1, 2, 3, \dots\}$, $\mathbb{Z}$ the set of integers, $\mathbb{Q}$ the set of rationals, $\mathbb{R}$ the set of real numbers and $\mathbb{C}$ the set of complex numbers. These sets are assumed to carry the usual algebraic and metric structures.
  \item $\mathbb{R}^n$ denotes the Euclidean space of dimension $n$. Subsets of $\mathbb{R}^n$ are assumed to carry the induced topology and metric. For a vector $v = (v_1, v_2, \dots, v_n) \in \mathbb{R}^n$ the norm $\|v\|$ is defined by $\|v\|^2 = v_1^2 + \dots + v_n^2$.
  \item $M_n(\mathbb{R})$ denotes the real vector space of $n \times n$ real matrices with the Euclidean metric.
  \item All logarithms are natural logarithms.
\end{itemize}

\subsection*{Part I}
Answer whether the following statements are True or False. Mark your answer on the machine checkable answer sheet that is provided. (Note: $+2$ marks for a correct answer, $-1$ mark for a wrong answer, $0$ marks for not answering.)

\begin{enumerate}
  \item Let $f:[0,1] \to \mathbb{R}$ be a continuous function such that $f(x) \ge x^3$ for all $x \in [0,1]$ with $\int_0^1 f(x)\,dx = \frac{1}{4}$. Then $f(x) = x^3$ for all $x \in \mathbb{R}$.
  \item Suppose $a, b, c$ are positive real numbers such that
  \[
    (1 + a + b + c)\left(1 + \frac{1}{a} + \frac{1}{b} + \frac{1}{c}\right) = 16.
  \]
  Then $a + b + c = 3$.
  \item There exists a function $f:\mathbb{R} \to \mathbb{R}$ satisfying $f(-1) = -1$, $f(1) = 1$ and $|f(x) - f(y)| \le |x - y|^{3/2}$ for all $x, y \in \mathbb{R}$.
  \item Over the real line,
  \[
    \lim_{x \to \infty} \log\left(1 + \sqrt{4 + x} - \sqrt{1 + x}\right) = \log(2).
  \]
  \item Suppose $f$ is a continuously differentiable function on $\mathbb{R}$ such that $f(x) \to 1$ and $f'(x) \to b$ as $x \to \infty$. Then $b = 1$.
  \item If $f:\mathbb{R} \to \mathbb{R}$ is differentiable and bijective, then $f^{-1}$ is also differentiable.
  \item Let $H_1, H_2, H_3, H_4$ be four hyperplanes in $\mathbb{R}^3$. The maximum possible number of connected components of $\mathbb{R}^3 - (H_1 \cup H_2 \cup H_3 \cup H_4)$ is $14$.
  \item Let $n \ge 2$ be a natural number. Let $S$ be the set of all $n \times n$ real matrices whose entries are only $0$, $1$ or $2$. Then the average determinant of a matrix in $S$ is greater than or equal to $1$.
  \item For any metric space $(X, d)$ with $X$ finite, there exists an isometric embedding $f: X \to \mathbb{R}^4$.
  \item There exists a non-negative continuous function $f:[0,1] \to \mathbb{R}$ such that $\int_0^1 f^n\,dx \to 2$ as $n \to \infty$.
  \item There exists a subset $A$ of $\mathbb{N}$ with exactly five elements such that the sum of any three elements of $A$ is a prime number.
  \item There exists a finite abelian group $G$ containing exactly $60$ elements of order $2$.
  \item Let $\alpha, \beta$ be complex numbers with non-positive real parts. Then $|e^\alpha - e^\beta| \le |\alpha - \beta|$.
  \item Every $2 \times 2$-matrix over $\mathbb{C}$ is a square of some matrix.
  \item Under the projection map $\mathbb{R}^2 \to \mathbb{R}$ sending $(x, y)$ to $x$, the image of any closed set is closed.
  \item The number of ways a $2 \times 8$ rectangle can be tiled with rectangular tiles of size $2 \times 1$ is $34$.
  \item Over the real line,
  \[
    \lim_{x \to \infty} \left(\frac{x + \log 9}{x - \log 9}\right)^x = 81.
  \]
  \item Let $f:[0,\infty) \to \mathbb{R}$ be a continuous function with $\lim_{x \to \infty} f(x) = 0$. Then $f$ has a maximum value in $[0,\infty)$.
  \item Given a continuous function $f:\mathbb{Q} \to \mathbb{Q}$, there exists a continuous function $g:\mathbb{R} \to \mathbb{R}$ such that the restriction of $g$ to $\mathbb{Q}$ is $f$.
  \item For all positive integers $m$ and $n$, if $A$ is an $m \times n$ real matrix, and $B$ is an $n \times m$ real matrix such that $AB = I$, then $BA = I$.
  \item There is a continuous onto function $f: S^2 \to S^1$ from the unit sphere in $\mathbb{R}^3$ to the unit sphere in $\mathbb{R}^2$, where $S^n = \{v \in \mathbb{R}^{n+1} \mid \|v\| = 1\}$ denotes the unit sphere in $\mathbb{R}^{n+1}$.
  \item Let $P$ be a monic, non-zero, polynomial of even degree, and $K > 0$. Then the function $P(x) - K e^x$ has a real zero.
  \item A $p$-Sylow subgroup of the underlying additive group of a finite commutative ring $R$ is an ideal in $R$.
  \item Suppose $A$ is an $n \times n$ real matrix, all whose eigenvalues have absolute value less than $1$. Then for any $v \in \mathbb{R}^n$, $\|Av\| \le \|v\|$.
  \item For any $x \in \mathbb{R}$, the sequence $\{a_n\}$, where $a_1 = x$ and $a_{n+1} = \cos(a_n)$ for all $n$, is convergent.
  \item Suppose $A_1, \dots, A_m$ are distinct $n \times n$ real matrices such that $A_i A_j = 0$ for all $i \ne j$. Then $m \le n$.
  \item In the symmetric group $S_n$ any two elements of the same order are conjugate.
  \item If a particle moving on the Euclidean line traverses distance $1$ in time $1$ starting and ending at rest, then at some time $t \in [0,1]$, the absolute value of its acceleration should be at least $4$.
  \item Let $y(t)$ be a real valued function defined on the real line such that $y' = y(1 - y)$, with $y(0) \in [0,1]$. Then $\lim_{t \to \infty} y(t) = 1$.
  \item The matrices
  \[
    \begin{pmatrix} x & 0 \\ 0 & y \end{pmatrix} \quad \text{and} \quad \begin{pmatrix} x & 1 \\ 0 & y \end{pmatrix}, \qquad x \ne y,
  \]
  for any $x, y \in \mathbb{R}$ are conjugate in $M_2(\mathbb{R})$.
\end{enumerate}

\subsection*{Part II}
Write your solutions in the answer booklet provided. All questions carry equal marks. There are no negative marks, and partial credit will be given for partial solutions.

\begin{enumerate}
  \item Show that the subset $GL_n(\mathbb{R})$ of $M_n(\mathbb{R})$ consisting of all invertible matrices is dense in $M_n(\mathbb{R})$.
  \item Let $f$ be a continuous function on $\mathbb{R}$ satisfying the relation $f(f(f(x))) = x$ for all $x \in \mathbb{R}$. Prove or disprove that $f$ is the identity function.
  \item Prove or disprove: the group of positive rationals under multiplication is isomorphic to its subgroup consisting of rationals which can be expressed as $p/q$ where both $p$ and $q$ are odd positive integers.
  \item Show that the only elements in $M_n(\mathbb{R})$ commuting with every idempotent matrix are the scalar matrices. (A matrix $P$ in $M_n(\mathbb{R})$ is said to be idempotent if $P^2 = P$.)
  \item Prove or disprove the following: let $f: X \to X$ be a continuous function from a complete metric space $(X, d)$ into itself such that $d(f(x), f(y)) < d(x, y)$ whenever $x \ne y$. Then $f$ has a fixed point.
  \item How many isomorphism classes of associative rings (with identity) are there with $35$ elements? Prove your answer.
  \item Prove or disprove: If $G$ is a finite group and $g, h \in G$, then $g, h$ have the same order if and only if there exists a group $H$ containing $G$ such that $g$ and $h$ are conjugate in $H$.
  \item Prove or disprove: there exists $A \subset \mathbb{N}$ with exactly five elements, such that the sum of any three elements of $A$ is a prime number.
  \item Show that there does not exist any continuous function $f: \mathbb{R} \to \mathbb{R}$ that takes every value exactly twice.
  \item For which positive integers $n$ does there exist an $\mathbb{R}$-linear ring homomorphism $f: \mathbb{C} \to M_n(\mathbb{R})$? Justify your answer.
\end{enumerate}

\newpage

% ==============================================================================
\section{GS-2018 (Mathematics)}
% ==============================================================================

\begin{center}
  \textbf{TATA INSTITUTE OF FUNDAMENTAL RESEARCH}\\
  \textbf{Written Test in MATHEMATICS --- December 10, 2017}\\
  For the Ph.D.\ Programs at TIFR, Mumbai and CAM \& ICTS, Bangalore and for the Int.\ Ph.D.\ Programs at TIFR, Mumbai and CAM, Bangalore.\\
  \textbf{Duration:} Two hours (2 hours)
\end{center}

\subsection*{Test Structure \& Marking Scheme}
The duration of this test is two hours. It has two parts, Part A and Part B. Part A has 25 `True or False' questions. Part B has 15 multiple choice questions. Each multiple choice question comes with four options, of which exactly one is correct. In both Part A and Part B, a correct answer gets $+1$ point, a wrong or invalid answer gets $-1$ point, and not attempting gets $0$ points.

\subsection*{Notation and Conventions}
\begin{itemize}
  \item $\mathbb{N}$ denotes the set of natural numbers $\{0, 1, 2, 3, \dots\}$, $\mathbb{Z}$ the set of integers, $\mathbb{Q}$ the set of rationals, $\mathbb{R}$ the set of real numbers and $\mathbb{C}$ the set of complex numbers. These sets are assumed to carry the usual algebraic and metric structures.
  \item $\mathbb{R}^n$ denotes the Euclidean space of dimension $n$. Subsets of $\mathbb{R}^n$ are assumed to carry the induced topology and metric.
  \item $M_n(\mathbb{R})$ denotes the real vector space of $n \times n$ real matrices with the Euclidean metric, and $I$ denotes the identity matrix.
  \item For any prime number $p$, $\mathbb{F}_p$ denotes the finite field with $p$ elements.
  \item All rings are associative, with a multiplicative identity.
  \item All logarithms are natural logarithms.
\end{itemize}

\subsection*{Part A (True/False)}
\begin{enumerate}
  \item Let $A$ be a countable subset of $\mathbb{R}$ which is well-ordered with respect to the usual ordering on $\mathbb{R}$ (where `well-ordered' means that every nonempty subset has a minimum element in it). Then $A$ has an order preserving bijection with a subset of $\mathbb{N}$.
  \item $\displaystyle \lim_{x \to 0} \frac{\sin x}{\log(1 + \tan x)} = 1$.
  \item For any closed subset $A \subset \mathbb{R}$, there exists a continuous function $f$ on $\mathbb{R}$ which vanishes exactly on $A$.
  \item Let $f$ be a nonnegative continuous function on $\mathbb{R}$ such that $\int_0^\infty f(t)\,dt$ is finite. Then $\lim_{x \to \infty} f(x) = 0$.
  \item The function $f(x) = \cos(e^x)$ is not uniformly continuous on $\mathbb{R}$.
  \item Let $A$ be a $3 \times 3$ real symmetric matrix such that $A^6 = I$. Then, $A^2 = I$.
  \item In the vector space $\{f \mid f:[0,1] \to \mathbb{R}\}$ of real-valued functions on the closed interval $[0,1]$, the set $S = \{\sin(x), \cos(x), \tan(x)\}$ is linearly independent.
  \item Let $f$ be a twice differentiable function on $\mathbb{R}$ such that both $f$ and $f''$ are strictly positive on $\mathbb{R}$. Then $\lim_{x \to \infty} f(x) = \infty$.
  \item Let $G, H$ be finite groups. Then any subgroup of $G \times H$ is equal to $A \times B$ for some subgroups $A < G$ and $B < H$.
  \item Let $g$ be a continuous function on $[0,1]$ such that $g(1) = 0$. Then the sequence of functions $f_n(x) = x^n g(x)$ converges uniformly on $[0,1]$.
  \item Let $A, B, C \in M_3(\mathbb{R})$ be such that $A$ commutes with $B$, $B$ commutes with $C$ and $B$ is not a scalar matrix. Then $A$ commutes with $C$.
  \item If $A \in M_n(\mathbb{R})$ (with $n \ge 2$) has rank $1$, then the minimal polynomial of $A$ has degree $2$.
  \item Let $V$ be the vector space over $\mathbb{R}$ consisting of polynomials of degree less than or equal to $3$. Let $T: V \to V$ be the operator sending $f(t)$ to $f(t+1)$ and $D: V \to V$ the operator sending $f(t)$ to $df(t)/dt$. Then $T$ is a polynomial in $D$.
  \item Let $V$ be the subspace of the real vector space of real valued functions on $\mathbb{R}$, spanned by $\cos t$ and $\sin t$. Let $D: V \to V$ be the linear map sending $f(t) \in V$ to $df(t)/dt$. Then $D$ has a real eigenvalue.
  \item The set of nilpotent matrices in $M_3(\mathbb{R})$ spans $M_3(\mathbb{R})$ considered as an $\mathbb{R}$-vector space (a matrix $A$ is said to be nilpotent if there exists $n \in \mathbb{N}$ such that $A^n = 0$).
  \item Let $G$ be a finite group with a normal subgroup $H$ such that $G/H$ has order $7$. Then $G \cong H \times G/H$.
  \item The multiplicative group $\mathbb{F}_7^\times$ is isomorphic to a subgroup of the multiplicative group $\mathbb{F}_{31}^\times$.
  \item Any linear transformation $A: \mathbb{R}^4 \to \mathbb{R}^4$ has a proper non-zero invariant subspace.
  \item Let $A, B \in M_n(\mathbb{R})$ be such that $A + B = AB$. Then $AB = BA$.
  \item Let $A \in M_n(\mathbb{R})$ be upper triangular with all diagonal entries $1$ such that $A \ne I$. Then $A$ is not diagonalizable.
  \item A countable group can have only countably many distinct subgroups.
  \item There exists a continuous surjection from $\mathbb{R}^3 - S^2$ to $\mathbb{R}^2 - \{(0,0)\}$ (here $S^2 \subset \mathbb{R}^3$ denotes the unit sphere defined by the equation $x^2 + y^2 + z^2 = 1$).
  \item The permutation group $S_{10}$ has an element of order $30$.
  \item Let $G$ be a finite group and $g \in G$ an element of even order. Then we can colour the elements of $G$ with two colours in such a way that $x$ and $gx$ have different colours for each $x \in G$.
  \item Let $f(x)$ and $g(x)$ be uniformly continuous functions from $\mathbb{R}$ to $\mathbb{R}$. Then their pointwise product $f(x)g(x)$ is uniformly continuous.
\end{enumerate}

\subsection*{Part B (Multiple Choice Questions)}
\begin{enumerate}
  \item The set of real numbers in the open interval $(0,1)$ which have more than one decimal expansion is
  \begin{enumerate}[label=(\alph*)]
    \item empty.
    \item non-empty but finite.
    \item countably infinite.
    \item uncountable.
  \end{enumerate}

  \item How many zeroes does the function $f(x) = e^x - 3x^2$ have in $\mathbb{R}$?
  \begin{enumerate}[label=(\alph*)]
    \item $0$
    \item $1$
    \item $2$
    \item $3$.
  \end{enumerate}

  \item Let $f: \mathbb{R} \to \mathbb{R}$ be defined as follows:
  \[
    f(x) = \begin{cases}
      1, & \text{if } x = 0, \\
      0, & \text{if } x \in \mathbb{R} \setminus \mathbb{Q}, \text{ and} \\
      \frac{1}{n}, & \text{if } x = \frac{m}{n} \text{ with } m, n \in \mathbb{Z}, \, n > 0 \text{ and } \gcd(m,n) = 1.
    \end{cases}
  \]
  Which of the following statements is true?
  \begin{enumerate}[label=(\alph*)]
    \item $f$ is continuous everywhere except at $0$.
    \item $f$ is continuous only at the irrationals.
    \item $f$ is continuous only at the non-zero rationals.
    \item $f$ is not continuous anywhere.
  \end{enumerate}

  \item Suppose $p$ is a degree $3$ polynomial such that $p(0) = 1$, $p(1) = 2$ and $p(2) = 5$. Which of the following numbers cannot equal $p(3)$?
  \begin{enumerate}[label=(\alph*)]
    \item $0$
    \item $2$
    \item $6$
    \item $10$.
  \end{enumerate}

  \item Let $A$ be the set of all functions $f: \mathbb{R} \to \mathbb{R}$ that satisfy the following two properties:
  \begin{itemize}
    \item $f$ has derivatives of all orders, and
    \item for all $x, y \in \mathbb{R}$, $f(x+y) - f(y-x) = 2x f'(y)$.
  \end{itemize}
  Which of the following sentences is true?
  \begin{enumerate}[label=(\alph*)]
    \item Any $f \in A$ is a polynomial of degree less than or equal to $1$.
    \item Any $f \in A$ is a polynomial of degree less than or equal to $2$.
    \item There exists $f \in A$ which is not a polynomial.
    \item There exists $f \in A$ which is a polynomial of degree $4$.
  \end{enumerate}

  \item Denote by $\mathfrak{A}$ the set of all $n \times n$ complex matrices $A$ ($n \ge 2$ a natural number) having the property that $4$ is the only eigenvalue of $A$. Consider the following four statements:
  \begin{itemize}
    \item $(A - 4I)^n = 0$,
    \item $A^n = 4^n I$,
    \item $(A^2 - 5A + 4I)^n = 0$,
    \item $A^n = 4nI$.
  \end{itemize}
  How many of the above statements are true for all $A \in \mathfrak{A}$?
  \begin{enumerate}[label=(\alph*)]
    \item $0$
    \item $1$
    \item $2$
    \item $3$.
  \end{enumerate}

  \item Let $A$ be the set of all continuous functions $f: [0,1] \to [0,\infty)$ satisfying the following condition:
  \[
    \int_0^x f(t)\,dt \ge f(x), \quad \text{for all } x \in [0,1].
  \]
  Then which of the following statements is true?
  \begin{enumerate}[label=(\alph*)]
    \item $A$ has cardinality $1$.
    \item $A$ has cardinality $2$.
    \item $A$ is infinite.
    \item $A$ is empty.
  \end{enumerate}

  \item Consider the following four sets of maps $f: \mathbb{Z} \to \mathbb{Q}$:
  \begin{enumerate}[label=(\roman*)]
    \item $\{f: \mathbb{Z} \to \mathbb{Q} \mid f \text{ is bijective and increasing}\}$,
    \item $\{f: \mathbb{Z} \to \mathbb{Q} \mid f \text{ is onto and increasing}\}$,
    \item $\{f: \mathbb{Z} \to \mathbb{Q} \mid f \text{ is bijective, and satisfies that } \forall n \le 0, \, f(n) \ge 0\}$, and
    \item $\{f: \mathbb{Z} \to \mathbb{Q} \mid f \text{ is onto and decreasing}\}$.
  \end{enumerate}
  How many of these sets are empty?
  \begin{enumerate}[label=(\alph*)]
    \item $0$
    \item $1$
    \item $2$
    \item $3$.
  \end{enumerate}

  \item What are the last $3$ digits of $2^{2017}$?
  \begin{enumerate}[label=(\alph*)]
    \item $072$
    \item $472$
    \item $512$
    \item $912$.
  \end{enumerate}

  \item The minimal polynomial of
  \[
    \begin{pmatrix}
      2 & 1 & 0 & 0 \\
      0 & 2 & 0 & 0 \\
      0 & 0 & 2 & 0 \\
      0 & 0 & 1 & 5
    \end{pmatrix}
  \]
  is
  \begin{enumerate}[label=(\alph*)]
    \item $(x - 2)(x - 5)$.
    \item $(x - 2)^2(x - 5)$.
    \item $(x - 2)^3(x - 5)$.
    \item none of the above.
  \end{enumerate}

  \item Consider a cube $C$ centered at the origin in $\mathbb{R}^3$. The number of invertible linear transformations of $\mathbb{R}^3$ which map $C$ onto itself is
  \begin{enumerate}[label=(\alph*)]
    \item $72$
    \item $48$
    \item $24$
    \item $12$.
  \end{enumerate}

  \item The number of rings of order $4$, up to isomorphism, is:
  \begin{enumerate}[label=(\alph*)]
    \item $1$
    \item $2$
    \item $3$
    \item $4$.
  \end{enumerate}

  \item For a sequence $\{a_n\}$ of real numbers, which of the following is a negation of the statement `$\lim_{n \to \infty} a_n = 0$'?
  \begin{enumerate}[label=(\alph*)]
    \item There exists $\epsilon > 0$ such that the set $\{n \in \mathbb{N} \mid |a_n| > \epsilon\}$ is infinite.
    \item For any $M > 0$ there exists $N \in \mathbb{N}$ such that $|a_n| > M$ for all $n \ge N$.
    \item There exists a nonzero real number $a$ such that for every $\epsilon > 0$, there exists $N \in \mathbb{N}$ with $|a_n - a| < \epsilon$ for all $n \ge N$.
    \item For any $a \in \mathbb{R}$ and every $\epsilon > 0$ there exist infinitely many $n$ such that $|a_n - a| > \epsilon$.
  \end{enumerate}

  \item Let $f: \mathbb{R} \to \mathbb{R}$ be continuous. Then which of the following statements implies that $f(0) = 0$?
  \begin{enumerate}[label=(\alph*)]
    \item $\lim_{n \to \infty} \int_0^1 f(x)^n\,dx = 0$.
    \item $\lim_{n \to \infty} \int_0^1 f(x/n)\,dx = 0$.
    \item $\lim_{n \to \infty} \int_0^1 f(nx)\,dx = 0$.
    \item None of the above.
  \end{enumerate}

  \item Consider the following maps from $\mathbb{R}^2$ to $\mathbb{R}^2$:
  \begin{enumerate}[label=(\roman*)]
    \item the map $(x, y) \mapsto (2x + 5y + 1, x + 3y)$,
    \item the map $(x, y) \mapsto (x + y^2, y + x^2)$, and
    \item the map given in polar coordinates as $(r, \theta) \mapsto (r, \theta + r^3)$ for $r \ne 0$, with the origin mapping to the origin.
  \end{enumerate}
  The number of maps in the above list that preserve areas is:
  \begin{enumerate}[label=(\alph*)]
    \item $0$
    \item $1$
    \item $2$
    \item $3$.
  \end{enumerate}
\end{enumerate}

\newpage

% ==============================================================================
\section{GS-2019 (Mathematics)}
% ==============================================================================

\subsection*{Notation and Conventions}
\begin{itemize}
  \item $\mathbb{N}$ denotes the set of natural numbers $\{0, 1, \dots\}$, $\mathbb{Z}$ the set of integers, $\mathbb{Q}$ the set of rational numbers, $\mathbb{R}$ the set of real numbers, and $\mathbb{C}$ the set of complex numbers. These sets are assumed to carry the usual algebraic and metric structures.
  \item $\mathbb{R}^n$ denotes the Euclidean space of dimension $n$. Subsets of $\mathbb{R}^n$ are viewed as metric spaces using the standard Euclidean distance on $\mathbb{R}^n$.
  \item $M_n(\mathbb{R})$ denotes the real vector space of $n \times n$ real matrices with the Euclidean metric, and $I$ denotes the identity matrix in $M_n(\mathbb{R})$.
  \item All rings are associative, with a multiplicative identity.
  \item For any prime number $p$, $\mathbb{F}_p$ denotes the finite field with $p$ elements.
  \item If $A$ and $B$ are sets, then $A - B$ refers to $\{x \in A \mid x \notin B\}$.
  \item For a ring $R$, $R[x]$ denotes the polynomial ring in one variable over $R$, and $R[x, y]$ denotes the polynomial ring in two variables over $R$.
\end{itemize}

\subsection*{Part A (Multiple Choice Questions)}
\begin{enumerate}
  \item The following sum of numbers (expressed in decimal notation)
  \[
    1 + 11 + 111 + \dots + \underbrace{11\dots 1}_{n}
  \]
  is equal to
  \begin{enumerate}[label=(\alph*)]
    \item $(10^{n+1} - 10 - 9n)/81$
    \item $(10^{n+1} - 10 + 9n)/81$
    \item $(10^{n+1} - 10 - n)/81$
    \item $(10^{n+1} - 10 + n)/81$
  \end{enumerate}

  \item For $n \ge 1$, the sequence $\{x_n\}_{n=1}^\infty$, where
  \[
    x_n = 1 + \frac{1}{\sqrt{2}} + \dots + \frac{1}{\sqrt{n}} - 2\sqrt{n},
  \]
  is
  \begin{enumerate}[label=(\alph*)]
    \item decreasing
    \item increasing
    \item constant
    \item oscillating
  \end{enumerate}

  \item Define a function:
  \[
    f(x) = \begin{cases}
      x + x^2 \cos(\pi/x), & x \ne 0, \\
      0, & x = 0.
    \end{cases}
  \]
  Consider the following statements:
  \begin{enumerate}[label=(\roman*)]
    \item $f'(0)$ exists and is equal to $1$
    \item $f$ is not increasing in any neighborhood of $0$
    \item $f'(0)$ does not exist
    \item $f$ is increasing on $\mathbb{R}$.
  \end{enumerate}
  How many of the above statements is/are true?
  \begin{enumerate}[label=(\alph*)]
    \item $0$
    \item $1$
    \item $2$
    \item $3$
  \end{enumerate}

  \item Consider differentiable functions $f: \mathbb{R} \to \mathbb{R}$ with the property that for all $a, b \in \mathbb{R}$ we have:
  \[
    f(b) - f(a) = (b - a) f'\left(\frac{a + b}{2}\right).
  \]
  Then which one of the following sentences is true?
  \begin{enumerate}[label=(\alph*)]
    \item Every such $f$ is a polynomial of degree less than or equal to $2$
    \item There exists such a function $f$ which is a polynomial of degree bigger than $2$
    \item There exists such a function $f$ which is not a polynomial
    \item Every such $f$ satisfies the condition $f\left(\frac{a+b}{2}\right) \le \frac{f(a) + f(b)}{2}$ for all $a, b \in \mathbb{R}$
  \end{enumerate}

  \item Let $V$ be an $n$-dimensional vector space and let $T: V \to V$ be a linear transformation such that $\operatorname{Rank} T = \operatorname{Rank} T^3$. Then which one of the following statements is necessarily true?
  \begin{enumerate}[label=(\alph*)]
    \item $\operatorname{Null~space}(T) = \operatorname{Range}(T)$
    \item $\operatorname{Null~space}(T) \cap \operatorname{Range}(T) = \{0\}$
    \item There exists a nonzero subspace $W$ of $V$ such that $\operatorname{Null~space}(T) \cap \operatorname{Range}(T) = W$
    \item $\operatorname{Null~space}(T) \subseteq \operatorname{Range}(T)$
  \end{enumerate}

  \item The limit
  \[
    \lim_{n \to \infty} n^2 \int_0^1 \frac{1}{(1 + x^2)^n}\,dx
  \]
  is equal to
  \begin{enumerate}[label=(\alph*)]
    \item $1$
    \item $0$
    \item $+\infty$
    \item $1/2$
  \end{enumerate}

  \item Let $A$ be an $n \times n$ matrix with rank $k$. Consider the following statements:
  \begin{enumerate}[label=(\roman*)]
    \item If $A$ has real entries, then $AA^t$ necessarily has rank $k$
    \item If $A$ has complex entries, then $AA^t$ necessarily has rank $k$.
  \end{enumerate}
  Then
  \begin{enumerate}[label=(\alph*)]
    \item (i) and (ii) are true
    \item (i) and (ii) are false
    \item (i) is true and (ii) is false
    \item (i) is false and (ii) is true
  \end{enumerate}

  \item Consider the following two statements:
  \begin{enumerate}[label=(\Alph*)]
    \item[(E)] Continuous functions on $[1,2]$ can be approximated uniformly by a sequence of even polynomials (i.e., polynomials $p(x) \in \mathbb{R}[x]$ such that $p(-x) = p(x)$).
    \item[(O)] Continuous functions on $[1,2]$ can be approximated uniformly by a sequence of odd polynomials (i.e., polynomials $p(x) \in \mathbb{R}[x]$ such that $p(-x) = -p(x)$).
  \end{enumerate}
  Choose the correct option below.
  \begin{enumerate}[label=(\alph*)]
    \item (E) and (O) are both false
    \item (E) and (O) are both true
    \item (E) is true but (O) is false
    \item (E) is false but (O) is true
  \end{enumerate}

  \item Let $f: (0,\infty) \to \mathbb{R}$ be defined by $f(x) = \frac{\sin(x^3)}{x}$. Then $f$ is
  \begin{enumerate}[label=(\alph*)]
    \item bounded and uniformly continuous
    \item bounded but not uniformly continuous
    \item not bounded but uniformly continuous
    \item not bounded and not uniformly continuous
  \end{enumerate}

  \item Let $S = \{x \in \mathbb{R} \mid x = \operatorname{Trace}(A) \text{ for some } A \in M_4(\mathbb{R}) \text{ such that } A^2 = A\}$. Then which of the following describes $S$?
  \begin{enumerate}[label=(\alph*)]
    \item $S = \{0, 2, 4\}$
    \item $S = \{0, 1/2, 1, 3/2, 2, 5/2, 3, 7/2, 4\}$
    \item $S = \{0, 1, 2, 3, 4\}$
    \item $S = [0, 4]$
  \end{enumerate}

  \item Let $f$ be a continuous function on $[0,1]$. Then the limit
  \[
    \lim_{n \to \infty} \int_0^1 n x^n f(x)\,dx
  \]
  is equal to
  \begin{enumerate}[label=(\alph*)]
    \item $f(0)$
    \item $f(1)$
    \item $\sup_{x \in [0,1]} f(x)$
    \item The limit need not exist
  \end{enumerate}

  \item Let $\{f_n\}_{n=1}^\infty$ be a sequence of functions from $\mathbb{R}$ to $\mathbb{R}$, defined by
  \[
    f_n(x) = \frac{1}{n} \exp(-n^2 x^2).
  \]
  Then which one of the following statements is true?
  \begin{enumerate}[label=(\alph*)]
    \item Both the sequences $\{f_n\}$ and $\{f_n'\}$ converge uniformly on $\mathbb{R}$
    \item Neither $\{f_n\}$ nor $\{f_n'\}$ converges uniformly on $\mathbb{R}$
    \item $\{f_n\}$ converges pointwise but not uniformly on any interval containing the origin
    \item $\{f_n'\}$ converges pointwise but not uniformly on any interval containing the origin
  \end{enumerate}

  \item Let the sequence $\{x_n\}_{n=1}^\infty$ be defined by $x_1 = \sqrt{2}$ and $x_{n+1} = (\sqrt{2})^{x_n}$ for $n \ge 1$. Then which one of the following statements is true?
  \begin{enumerate}[label=(\alph*)]
    \item The sequence $\{x_n\}$ is monotonically increasing and $\lim_{n \to \infty} x_n = 2$
    \item The sequence $\{x_n\}$ is neither monotonically increasing nor monotonically decreasing
    \item $\lim_{n \to \infty} x_n$ does not exist
    \item $\lim_{n \to \infty} x_n = \infty$
  \end{enumerate}

  \item Consider functions $f: \mathbb{R} \to \mathbb{R}$ with the property that $|f(x) - f(y)| \le 4321 |x - y|$ for all real numbers $x, y$. Then which one of the following statements is true?
  \begin{enumerate}[label=(\alph*)]
    \item $f$ is always differentiable
    \item There exists at least one such $f$ that is continuous and such that $\lim_{x \to \pm\infty} \frac{f(x)}{|x|} = \infty$
    \item There exists at least one such $f$ that is continuous, but is non-differentiable at exactly $2018$ points, and satisfies $\lim_{x \to \pm\infty} \frac{f(x)}{|x|} = 2018$
    \item It is not possible to find a sequence $\{x_n\}$ of real numbers such that $\lim_{n \to \infty} x_n = \infty$ and further satisfying $\lim_{n \to \infty} \left|\frac{f(x_n)}{x_n}\right| \le 10000$
  \end{enumerate}

  \item Let $\{f_n\}_{n=1}^\infty$ be a sequence of functions from $\mathbb{R}$ to $\mathbb{R}$, defined by
  \[
    f_n(x) = \frac{\sqrt{1 + (nx)^2}}{n}.
  \]
  Then which one of the following statements is true?
  \begin{enumerate}[label=(\alph*)]
    \item $\{f_n\}$ and $\{f_n'\}$ converge uniformly on $\mathbb{R}$
    \item $\{f_n'\}$ converges uniformly on $\mathbb{R}$ but $\{f_n\}$ does not
    \item $\{f_n\}$ converges uniformly on $\mathbb{R}$ but $\{f_n'\}$ does not
    \item $\{f_n\}$ converges uniformly to a differentiable function on $\mathbb{R}$
  \end{enumerate}

  \item The number of ring homomorphisms from $\mathbb{Z}[x, y]$ to $\mathbb{F}_2[x]/(x^3 + x^2 + x + 1)$ equals
  \begin{enumerate}[label=(\alph*)]
    \item $2^6$
    \item $2^{18}$
    \item $1$
    \item $2^9$
  \end{enumerate}

  \item Let $X \subset \mathbb{R}^2$ be the subset
  \[
    X = \{(x, y) \mid x = 0, \, |y| \le 1\} \cup \left\{(x, y) \mid 0 < x \le 1, \, y = \sin\frac{1}{x}\right\}.
  \]
  Consider the following statements:
  \begin{enumerate}[label=(\roman*)]
    \item $X$ is compact
    \item $X$ is connected
    \item $X$ is path connected.
  \end{enumerate}
  How many of the statements (i)--(iii) is/are true?
  \begin{enumerate}[label=(\alph*)]
    \item $0$
    \item $1$
    \item $2$
    \item $3$
  \end{enumerate}

  \item Consider the different ways to colour the faces of a cube with six given colours, such that each face is given exactly one colour and all the six colours are used. Define two such colouring schemes to be equivalent if the resulting configurations can be obtained from one another by a rotation of the cube. Then the number of inequivalent colouring schemes is
  \begin{enumerate}[label=(\alph*)]
    \item $15$
    \item $24$
    \item $30$
    \item $48$
  \end{enumerate}

  \item Let $C^\infty(0,1)$ stand for the set of all real-valued functions on $(0,1)$ that have derivatives of all orders. Then the map $C^\infty(0,1) \to C^\infty(0,1)$ given by $f \mapsto f + \frac{df}{dx}$ is
  \begin{enumerate}[label=(\alph*)]
    \item injective but not surjective
    \item surjective but not injective
    \item neither injective nor surjective
    \item both injective and surjective
  \end{enumerate}

  \item A stick of length $1$ is broken into two pieces by cutting at a randomly chosen point. What is the expected length of the smaller piece?
  \begin{enumerate}[label=(\alph*)]
    \item $1/8$
    \item $1/4$
    \item $1/e$
    \item $1/\pi$
  \end{enumerate}
\end{enumerate}

\subsection*{Part B (True/False)}
\begin{enumerate}
  \item There exists a continuous function $f: \mathbb{R} \to \mathbb{R}$ such that $f(\mathbb{Q}) \subseteq \mathbb{R} - \mathbb{Q}$ and $f(\mathbb{R} - \mathbb{Q}) \subseteq \mathbb{Q}$.
  \item If $A \in M_{10}(\mathbb{R})$ satisfies $A^2 + A + I = 0$, then $A$ is invertible.
  \item Let $X \subseteq \mathbb{Q}^2$. Suppose each continuous function $f: X \to \mathbb{R}^2$ is bounded. Then $X$ is necessarily finite.
  \item If $A$ is a $2 \times 2$ complex matrix that is invertible and diagonalizable, and such that $A$ and $A^2$ have the same characteristic polynomial, then $A$ is the $2 \times 2$ identity matrix.
  \item Suppose $A, B, C$ are $3 \times 3$ real matrices with $\operatorname{Rank} A = 2$, $\operatorname{Rank} B = 1$, $\operatorname{Rank} C = 2$. Then $\operatorname{Rank}(ABC) = 1$.
  \item For any $n \ge 2$, there exists an $n \times n$ real matrix $A$ such that the set $\{A^p \mid p \ge 1\}$ spans the $\mathbb{R}$-vector space $M_n(\mathbb{R})$.
  \item The matrices
  \[
    \begin{pmatrix} 0 & i & 0 \\ 0 & 0 & 1 \\ 0 & 0 & 0 \end{pmatrix} \quad \text{and} \quad \begin{pmatrix} 0 & 0 & 0 \\ -i & 0 & 0 \\ 0 & 1 & 0 \end{pmatrix}
  \]
  are similar.
  \item Consider the set $A \subset M_3(\mathbb{R})$ of $3 \times 3$ real matrices with characteristic polynomial $x^3 - 3x^2 + 2x - 1$. Then $A$ is a compact subset of $M_3(\mathbb{R}) \cong \mathbb{R}^9$.
  \item There exists an injective ring homomorphism from the product ring $\mathbb{R} \times \mathbb{R}$ into $C(\mathbb{R})$ where $C(\mathbb{R})$ denotes the ring of all continuous functions $\mathbb{R} \to \mathbb{R}$ under pointwise addition and multiplication.
  \item $\mathbb{R}$ and $\mathbb{R} \oplus \mathbb{R}$ are isomorphic as vector spaces over $\mathbb{Q}$.
  \item If $0$ is a limit point of a set $A \subseteq (0,\infty)$, then the set of all $x \in (0,\infty)$ that can be expressed as a sum of (not necessarily distinct) elements of $A$ is dense in $(0,\infty)$.
  \item The only idempotents in the ring $\mathbb{Z}_{51}$ (i.e., $\mathbb{Z}/51\mathbb{Z}$) are $0$ and $1$. (An idempotent is an element $x$ such that $x^2 = x$.)
  \item Let $A$ be a commutative ring with $1$, and let $a, b, c \in A$. Suppose there exist $x, y, z \in A$ such that $ax + by + cz = 1$. Then there exist $x', y', z' \in A$ such that $a^{50} x' + b^{20} y' + c^{15} z' = 1$.
  \item The ring $\mathbb{R}[x]/(x^5 + x - 3)$ is an integral domain.
  \item Given any group $G$ of order $12$, and any $n$ that divides $12$, there exists a subgroup $H$ of $G$ of order $n$.
  \item Let $H, N$ be subgroups of a finite group $G$, with $N$ a normal subgroup of $G$. If the orders of $G/N$ and $H$ are relatively prime, then $H$ is necessarily contained in $N$.
  \item If every proper subgroup of an infinite group $G$ is cyclic, then $G$ is cyclic.
  \item Each solution of the differential equation $y'' + e^x y = 0$ remains bounded as $x \to \infty$.
  \item There exists a uniformly continuous function $f: (0,\infty) \to (0,\infty)$ such that $\sum_{n=1}^\infty \frac{1}{f(n)}$ converges.
  \item Let $v: \mathbb{R} \to \mathbb{R}^2$ be $C^\infty$ (i.e., has derivatives of all orders). Then there exists $t_0 \in (0,1)$ such that $v(1) - v(0)$ is a scalar multiple of $\left.\frac{dv}{dt}\right|_{t=t_0}$.
\end{enumerate}

\newpage

% ==============================================================================
\section{GS-2020 (Mathematics)}
% ==============================================================================

\subsection*{Notation and Conventions}
\begin{itemize}
  \item $\mathbb{N}$ denotes the set of natural numbers $\{0, 1, \dots\}$, $\mathbb{Z}$ the set of integers, $\mathbb{Q}$ the set of rational numbers, $\mathbb{R}$ the set of real numbers, and $\mathbb{C}$ the set of complex numbers. These sets are assumed to carry the usual algebraic and metric structures.
  \item $\mathbb{R}^n$ denotes the Euclidean space of dimension $n$. Subsets of $\mathbb{R}^n$ are viewed as metric spaces using the standard Euclidean distance on $\mathbb{R}^n$.
  \item $M_n(\mathbb{R})$ denotes the real vector space of $n \times n$ real matrices, and $M_n(\mathbb{C})$ the complex vector space of $n \times n$ complex matrices. $I$ denotes the identity matrix in $M_n(\mathbb{R}) \subset M_n(\mathbb{C})$.
  \item For any $A \in M_n(\mathbb{C})$, we denote by $\operatorname{tr}(A)$ the trace of $A$ and by $\det(A)$ the determinant of $A$.
  \item All rings are associative, with a multiplicative identity.
  \item For a ring $R$, $R[x]$ denotes the polynomial ring in one variable over $R$, and $R^\times$ denotes the multiplicative group of units of $R$.
  \item All logarithms are natural logarithms.
  \item If $B$ is a subset of a set $A$, we write $A \setminus B$ for the set $\{a \in A \mid a \notin B\}$.
\end{itemize}

\subsection*{Part A (Multiple Choice Questions)}
\begin{enumerate}
  \item Consider the sequences $\{a_n\}_{n=1}^\infty$ and $\{b_n\}_{n=1}^\infty$ defined by
  \[
    a_n = (2^n + 3^n)^{1/n} \quad \text{and} \quad b_n = \frac{n}{\sum_{i=1}^n \frac{1}{a_i}}.
  \]
  What is the limit of $\{b_n\}_{n=1}^\infty$?
  \begin{enumerate}[label=(\alph*)]
    \item $2$.
    \item $3$.
    \item $5$.
    \item The limit does not exist.
  \end{enumerate}

  \item Consider the set of continuous functions $f: [0,1] \to \mathbb{R}$ that satisfy:
  \[
    \int_0^1 f(x)(1 - f(x))\,dx = \frac{1}{4}.
  \]
  Then the cardinality of this set is:
  \begin{enumerate}[label=(\alph*)]
    \item $0$.
    \item $1$.
    \item $2$.
    \item more than $2$.
  \end{enumerate}

  \item Let $f: \mathbb{R} \to \mathbb{R}$ be defined as:
  \[
    f(x) = \begin{cases}
      x^2 \sin\frac{1}{x}, & \text{if } x \ne 0, \text{ and} \\
      0, & \text{if } x = 0.
    \end{cases}
  \]
  Which of the following statements is correct?
  \begin{enumerate}[label=(\alph*)]
    \item $f$ is a surjective function.
    \item $f$ is bounded.
    \item The derivative of $f$ exists and is continuous on $\mathbb{R}$.
    \item $\{x \in \mathbb{R} \mid f(x) = 0\}$ is a finite set.
  \end{enumerate}

  \item Let $\{a_n\}_{n=1}^\infty$ be a strictly increasing bounded sequence of real numbers such that $\lim_{n \to \infty} a_n = A$. Let $f: [a_1, A] \to \mathbb{R}$ be a continuous function such that for each positive integer $i$, $f|_{[a_i, a_{i+1}]}: [a_i, a_{i+1}] \to \mathbb{R}$ is either strictly increasing or strictly decreasing. Consider the set
  \[
    B = \{M \in \mathbb{R} \mid \text{there exist infinitely many } x \in [a_1, A] \text{ such that } f(x) = M\}.
  \]
  Then the cardinality of $B$ is:
  \begin{enumerate}[label=(\alph*)]
    \item necessarily $0$.
    \item at most $1$.
    \item possibly greater than $1$, but finite.
    \item possibly infinite.
  \end{enumerate}

  \item Let $f: \mathbb{R} \to \mathbb{R}$ be a function that satisfies:
  \[
    |f(x) - f(y)| \le |x - y| |\sin(x - y)|, \quad \text{for all } x, y \in \mathbb{R}.
  \]
  Which of the following statements is correct?
  \begin{enumerate}[label=(\alph*)]
    \item $f$ is continuous but need not be uniformly continuous.
    \item $f$ is uniformly continuous but not necessarily differentiable.
    \item $f$ is differentiable, but its derivative may not be continuous.
    \item $f$ is constant.
  \end{enumerate}

  \item Let $\mathcal{C} = \{f: \mathbb{R} \to \mathbb{R} \mid f \text{ is differentiable, and } \lim_{x \to \infty} (2f(x) + f'(x)) = 0\}$. Which of the following statements is correct?
  \begin{enumerate}[label=(\alph*)]
    \item For each $L$ with $0 \ne L < \infty$, there exists $f \in \mathcal{C}$ such that $\lim_{x \to \infty} f(x) = L$.
    \item For all $f \in \mathcal{C}$, $\lim_{x \to \infty} f(x) = 0$.
    \item There exists $f \in \mathcal{C}$ such that $\lim_{x \to \infty} f(x)$ does not exist.
    \item There exists $f \in \mathcal{C}$ such that $\lim_{x \to \infty} f(x) = \frac{1}{2}$.
  \end{enumerate}

  \item Let $f(x) = \frac{\log(2+x)}{\sqrt{1+x}}$ for $x \ge 0$, and $a_m = \frac{1}{m} \int_0^m f(t)\,dt$ for every positive integer $m$. Then the sequence $\{a_m\}_{m=1}^\infty$
  \begin{enumerate}[label=(\alph*)]
    \item diverges to $+\infty$.
    \item has more than one limit point.
    \item converges and satisfies $\lim_{m \to \infty} a_m = \frac{1}{2} \log 2$.
    \item converges and satisfies $\lim_{m \to \infty} a_m = 0$.
  \end{enumerate}

  \item Let $f: \mathbb{R} \to \mathbb{R}$ be a continuous function such that $|f(x) - f(y)| \ge \log(1 + |x - y|)$ for all $x, y \in \mathbb{R}$. Then:
  \begin{enumerate}[label=(\alph*)]
    \item $f$ is injective but not surjective.
    \item $f$ is surjective but not injective.
    \item $f$ is neither injective nor surjective.
    \item $f$ is bijective.
  \end{enumerate}

  \item What is the greatest integer less than or equal to $\displaystyle \sum_{n=1}^{9999} \frac{1}{\sqrt[4]{n}}$?
  \begin{enumerate}[label=(\alph*)]
    \item $1332$
    \item $1352$
    \item $1372$
    \item $1392$
  \end{enumerate}

  \item Consider the following sentences:
  \begin{enumerate}[label=(\Roman*)]
    \item For every connected subset $Y$ of a metric space $X$, its interior $Y^\circ$ is connected.
    \item For every connected subset $Y$ of a metric space $X$, its boundary $\partial Y$ is connected.
  \end{enumerate}
  Which of the following options is correct?
  \begin{enumerate}[label=(\alph*)]
    \item (I) is true, but (II) is false.
    \item (II) is true, but (I) is false.
    \item (I) and (II) are both true.
    \item (I) and (II) are both false.
  \end{enumerate}

  \item Consider a set $\{A_1, \dots, A_n\}$ of events, $n > 1$. Suppose that one of the events in $\{A_1, \dots, A_n\}$ is certain to occur, but that no more than two of them can occur. Suppose that for each $1 \le r, s \le n$ such that $r \ne s$, the probability of $A_r$ occurring is $p$, while the probability of both $A_r$ and $A_s$ occurring is $q$. Then:
  \begin{enumerate}[label=(\alph*)]
    \item $p \le 1/n$ and $q \le 2/n$.
    \item $p \le 1/n$ and $q \ge 2/n$.
    \item $p \ge 1/n$ and $q \le 2/n$.
    \item $p \ge 1/n$ and $q \ge 2/n$.
  \end{enumerate}

  \item Let $\{z_1, z_2, \dots, z_7\}$ be a set of seven complex numbers with unit modulus. Assume that they form the vertices of a regular heptagon in the complex plane. Define $w = \sum_{i < j} z_i z_j$. Then:
  \begin{enumerate}[label=(\alph*)]
    \item $w = 0$.
    \item $|w| = \sqrt{7}$.
    \item $|w| = 7$.
    \item $|w| = 1$.
  \end{enumerate}

  \item Consider $\mathbb{R}^3$ as the space of $3 \times 1$ real matrices. The multiplicative group $GL_3(\mathbb{R})$ of invertible $3 \times 3$ real matrices acts on this space by left multiplication. What is the number of orbits for this action?
  \begin{enumerate}[label=(\alph*)]
    \item $1$.
    \item $2$.
    \item $4$.
    \item $\infty$.
  \end{enumerate}

  \item Let $V$ be a finite dimensional vector space over $\mathbb{R}$, and $W \subset V$ a subspace. Then $W \cap T(W) \ne \{0\}$ for every linear automorphism $T: V \to V$ if and only if:
  \begin{enumerate}[label=(\alph*)]
    \item $W = V$.
    \item $\dim W < \frac{1}{2} \dim V$.
    \item $\dim W = \frac{1}{2} \dim V$.
    \item $\dim W > \frac{1}{2} \dim V$.
  \end{enumerate}

  \item Let $A \in M_n(\mathbb{C})$. Then $\begin{pmatrix} A & A \\ 0 & A \end{pmatrix}$ is diagonalizable if and only if:
  \begin{enumerate}[label=(\alph*)]
    \item $A = 0$.
    \item $A = I$.
    \item $n = 2$.
    \item None of the other three options.
  \end{enumerate}

  \item Let $T: \mathbb{C} \to \mathbb{R}$ be the map defined by $T(z) = z + \overline{z}$. For a $\mathbb{C}$-vector space $V$, consider the map
  \[
    \varphi: \{f: V \to \mathbb{C} \mid f \text{ is } \mathbb{C}\text{-linear}\} \to \{g: V \to \mathbb{R} \mid g \text{ is } \mathbb{R}\text{-linear}\},
  \]
  defined by $\varphi(f) = T \circ f$. Then this map is
  \begin{enumerate}[label=(\alph*)]
    \item injective, but not surjective.
    \item surjective, but not injective.
    \item bijective.
    \item neither injective nor surjective.
  \end{enumerate}

  \item Which of the following statements is correct for every linear transformation $T: \mathbb{R}^3 \to \mathbb{R}^3$ such that $T^3 - T^2 - T + I = 0$?
  \begin{enumerate}[label=(\alph*)]
    \item $T$ is invertible as well as diagonalizable.
    \item $T$ is invertible, but not necessarily diagonalizable.
    \item $T$ is diagonalizable, but not necessarily invertible.
    \item None of the other three statements.
  \end{enumerate}

  \item Let $n \ge 2$. Which of the following statements is true for every $n \times n$ real matrix $A$ of rank one?
  \begin{enumerate}[label=(\alph*)]
    \item There exist matrices $P, Q \in M_n(\mathbb{R})$ such that all the entries of the matrix $PAQ$ are equal to $1$.
    \item There exists an invertible matrix $P \in M_n(\mathbb{R})$ such that $PAP^{-1}$ is a diagonal matrix.
    \item $A$ has a nonzero eigenvalue.
    \item The vector $(1, 1, \dots, 1) \in \mathbb{R}^n$ is an eigenvector for $A$.
  \end{enumerate}

  \item Let $m, n$ be positive integers. Then the greatest common divisor ($\gcd$) of the polynomials $x^m - 1$ and $x^n - 1$ in the ring $\mathbb{C}[x]$ equals
  \begin{enumerate}[label=(\alph*)]
    \item $x^{\min(m,n)} - 1$.
    \item $x - 1$.
    \item $x^{\gcd(m,n)} - 1$.
    \item None of the other three options.
  \end{enumerate}

  \item Let $A_4$ denote the group of even permutations of $\{1, 2, 3, 4\}$. Consider the following statements:
  \begin{enumerate}[label=(\Roman*)]
    \item There exists a surjective group homomorphism $A_4 \to \mathbb{Z}/4\mathbb{Z}$.
    \item There exists a surjective group homomorphism $A_4 \to \mathbb{Z}/3\mathbb{Z}$.
  \end{enumerate}
  Which of the following statements is correct?
  \begin{enumerate}[label=(\alph*)]
    \item (I) is true and (II) is false.
    \item (II) is true and (I) is false.
    \item (I) and (II) are both true.
    \item (I) and (II) are both false.
  \end{enumerate}
\end{enumerate}

\subsection*{Part B (True/False)}
\begin{enumerate}
  \item There exists no monotone function $f: \mathbb{R} \to \mathbb{R}$ which is discontinuous at every rational number.
  \item Let $C([0,1])$ denote the set of continuous real valued functions on $[0,1]$, and $\mathbb{R}^{\mathbb{N}}$ the set of all sequences of real numbers. Then there exists an injective map from $C([0,1])$ to $\mathbb{R}^{\mathbb{N}}$.
  \item Let $\{a_n\}_{n=1}^\infty$ be a bounded sequence of positive real numbers. Then:
  \[
    \limsup_{n \to \infty} \frac{1}{a_n} = \frac{1}{\liminf_{n \to \infty} a_n}.
  \]
  \item Let $C([0,1])$ denote the metric space of continuous real valued functions on $[0,1]$ under the supremum metric, i.e., the distance between $f$ and $g$ in $C([0,1])$ equals $\sup\{|f(x) - g(x)| \mid x \in [0,1]\}$. Let $Q \subset C([0,1])$ be the set of all polynomials in $\mathbb{R}[x]$ in which the coefficient of $x^2$ is $0$. Then $Q$ is dense in $C([0,1])$.
  \item If $X$ is a metric space such that every continuous function $f: X \to \mathbb{R}$ is uniformly continuous, then $X$ is compact.
  \item Let $X$ be a metric space, and let $C(X)$ denote the $\mathbb{R}$-vector space of continuous real valued functions on $X$. Then $X$ is infinite if and only if $\dim_{\mathbb{R}} C(X) = \infty$.
  \item Let $A$ be a countable union of lines in $\mathbb{R}^3$. Then $\mathbb{R}^3 \setminus A$ is connected.
  \item An invertible linear map from $\mathbb{R}^2$ to itself takes parallel lines to parallel lines.
  \item For any matrix $C$ with entries in $\mathbb{C}$, let $m(C)$ denote the minimal polynomial of $C$, and $p(C)$ its characteristic polynomial. Then for any $n \in \mathbb{N}$ two matrices $A, B \in M_n(\mathbb{C})$ are similar if and only if $m(A) = m(B)$ and $p(A) = p(B)$.
  \item Let $A, B \in M_3(\mathbb{R})$. Then
  \[
    \det(AB - BA) = \frac{\operatorname{tr}[(AB - BA)^3]}{3}.
  \]
  \item There exist an integer $r \ge 1$ and a symmetric matrix $A \in M_r(\mathbb{R})$ such that for all $n \in \mathbb{N}$, we have:
  \[
    2^{\sqrt{n}} \le |\operatorname{tr}(A^n)| \le 2020 \cdot 2^{\sqrt{n}}.
  \]
  \item The polynomial $1 + x + \frac{x^2}{2!} + \dots + \frac{x^{101}}{101!}$ is irreducible in $\mathbb{Q}[x]$.
  \item There exists an integer $n > 3$ such that the group of units of the ring $\mathbb{Z}/2^n\mathbb{Z}$ is cyclic.
  \item For every surjective ring homomorphism $\varphi: R \to S$ we have $\varphi(R^\times) = S^\times$.
  \item Let $G$ be a finite group and $P$ a $p$-Sylow subgroup of $G$, where $p$ is a prime number. Then for every subgroup $H$ of $G$, $H \cap P$ is a $p$-Sylow subgroup of $H$.
  \item Let $G$ be an abelian group, with identity element $e$. If $\{g \in G \mid g = e \text{ or } g \text{ has infinite order}\}$ is a subgroup of $G$, then either all elements of $G \setminus \{e\}$ have infinite order, or all elements of $G$ have finite order.
  \item There exists a natural number $n$, with $1 < n \le 10$, such that $x^n$ and $x$ are conjugate for every element $x$ of $S_7$, the group of permutations of $\{1, \dots, 7\}$.
  \item Every noncommutative ring has at least $10$ elements.
  \item Let $\{a_n\}_{n=1}^\infty$ be a sequence of elements in $\{0, 1\}$ such that for all positive integers $n$, $\sum_{i=n}^{n+9} a_i$ is divisible by $3$. Then there exists a positive integer $k$ such that $a_{n+k} = a_n$ for all positive integers $n$.
  \item The interior of any strip bounded by two parallel lines in $\mathbb{R}^2$, of width strictly greater than $1$, contains a point with integer coordinates.
\end{enumerate}

\newpage

% ==============================================================================
\section{GS-2021 (Mathematics)}
% ==============================================================================

\subsection*{Notation and Conventions}
\begin{itemize}
  \item $\mathbb{N}$ denotes the set of natural numbers $\{0, 1, \dots\}$, $\mathbb{Z}$ the set of integers, $\mathbb{Q}$ the set of rational numbers, $\mathbb{R}$ the set of real numbers, and $\mathbb{C}$ the set of complex numbers. These sets are assumed to carry the usual algebraic and metric structures.
  \item $\mathbb{R}^n$ denotes the Euclidean space of dimension $n$. Subsets of $\mathbb{R}^n$ are viewed as metric spaces using the standard Euclidean distance on $\mathbb{R}^n$.
  \item $M_n(\mathbb{R})$ denotes the real vector space of $n \times n$ real matrices, and $M_n(\mathbb{C})$ the complex vector space of $n \times n$ complex matrices. $M_n(\mathbb{R})$ gets the topology transferred from any $\mathbb{R}$-linear isomorphism $M_n(\mathbb{R}) \cong \mathbb{R}^{n^2}$, and its subsets get the subspace topology.
  \item $\operatorname{Id}$ denotes the identity matrix in $M_n(\mathbb{R}) \subset M_n(\mathbb{C})$.
  \item A matrix $A \in M_n(\mathbb{R})$ is called idempotent if $A^2 = A$, and nilpotent if $A^m = 0$ for some positive integer $m$.
  \item All rings are associative, with a multiplicative identity.
  \item For a ring $R$, $R[x]$ denotes the polynomial ring in one variable over $R$, and $R^\times$ denotes the multiplicative group of units of $R$.
  \item All logarithms are natural logarithms.
  \item If $B$ is a subset of a set $A$, we write $A \setminus B$ for the set $\{a \in A \mid a \notin B\}$.
  \item An isometry from a metric space $(X, d)$ to a metric space $(X', d')$ is a (not necessarily surjective) map $f: X \to X'$ such that for all $a, b \in X$, we have $d'(f(a), f(b)) = d(a, b)$.
\end{itemize}

\subsection*{Part A (Multiple Choice Questions)}
\begin{enumerate}
  \item For each positive integer $n$, let
  \[
    s_n = \frac{1}{\sqrt{4n^2 - 1^2}} + \frac{1}{\sqrt{4n^2 - 2^2}} + \dots + \frac{1}{\sqrt{4n^2 - n^2}}.
  \]
  Then $\lim_{n \to \infty} s_n$ equals
  \begin{enumerate}[label=(\alph*)]
    \item $\pi/2$
    \item $\pi/6$
    \item $1/2$
    \item $\infty$
  \end{enumerate}

  \item The number of bijective maps $g: \mathbb{N} \to \mathbb{N}$ such that $\sum_{n=1}^\infty \frac{g(n)}{n^2} < \infty$ is
  \begin{enumerate}[label=(\alph*)]
    \item $0$
    \item $1$
    \item $2$
    \item $\infty$
  \end{enumerate}

  \item The value of $\displaystyle \lim_{n \to \infty} \prod_{k=2}^n \left(1 - \frac{1}{k^2}\right)$ is
  \begin{enumerate}[label=(\alph*)]
    \item $1/2$
    \item $1$
    \item $1/4$
    \item $0$
  \end{enumerate}

  \item The set $S = \{x \in \mathbb{R} \mid x > 0 \text{ and } (1 + x^2)\tan(2x) = x\}$ is
  \begin{enumerate}[label=(\alph*)]
    \item empty
    \item nonempty but finite
    \item countably infinite
    \item uncountable
  \end{enumerate}

  \item The dimension of the real vector space
  \[
    V = \{f: (-1, 1) \to \mathbb{R} \mid f \text{ is infinitely differentiable on } (-1, 1) \text{ and } f^{(n)}(0) = 0 \text{ for all } n \ge 0\}
  \]
  is
  \begin{enumerate}[label=(\alph*)]
    \item $0$
    \item $1$
    \item greater than one, but finite
    \item infinite
  \end{enumerate}

  \item For a positive integer $n$, let $a_n$ denote the unique positive real root of $x^n + x^{n-1} + \dots + x - 1 = 0$. Then
  \begin{enumerate}[label=(\alph*)]
    \item the sequence $\{a_n\}_{n=1}^\infty$ is unbounded
    \item $\lim_{n \to \infty} a_n = 0$
    \item $\lim_{n \to \infty} a_n = 1/2$
    \item $\lim_{n \to \infty} a_n$ does not exist
  \end{enumerate}

  \item Let $A$ be the set of all real numbers $\lambda \in [0,1]$ such that
  \[
    \lim_{p \to 0} \frac{\log(\lambda 2^p + (1 - \lambda) 3^p)}{p} = \lambda \log 2 + (1 - \lambda) \log 3.
  \]
  Then
  \begin{enumerate}[label=(\alph*)]
    \item $A = \{0, 1\}$
    \item $A = \{0, \frac{1}{2}, 1\}$
    \item $A = \{0, \frac{1}{3}, \frac{1}{2}, \frac{2}{3}, 1\}$
    \item $A = [0, 1]$
  \end{enumerate}

  \item Let $X \subseteq \mathbb{R}$ be a subset. Let $\{f_n\}_{n=1}^\infty$ be a sequence of functions $f_n: X \to \mathbb{R}$ that converges uniformly to a function $f: X \to \mathbb{R}$. For each positive integer $n$, let $D_n \subseteq X$ denote the set of points at which $f_n$ is not continuous. Let $D \subseteq X$ denote the set of points at which $f$ is not continuous. Which one of the following statements is correct?
  \begin{enumerate}[label=(\alph*)]
    \item If each $D_n$ is finite, then $D$ is finite.
    \item If each $D_n$ has at most $7$ elements, then $D$ has at most $7$ elements.
    \item If each $D_n$ is uncountable, then $D$ is uncountable.
    \item None of the other three statements is correct.
  \end{enumerate}

  \item Let $f: \mathbb{R} \to \mathbb{R}$ be an arbitrary function. Consider the following assertions:
  \begin{enumerate}[label=(\Roman*)]
    \item $f$ is continuous.
    \item The set $\operatorname{Graph}(f) = \{(x, f(x)) \in \mathbb{R}^2 \mid x \in \mathbb{R}\}$ is a connected subset of $\mathbb{R}^2$.
  \end{enumerate}
  Which one of the following statements is correct?
  \begin{enumerate}[label=(\alph*)]
    \item (I) implies (II) but (II) does not imply (I).
    \item (II) implies (I) but (I) does not imply (II).
    \item (I) implies (II) and (II) implies (I).
    \item (I) does not imply (II), and (II) does not imply (I).
  \end{enumerate}

  \item Let $\mathcal{C}$ denote the set of colorings of an $8 \times 8$ chessboard, where each square is colored either black or white. Let $\sim$ denote the equivalence relation on $\mathcal{C}$ defined as follows: two colorings are equivalent if and only if one of them can be obtained from the other by a rotation of the chessboard. The cardinality of the set $\mathcal{C}/\sim$ of equivalence classes of elements of $\mathcal{C}$ under $\sim$ is
  \begin{enumerate}[label=(\alph*)]
    \item $2^{62}$
    \item $2^{62} + 2^{30} + 2^{15}$
    \item $2^{64} - 2^{32} + 2^{16}$
    \item $2^{63} - 2^{31} + 2^{15}$
  \end{enumerate}

  \item What is the number of surjective maps from the set $\{1, \dots, 10\}$ to the set $\{1, 2\}$?
  \begin{enumerate}[label=(\alph*)]
    \item $90$
    \item $1022$
    \item $98$
    \item $1024$
  \end{enumerate}

  \item Let $V$ be a vector space over a field $F$. Consider the following assertions:
  \begin{enumerate}[label=(\Roman*)]
    \item $V$ is finite dimensional.
    \item For every linear transformation $T: V \to V$ there exists a nonzero polynomial $p(x) \in F[x]$ such that $p(T): V \to V$ is the zero map.
  \end{enumerate}
  Which one of the following statements is correct?
  \begin{enumerate}[label=(\alph*)]
    \item (I) implies (II) but (II) does not imply (I).
    \item (II) implies (I) but (I) does not imply (II).
    \item (I) implies (II) and (II) implies (I).
    \item (I) does not imply (II), and (II) does not imply (I).
  \end{enumerate}

  \item Let $T: \mathbb{C}[x] \to \mathbb{C}[x]$ be the $\mathbb{C}$-linear transformation defined on the complex vector space $\mathbb{C}[x]$ of one variable complex polynomials by $Tf(x) = f(x+1)$. How many eigenvalues does $T$ have?
  \begin{enumerate}[label=(\alph*)]
    \item $1$
    \item finite but more than $1$
    \item countably infinite
    \item uncountable
  \end{enumerate}

  \item Let $\mathbb{R}^{\mathbb{N}}$ denote the real vector space of sequences $(x_0, x_1, x_2, \dots)$ of real numbers. Define a linear transformation $T: \mathbb{R}^{\mathbb{N}} \to \mathbb{R}^{\mathbb{N}}$ by $(x_0, x_1, \dots) \mapsto (x_0 + x_1, x_1 + x_2, \dots)$. Which one of the following statements is correct?
  \begin{enumerate}[label=(\alph*)]
    \item The kernel of $T$ is infinite dimensional.
    \item The image of $T$ is infinite dimensional.
    \item The quotient vector space $\mathbb{R}^{\mathbb{N}}/T(\mathbb{R}^{\mathbb{N}})$ is infinite dimensional.
    \item None of the other three statements is correct.
  \end{enumerate}

  \item Which one of the following statements is correct?
  \begin{enumerate}[label=(\alph*)]
    \item There exists a $\mathbb{C}$-linear isomorphism $\mathbb{C}^2 \to \mathbb{C}$.
    \item There exists no $\mathbb{C}$-linear isomorphism $\mathbb{C}^2 \to \mathbb{C}$, but there exists an $\mathbb{R}$-linear isomorphism $\mathbb{C}^2 \to \mathbb{C}$.
    \item There exists no $\mathbb{R}$-linear isomorphism $\mathbb{C}^2 \to \mathbb{C}$, but there exists a $\mathbb{Q}$-linear isomorphism $\mathbb{C}^2 \to \mathbb{C}$.
    \item None of the other three statements is correct.
  \end{enumerate}

  \item The matrix
  \[
    \begin{pmatrix}
      4 & -3 & -3 \\
      3 & -2 & -3 \\
      -1 & 1 & 2
    \end{pmatrix}
  \]
  is
  \begin{enumerate}[label=(\alph*)]
    \item diagonalizable
    \item nilpotent
    \item idempotent
    \item none of the other three options
  \end{enumerate}

  \item Which of the following is a necessary and sufficient condition for two real $3 \times 3$ matrices $A$ and $B$ to be similar (i.e., $PAP^{-1} = B$ for an invertible real $3 \times 3$ matrix $P$)?
  \begin{enumerate}[label=(\alph*)]
    \item They have the same characteristic polynomial.
    \item They have the same minimal polynomial.
    \item They have the same minimal and characteristic polynomials.
    \item None of the other three conditions.
  \end{enumerate}

  \item Consider the following two subgroups $A, B$ of the group $\mathbb{Q}[x]$ of one variable rational polynomials under addition:
  \begin{align*}
    A &= \{p(x) \in \mathbb{Z}[x] \mid p \text{ has degree at most } 2\}, \text{ and} \\
    B &= \{p(x) \in \mathbb{Q}[x] \mid p \text{ has degree at most } 2, \text{ and } p(\mathbb{Z}) \subseteq \mathbb{Z}\}.
  \end{align*}
  Then the index $[B : A]$ of $A$ in $B$ equals
  \begin{enumerate}[label=(\alph*)]
    \item $1$
    \item $2$
    \item $4$
    \item none of the other three options
  \end{enumerate}

  \item Let $G$ be any finite group of order $2021$. For which of the following positive integers $m$ is the map $G \to G$, given by $g \mapsto g^m$, a bijection?
  \begin{enumerate}[label=(\alph*)]
    \item $43$
    \item $45$
    \item $47$
    \item none of the other three options
  \end{enumerate}

  \item How many subgroups does $(\mathbb{Z}/13\mathbb{Z}) \times (\mathbb{Z}/13\mathbb{Z})$ have?
  \begin{enumerate}[label=(\alph*)]
    \item $13$
    \item $16$
    \item $4$
    \item $25$
  \end{enumerate}
\end{enumerate}

\subsection*{Part B (True/False)}
\begin{enumerate}
  \item Let $f_n: [0,1] \to \mathbb{R}$ be a continuous function for each positive integer $n$. If $\lim_{n \to \infty} \int_0^1 f_n(x)^2\,dx = 0$, then $\lim_{n \to \infty} f_n\left(\frac{1}{2}\right) = 0$.
  \item Let $(X, d)$ be an infinite compact metric space. Then there exists no function $f: X \to X$, continuous or otherwise, with the property that $d(f(x), f(y)) > d(x, y)$ for all $x \ne y$.
  \item Every infinite closed subset of $\mathbb{R}^n$ is the closure of a countable set.
  \item If $X$ is a compact metric space, there exists a surjective (not necessarily continuous) function $\mathbb{R} \to X$.
  \item If $X$ is a compact metric space, then every isometry $f: X \to X$ is surjective.
  \item Define a metric on the set of finite subsets of $\mathbb{Z}$ as follows: $d(A, B) = \text{the cardinality of } (A \cup B) \setminus (A \cap B)$. The resulting metric space admits an isometry into $\mathbb{R}^n$, for some positive integer $n$.
  \item There exists a continuous function $f: [0,1] \to \{A \in M_2(\mathbb{R}) \mid A^2 = A\}$ such that $f(0) = 0$ and $f(1) = \operatorname{Id}$.
  \item Let $f: [0,1] \to \mathbb{R}$ be a monotone increasing (not necessarily continuous) function such that $f(0) > 0$ and $f(1) < 1$. Then there exists $x \in [0,1]$ such that $f(x) = x$.
  \item The set $\{(x, y) \in \mathbb{N} \times \mathbb{N} \mid x \text{ divides } y^x, \, x \ne y, \, xy \ne 0, \, x \ne 1\}$ is finite.
  \item Suppose a line segment of a fixed length $L$ is given. It is possible to construct a triangle of perimeter $L$, whose angles are $105^\circ$, $45^\circ$ and $30^\circ$, using only a straight edge and a compass.
  \item The real vector space $M_n(\mathbb{R})$ cannot be spanned by nilpotent matrices, for any positive integer $n$.
  \item Let $S \subseteq M_n(\mathbb{R})$ be a nonempty finite set closed under matrix multiplication. Then there exists $A \in S$ such that the trace of $A$ is an integer.
  \item Given a linear transformation $T: \mathbb{Q}^4 \to \mathbb{Q}^4$, there exists a nonzero proper subspace $V$ of $\mathbb{Q}^4$ such that $T(V) \subseteq V$.
  \item If $G$ is a finite group such that the group $\operatorname{Aut}(G)$ of automorphisms of $G$ is cyclic, then $G$ is abelian.
  \item There exists a countable group having uncountably many subgroups.
  \item There exists a nonzero ideal $I \subseteq \mathbb{Z}[i]$ such that the quotient ring $\mathbb{Z}[i]/I$ is infinite (here $i$ is a square root of $-1$ in $\mathbb{C}$).
  \item There exists an injective ring homomorphism from the ring $\mathbb{Q}[x, y]/(x^2 - y^2)$ into the ring $\mathbb{Q}[x, y]/(x - y^2)$.
  \item The set $\{n \in \mathbb{N} \mid n \text{ divides } a^3 - 1, \text{ for all integers } a \text{ such that } \gcd(a, n) = 1\}$ is infinite.
  \item The set of polynomials in the ring $\mathbb{Z}[x]$, the sum of whose coefficients is zero, forms an ideal of the ring $\mathbb{Z}[x]$.
  \item Let $c_1, c_2 > 0$, and let $f, g: \mathbb{R} \to \mathbb{R}$ be functions (not assumed to be continuous) such that for all $x \in \mathbb{R}$, $f(x + c_1) = f(x)$ and $g(x + c_2) = g(x)$. Further, assume that $\lim_{x \to \infty} (f(x) - g(x)) = 0$. Then $f = g$.
\end{enumerate}

\newpage

% ==============================================================================
\section{GS-2022 (Mathematics)}
% ==============================================================================

\subsection*{Notation and Conventions}
\begin{itemize}
  \item $\mathbb{N}$ denotes the set of natural numbers $\{0, 1, \dots\}$, $\mathbb{Z}$ the set of integers, $\mathbb{Q}$ the set of rational numbers, $\mathbb{R}$ the set of real numbers, and $\mathbb{C}$ the set of complex numbers. These sets are assumed to carry the usual algebraic and metric structures.
  \item $\mathbb{R}^n$ denotes the Euclidean space of dimension $n$. Subsets of $\mathbb{R}^n$ are viewed as metric spaces using the standard Euclidean distance on $\mathbb{R}^n$.
  \item $M_n(\mathbb{R})$ denotes the real vector space of $n \times n$ real matrices, and $M_n(\mathbb{C})$ the complex vector space of $n \times n$ complex matrices. $M_n(\mathbb{R})$ gets the topology transferred from any $\mathbb{R}$-linear isomorphism $M_n(\mathbb{R}) \cong \mathbb{R}^{n^2}$, and its subsets get the subspace topology.
  \item $\operatorname{Id}$ denotes the identity matrix in $M_n(\mathbb{R}) \subset M_n(\mathbb{C})$.
  \item For a metric space $X$, $C(X, \mathbb{R})$ denotes the set of continuous functions from $X$ to $\mathbb{R}$, viewed as a ring under pointwise addition and multiplication. Similarly, $C(X, \mathbb{C})$ will denote the ring of continuous functions from $X$ to $\mathbb{C}$, again with pointwise addition and multiplication.
  \item A matrix $T \in M_n(\mathbb{C})$, or a linear transformation $T: V \to V$ from a vector space $V$ to itself, is called idempotent if $T^2 = T$, and nilpotent if $T^m = 0$ for some positive integer $m$.
  \item All rings are associative, with a multiplicative identity. A subring of a ring $R$ is assumed by definition to contain the multiplicative identity of $R$.
  \item If $q$ is a power of a prime number, $\mathbb{F}_q$ will stand for the finite field with $q$ elements.
  \item For a ring $R$, $R[x_1, \dots, x_n]$ denotes the polynomial ring in $n$ variables $x_1, \dots, x_n$ over $R$, and $R^\times$ denotes the multiplicative group of units of $R$.
  \item All logarithms are natural logarithms.
  \item If $B$ is a subset of a set $A$, we write $A \setminus B$ for the set $\{a \in A \mid a \notin B\}$.
  \item By the group of isometries of a metric space $(X, d)$, we mean the group whose elements are bijective maps $f: X \to X$ satisfying $d(x, y) = d(f(x), f(y))$ for all $x, y \in X$, and whose multiplication is defined by composition.
  \item If $f$ is a real or complex valued function defined on an open interval $(a, b)$ of $\mathbb{R}$, then $f$ is said to be continuously differentiable on $(a, b)$ if its derivative exists and is continuous on $(a, b)$.
  \item If $f$ is a real or complex valued function defined on an open interval of $\mathbb{R}$, then $f'$ will stand for the first derivative of $f$ (wherever it exists), and $f''$ for the second derivative of $f$ (wherever it exists).
\end{itemize}

\subsection*{Part A (Multiple Choice Questions)}
\begin{enumerate}
  \item Consider the following properties of a metric space $(X, d)$:
  \begin{enumerate}[label=(\Roman*)]
    \item $(X, d)$ is complete as a metric space.
    \item For any sequence $\{Z_n\}_{n \in \mathbb{N}}$ of closed nonempty subsets of $X$, such that $Z_1 \supseteq Z_2 \supseteq \cdots$ and $\lim_{n \to \infty} (\sup_{x, y \in Z_n} d(x, y)) = 0$, $\bigcap_{n=1}^\infty Z_n$ is a singleton set.
  \end{enumerate}
  Which of the following sentences is true?
  \begin{enumerate}[label=(\alph*)]
    \item (I) implies (II) and (II) implies (I).
    \item (I) implies (II) but (II) does not imply (I).
    \item (I) does not imply (II) but (II) implies (I).
    \item (I) does not imply (II) and (II) does not imply (I).
  \end{enumerate}

  \item Consider the following assertions:
  \begin{enumerate}[label=(\Roman*)]
    \item $\{(x, y) \in \mathbb{R}^2 \mid xy = 1\}$ is connected.
    \item $\{(x, y) \in \mathbb{C}^2 \mid xy = 1\}$ is connected.
  \end{enumerate}
  Which of the following sentences is true?
  \begin{enumerate}[label=(\alph*)]
    \item Both (I) and (II) are true.
    \item (I) is true but (II) is false.
    \item (I) is false but (II) is true.
    \item Both (I) and (II) are false.
  \end{enumerate}

  \item What is the number of solutions of $x = \frac{x^2}{50} - \cos\frac{x}{2} + 2$ in $[0, 10]$?
  \begin{enumerate}[label=(\alph*)]
    \item $0$
    \item $1$
    \item $2$
    \item $\infty$
  \end{enumerate}

  \item Let $A$ be an element of $M_4(\mathbb{R})$ with characteristic polynomial $t^4 - t$. What is the characteristic polynomial of $A^2$?
  \begin{enumerate}[label=(\alph*)]
    \item $t^4 - t$
    \item $t^4 - 2t^3 + t^2$
    \item $t^4 - t^2$
    \item None of the other three options
  \end{enumerate}

  \item Let $n$ be a positive integer, and let $V = \{f \in \mathbb{R}[x] \mid \deg f \le n\}$ be the real vector space of real polynomials of degree at most $n$. Let $\operatorname{End}_{\mathbb{R}}(V)$ denote the real vector space of linear transformations from $V$ to itself. For $m \in \mathbb{Z}$, let $T_m \in \operatorname{End}_{\mathbb{R}}(V)$ be such that $(T_m f)(x) = f(x + m)$ for all $f \in V$. Then the dimension of the vector subspace of $\operatorname{End}_{\mathbb{R}}(V)$ given by $\operatorname{Span}(\{T_m \mid m \in \mathbb{Z}\})$ is
  \begin{enumerate}[label=(\alph*)]
    \item $1$
    \item $n$
    \item $n + 1$
    \item $n^2$
  \end{enumerate}

  \item Let $T$ be the linear transformation from the real vector space $\mathbb{R}[x]$ to itself, given by $T(f) = f'$, where $f'$ is the derivative of $f$. Consider the following statements about $T$:
  \begin{enumerate}[label=(\Roman*)]
    \item $T$ is nilpotent.
    \item The only eigenvalue of $T$ is $0$.
  \end{enumerate}
  Which of the following sentences is true?
  \begin{enumerate}[label=(\alph*)]
    \item Both (I) and (II) are true.
    \item (I) is true but (II) is false.
    \item (I) is false but (II) is true.
    \item Both (I) and (II) are false.
  \end{enumerate}

  \item What is the cardinality of the set of $\theta \in [0, 2\pi)$ such that the linear map $\mathbb{R}^2 \to \mathbb{R}^2$ given by the matrix
  \[
    \begin{pmatrix} \cos\theta & -\sin\theta \\ \sin\theta & \cos\theta \end{pmatrix}
  \]
  has an eigenvector in $\mathbb{R}^2$?
  \begin{enumerate}[label=(\alph*)]
    \item $1$
    \item $2$
    \item $4$
    \item $\infty$
  \end{enumerate}

  \item Let $p$ be a prime number, and let $A$ equal $\begin{pmatrix} 2 & -1 \\ 1 & 0 \end{pmatrix}$ viewed as a $2 \times 2$ matrix with integer entries. What is the smallest positive integer $n$ such that the matrix $A^n$ is congruent to the $2 \times 2$ identity matrix modulo $p$?
  \begin{enumerate}[label=(\alph*)]
    \item $p^2 - 1$
    \item $p - 1$
    \item $p$
    \item $p + 1$
  \end{enumerate}

  \item What is the largest value of $n$ for which there exists a set $\{A_1, \dots, A_n\}$ of (distinct) nonzero matrices in $M_2(\mathbb{C})$ such that $A_i^* A_j$ has trace zero for all $1 \le i < j \le n$?
  \begin{enumerate}[label=(\alph*)]
    \item $1$
    \item Greater than $1$ but at most $4$
    \item Greater than $4$ but finite
    \item $\infty$
  \end{enumerate}

  \item Let $p$ be a prime number. What is the number of elements in the group $\mathbb{Q}/\mathbb{Z}$ that have order exactly $p$?
  \begin{enumerate}[label=(\alph*)]
    \item $0$
    \item $p - 1$
    \item $p$
    \item $\infty$
  \end{enumerate}

  \item Consider the real polynomial $f(x) = x^{11} - x^7 + x^2 - 1$. Which of the following sentences is correct?
  \begin{enumerate}[label=(\alph*)]
    \item $f(x)$ has exactly one positive root.
    \item $f(x)$ has exactly two positive roots.
    \item $f(x)$ has at least three positive roots.
    \item None of the other three options.
  \end{enumerate}

  \item Consider polynomials $f_1(x, y) = \sum_{i,j=0}^\infty a_{ij} x^i y^j$ and $f_2(x, y) = \sum_{i,j=0}^\infty b_{ij} x^i y^j \in \mathbb{R}[x, y]$ (where $a_{ij} = b_{ij} = 0$ for all but finitely many $(i, j) \in \mathbb{N}^2$), such that $f_1(p, q) = f_2(p, q)$ for all $(p, q) \in \mathbb{R}^2$ satisfying $p^2 = q^2$. Which of the following sentences is true for all such $f_1$ and $f_2$?
  \begin{enumerate}[label=(\alph*)]
    \item $a_{00} = b_{00}$, but we may not have $a_{ij} = b_{ij}$ for all $(i, j)$ with $i + j = 1$.
    \item $a_{ij} = b_{ij}$ if $i + j \le 1$, but we may not have $a_{ij} = b_{ij}$ for all $(i, j)$ with $i + j = 2$.
    \item $a_{ij} = b_{ij}$ if $i + j \le 2$, but we may not have $a_{ij} = b_{ij}$ for all $(i, j)$ with $i + j = 3$.
    \item $a_{ij} = b_{ij}$ if $i + j \le 3$, but we may not have $f_1 = f_2$.
  \end{enumerate}

  \item What is the number of bijections $f: \{1, 2, \dots, 9\} \to \{1, 2, \dots, 9\}$ such that, for all distinct $i, j \in \{1, \dots, 9\}$, whenever the squares labelled $i$ and $j$ in the diagram below share an edge, the squares labelled $f(i)$ and $f(j)$ share an edge too?
  \begin{center}
    \begin{tabular}{|c|c|c|}
      \hline
      1 & 2 & 3 \\ \hline
      4 & 5 & 6 \\ \hline
      7 & 8 & 9 \\ \hline
    \end{tabular}
  \end{center}
  \begin{enumerate}[label=(\alph*)]
    \item $8$
    \item $9$
    \item $4$
    \item $24$
  \end{enumerate}

  \item Let $m$ be the number of positive integers $n$ such that $1 \le n \le 2022$ and such that $n$ has an odd number of (positive integer) divisors. Then $m$ is
  \begin{enumerate}[label=(\alph*)]
    \item $22$
    \item $33$
    \item $44$
    \item $55$
  \end{enumerate}

  \item Let $S$ be the set of nonnegative continuous functions $f$ on $[0,1]$ satisfying
  \[
    \int_0^1 \sin^2(x) f(x)\,dx = \int_0^1 \sin(x)\cos(x) f(x)\,dx = \int_0^1 \cos^2(x) f(x)\,dx = 1.
  \]
  Then $S$ is:
  \begin{enumerate}[label=(\alph*)]
    \item an uncountable set
    \item a countably infinite set
    \item a finite and nonempty set
    \item the empty set
  \end{enumerate}

  \item Let $f(x) = 1 - \sin x$ for $x \in \mathbb{R}$. Define
  \[
    a_n = \sqrt[n]{f\left(\frac{1}{n}\right) f\left(\frac{2}{n}\right) \cdots f\left(\frac{n}{n}\right)}.
  \]
  Then
  \begin{enumerate}[label=(\alph*)]
    \item $\{a_n\}_n$ converges to $0$
    \item $\{a_n\}_n$ diverges to $\infty$
    \item $\{a_n\}_n$ converges and $\lim_{n \to \infty} a_n > 0$
    \item none of the other three options is correct
  \end{enumerate}

  \item Which of the following is true for every function $u: \mathbb{R} \to \mathbb{R}$ which is continuously differentiable on $\mathbb{R}$ (i.e., $u$ is differentiable on $\mathbb{R}$ and its derivative $u'$ is continuous on $\mathbb{R}$), and satisfies $u(y) \ge u(x) + u'(x)(y - x)$ for all $x, y \in \mathbb{R}$?
  \begin{enumerate}[label=(\alph*)]
    \item $u'$ is nonnegative.
    \item $u$ attains a minimum at some $x \in \mathbb{R}$.
    \item $u'$ is nondecreasing.
    \item $u'$ is nonincreasing.
  \end{enumerate}

  \item Let $\{x_n\}$ be a sequence of positive numbers such that $\lim_{n \to \infty} x_n = x$. Define
  \[
    z_n = \frac{1}{n}\left[x_1 \left(1 + \frac{x}{n}\right)^n + x_2 \left(1 + \frac{x}{n-1}\right)^{n-1} + \dots + x_n (1 + x)\right].
  \]
  Then
  \begin{enumerate}[label=(\alph*)]
    \item $\{z_n\}$ converges to $xe$
    \item $\{z_n\}$ converges to $e^x$
    \item $\{z_n\}$ does not have a limit
    \item $\{z_n\}$ converges to $x e^x$
  \end{enumerate}

  \item The value of $\displaystyle \lim_{n \to \infty} \int_0^1 \frac{n e^x}{1 + n^2 x^2}\,dx$ is
  \begin{enumerate}[label=(\alph*)]
    \item $\pi e$
    \item $\frac{\pi}{2}$
    \item $\frac{\pi}{2} e$
    \item $\pi$
  \end{enumerate}

  \item What is the number of real solutions of the equation $e^{\sin x} = \pi^x$?
  \begin{enumerate}[label=(\alph*)]
    \item $0$
    \item $1$
    \item Countably infinite
    \item Uncountable
  \end{enumerate}
\end{enumerate}

\subsection*{Part B (True/False)}
\begin{enumerate}
  \item $\mathbb{R}^2 \setminus \mathbb{Q}^2$ is connected but not path-connected.
  \item If $X$ is a connected metric space, and $F$ is a subring of $C(X, \mathbb{R})$ that is a field, then every element of $C(X, \mathbb{R})$ that belongs to $F$ is a constant function.
  \item Let $K \subseteq [0,1]$ be the Cantor set. Then there exists no injective ring homomorphism $C([0,1], \mathbb{R}) \to C(K, \mathbb{R})$.
  \item There exists a metric space $(X, d)$ such that the group of isometries of $X$ is isomorphic to $\mathbb{Z}$.
  \item Let $A \subset \mathbb{R}^2$ be a nonempty subset such that any continuous function $f: A \to \mathbb{R}$ is constant. Then $A$ is a singleton set.
  \item For a nilpotent matrix $A \in M_n(\mathbb{R})$, let
  \[
    \exp(A) := \sum_{k=0}^\infty \frac{A^k}{k!} = \operatorname{Id} + \frac{A}{1!} + \frac{A^2}{2!} + \dots \in M_n(\mathbb{R}).
  \]
  If $A$ is a nilpotent matrix such that $\exp(A) = \operatorname{Id}$, then $A$ is the zero matrix.
  \item There exists $A = \begin{pmatrix} a & b \\ c & d \end{pmatrix} \in M_2(\mathbb{R})$, with $A^2 = A \ne 0$, such that $|a| + |b| < 1$ and $|c| + |d| < 1$.
  \item If $A \in M_3(\mathbb{C})$ is such that $A^i$ has trace zero for all positive integers $i$, then $A$ is nilpotent.
  \item For any finite cyclic group $G$, there exists a prime power $q$ such that $G$ is a subgroup of $\mathbb{F}_q^\times$.
  \item There are only finitely many isomorphism classes of finite nonabelian groups, all of whose proper subgroups are abelian.
  \item Every subring of a unique factorization domain is a unique factorization domain.
  \item Let $f_1, f_2, f_3, f_4 \in \mathbb{R}[x]$ be monic polynomials each of degree exactly two. Then there exist a real polynomial $p \in \mathbb{R}[x]$ and a subset $\{i, j\} \subset \{1, 2, 3, 4\}$ with $i \ne j$, such that $f_i \circ p = c f_j$ for some $c \in \mathbb{R}$.
  \item There exists a finite abelian group $G$ such that the group $\operatorname{Aut}(G)$ of automorphisms of $G$ is isomorphic to $\mathbb{Z}/7\mathbb{Z}$.
  \item There exists an integral domain $R$ and a surjective homomorphism $R \to R$ of rings that is not injective.
  \item There exists $f \in C([0,1], \mathbb{R})$ satisfying the following two conditions:
  \begin{enumerate}[label=(\roman*)]
    \item $\int_0^1 f(x)\,dx = 1$, and
    \item $\lim_{n \to \infty} \int_0^1 f(x)^n\,dx = 0$.
  \end{enumerate}
  \item Let $a_n \ge 0$ for each positive integer $n$. If the series $\sum_{n=1}^\infty \sqrt{a_n}$ converges, then so does the series $\sum_{n=1}^\infty \frac{a_n}{n^{1/4}}$.
  \item There exists a differentiable function $f: \mathbb{R} \to \mathbb{R}$ such that $\lim_{x \to \infty} f(x) = 2$ and $\lim_{x \to \infty} f'(x) = 1$.
  \item Let $f: [0,1] \to [0,\infty)$ be continuous on $[0,1]$ and twice differentiable in $(0,1)$. If $f''(x) = 7f(x)$ for all $x \in (0,1)$, then $f(x) \le \max\{f(0), f(1)\}$ for all $x \in [0,1]$.
  \item There are $N$ balls in a box, out of which $n$ are blue ($1 < n < N$) and the rest are red. Balls are drawn from the box one by one at random, and discarded. Then the probability of picking all the blue balls in the first $n$ draws is the same as the probability of picking all the red balls in the first $(N - n)$ draws.
  \item The set $\{f(x) \in \mathbb{R}[x] \mid f(n) \in \mathbb{Z} \text{ for all } n \in \mathbb{Z}\}$ is uncountable.
\end{enumerate}

\newpage

% ==============================================================================
\section{GS-2023 (Mathematics)}
% ==============================================================================

\subsection*{Notation and Conventions}
\begin{itemize}
  \item $\mathbb{N}$ denotes the set of natural numbers $\{0, 1, \dots\}$, $\mathbb{Z}$ the set of integers, $\mathbb{Q}$ the set of rational numbers, $\mathbb{R}$ the set of real numbers, and $\mathbb{C}$ the set of complex numbers. These sets are assumed to carry the usual algebraic and metric structures.
  \item $\mathbb{R}^n$ denotes the Euclidean space of dimension $n$. Subsets of $\mathbb{R}^n$ are viewed as metric spaces using the standard Euclidean distance on $\mathbb{R}^n$. For $x \in \mathbb{R}^n$, $\|x\|$ denotes the standard Euclidean norm of $x$, i.e., the distance from $x$ to $0$.
  \item All rings are associative, with a multiplicative identity.
  \item For any ring $R$, $M_n(R)$ denotes the ring of $n \times n$ matrices with entries in $R$. The identity matrix in $M_n(R)$ will be denoted by $\operatorname{Id}$.
  \item $M_n(\mathbb{R})$ will also be viewed as a real vector space, and $M_n(\mathbb{C})$ as a complex vector space. $M_n(\mathbb{R})$ is given the topology such that any $\mathbb{R}$-linear isomorphism $M_n(\mathbb{R}) \to \mathbb{R}^{n^2}$ is a homeomorphism. Subsets of $M_n(\mathbb{R})$ are given the subspace topology.
  \item For a ring $R$, $R[x_1, \dots, x_n]$ denotes the polynomial ring in $n$ variables $x_1, \dots, x_n$ over $R$.
  \item All logarithms are natural logarithms.
  \item If $B$ is a subset of a set $A$, we write $A \setminus B$ for the set $\{a \in A \mid a \notin B\}$.
  \item Let $G$ be a finite group, and let $S \subset G$. We say that $S$ generates $G$ if no proper subgroup of $G$ contains $S$.
  \item If $f: X \to Y$ is a map of sets, and $X_1 \subset X$, then $f|_{X_1}$ denotes the restriction of $f$ to $X_1$.
\end{itemize}

\subsection*{Part A (Multiple Choice Questions)}
\begin{enumerate}
  \item Define $f: \mathbb{R} \to \mathbb{R}$ by $f(x) = (3x^2 + 1)/(x^2 + 3)$. Let $f^{\circ 1} = f$ and let $f^{\circ n} = f^{\circ (n-1)} \circ f$ for all integers $n \ge 2$. Which of the following statements is correct?
  \begin{enumerate}[label=(\alph*)]
    \item $\lim_{n \to \infty} f^{\circ n}(1/2) = 1$, and $\lim_{n \to \infty} f^{\circ n}(2) = 1$.
    \item $\lim_{n \to \infty} f^{\circ n}(1/2) = 1$, but $\lim_{n \to \infty} f^{\circ n}(2)$ does not exist.
    \item $\lim_{n \to \infty} f^{\circ n}(1/2)$ does not exist, but $\lim_{n \to \infty} f^{\circ n}(2) = 1$.
    \item Neither $\lim_{n \to \infty} f^{\circ n}(1/2)$ nor $\lim_{n \to \infty} f^{\circ n}(2)$ exists.
  \end{enumerate}

  \item Consider the following properties of a sequence $\{a_n\}_n$ of real numbers.
  \begin{enumerate}[label=(\Roman*)]
    \item $\lim_{n \to \infty} a_n = 0$.
    \item There exists a sequence $\{i_n\}_n$ of positive integers such that $\sum_{n=1}^\infty a_{i_n}$ converges.
  \end{enumerate}
  Which of the following statements is correct?
  \begin{enumerate}[label=(\alph*)]
    \item (I) implies (II), and (II) implies (I).
    \item (I) implies (II), but (II) does not imply (I).
    \item (I) does not imply (II), but (II) implies (I).
    \item (I) does not imply (II), and (II) does not imply (I).
  \end{enumerate}

  \item Consider sequences $\{x_n\}_n$ of real numbers such that $\lim_{n \to \infty} (x_{2n-1} + x_{2n}) = 2$ and $\lim_{n \to \infty} (x_{2n} + x_{2n+1}) = 3$. Which of the following statements is correct?
  \begin{enumerate}[label=(\alph*)]
    \item For every such sequence $\{x_n\}_n$, $\lim_{n \to \infty} \frac{x_{2n+1}}{x_{2n}} = 1$.
    \item For every such sequence $\{x_n\}_n$, $\lim_{n \to \infty} \frac{x_{2n+1}}{x_{2n}} = -1$.
    \item For every such sequence $\{x_n\}_n$, $\lim_{n \to \infty} \frac{x_{2n+1}}{x_{2n}} = 3/2$.
    \item There exists such a sequence $\{x_n\}_n$, for which $\lim_{n \to \infty} \frac{x_{2n+1}}{x_{2n}}$ does not exist.
  \end{enumerate}

  \item Consider the function $f: (0,\infty) \to (0,\infty)$ given by $f(x) = x e^x$. Let $L: (0,\infty) \to (0,\infty)$ be its inverse function. Which of the following statements is correct?
  \begin{enumerate}[label=(\alph*)]
    \item $\lim_{x \to \infty} \frac{L(x)}{\log x} = 1$.
    \item $\lim_{x \to \infty} \frac{L(x)}{(\log x)^2} = 1$.
    \item $\lim_{x \to \infty} \frac{L(x)}{\sqrt{\log x}} = 1$.
    \item None of the remaining three options is correct.
  \end{enumerate}

  \item Let $\{b_n\}_n$ be a monotonically increasing sequence of positive real numbers such that $\lim_{n \to \infty} b_n = \infty$. Which of the following statements is true about $\displaystyle \lim_{n \to \infty} \frac{1}{b_n} \sum_{k=1}^n \frac{b_k}{k^2}$?
  \begin{enumerate}[label=(\alph*)]
    \item The limit exists for all such sequences, and its value is always $+\infty$.
    \item The limit exists for all such sequences, and its value is always $0$.
    \item The limit exists for all such sequences, and its value is always $1$.
    \item None of the remaining three options is correct.
  \end{enumerate}

  \item For every positive integer $n$, define $f_n: [0,1] \to \mathbb{R}$ by
  \[
    f_n(x) = \frac{\sin(n^2 x) + \cos(e^n x)}{1 + n^2 x^2}.
  \]
  Then $\displaystyle \lim_{n \to \infty} \int_0^{1 - \sin(1/n)} f_n(x)\,dx$ equals
  \begin{enumerate}[label=(\alph*)]
    \item $1$.
    \item $0$.
    \item $\infty$.
    \item $1/2$.
  \end{enumerate}

  \item Consider the functions $f_1, f_2: (0,\infty) \to \mathbb{R}$ defined by $f_1(x) = \sqrt{x}$ and $f_2(x) = \sqrt{x} \sin x$. Which of the following statements is correct?
  \begin{enumerate}[label=(\alph*)]
    \item $f_1$ and $f_2$ are uniformly continuous.
    \item $f_1$ is uniformly continuous, but $f_2$ is not.
    \item $f_2$ is uniformly continuous, but $f_1$ is not.
    \item Neither $f_1$ nor $f_2$ is uniformly continuous.
  \end{enumerate}

  \item Let $x_1 \in \mathbb{R}^2 \setminus \{0\}$ be fixed, and inductively define $x_{n+1} = A x_n$ for $n \ge 1$, where $A$ is the $2 \times 2$ real matrix given by
  \[
    A := \begin{pmatrix} \frac{\sqrt{3}}{2} & \frac{1}{2} \\ -\frac{1}{2} & \frac{\sqrt{3}}{2} \end{pmatrix}.
  \]
  Which of the following statements is correct?
  \begin{enumerate}[label=(\alph*)]
    \item $\{x_n\}_n$ is a convergent sequence.
    \item $\{x_n\}_n$ is not a convergent sequence, but it has a convergent subsequence.
    \item $\lim_{n \to \infty} \|x_n\| = 0$.
    \item None of the remaining three options is correct.
  \end{enumerate}

  \item Let $T: M_3(\mathbb{R}) \to \mathbb{R}^3$ be the linear map defined by $T(A) = A \begin{pmatrix} 1 \\ 0 \\ -1 \end{pmatrix}$. Then the dimension of the kernel of $T$ equals
  \begin{enumerate}[label=(\alph*)]
    \item $2$.
    \item $8$.
    \item $1$.
    \item None of the remaining three options.
  \end{enumerate}

  \item Let $V = \{f(x) \in \mathbb{R}[x] \mid f(0) = 0\}$, viewed as a real vector space. Consider the following assertions:
  \begin{enumerate}[label=(\Roman*)]
    \item $V$ contains three linearly independent polynomials of degree $2$.
    \item $V$ contains two linearly independent polynomials of degree $3$.
  \end{enumerate}
  Which of the following statements is correct?
  \begin{enumerate}[label=(\alph*)]
    \item Both (I) and (II) are true.
    \item (I) is true, but (II) is false.
    \item (I) is false, but (II) is true.
    \item Neither (I) nor (II) is true.
  \end{enumerate}

  \item Let $C([-1, 1], \mathbb{R})$ denote the real vector space of continuous functions from $[-1, 1]$ to $\mathbb{R}$, and consider the subspace
  \[
    V = \{f \in C([-1, 1], \mathbb{R}) \mid f(-x) = f(x) \text{ for all } x \in [-1, 1]\}.
  \]
  Define an inner product on $C([-1, 1], \mathbb{R})$ by $\langle f, g \rangle = \int_{-1}^1 f(t) g(t)\,dt$. What is the orthogonal complement of $V$ in $C([-1, 1], \mathbb{R})$?
  \begin{enumerate}[label=(\alph*)]
    \item $\{f \in C([-1, 1], \mathbb{R}) \mid f(-x) = -f(x) \text{ for all } x \in [-1, 1]\}$.
    \item $\{f \in C([-1, 1], \mathbb{R}) \mid f(0) = 0\}$.
    \item $V$ does not have an orthogonal complement in $C([-1, 1], \mathbb{R})$.
    \item None of the remaining three options.
  \end{enumerate}

  \item Consider pairs $(X, d)$, where $X$ is a set with $100$ elements, and $d: X \times X \to \mathbb{R}$ a function such that $d(x, y) = d(y, x) > 0$ if $x, y \in X$ are distinct, and $d(x, x) = 0$ for all $x \in X$. For $n < 100$, let $A_n$ be the statement:
  \begin{quote}
    For every such pair $(X, d)$, there exists a subset $X_1$ of $X$, with $n$ elements, such that $(X_1, d|_{X_1 \times X_1})$ is a metric space.
  \end{quote}
  Which of the following statements is correct?
  \begin{enumerate}[label=(\alph*)]
    \item $A_2$ is true, but $A_3$ is not true.
    \item $A_3$ is true, but $A_4$ is not true.
    \item $A_n$ is true for all $n \le 10$, but not for all $n \le 25$.
    \item $A_n$ is true for all $n \le 25$.
  \end{enumerate}

  \item Let $\{x_n\}_n$ be a sequence in a metric space $(X, d)$. Let $f: X \to \mathbb{R}$ be defined by $f(x) = \inf\{d(x, x_n) \mid n \in \mathbb{N}\}$. Which of the following statements is correct?
  \begin{enumerate}[label=(\alph*)]
    \item $f$ is uniformly continuous on $X$.
    \item $f$ is continuous on $X$, but not necessarily uniformly continuous.
    \item $f$ is continuous on $X$ if and only if $X$ is compact.
    \item None of the remaining three options is correct.
  \end{enumerate}

  \item The number of finite groups, up to isomorphism, with exactly two conjugacy classes, equals
  \begin{enumerate}[label=(\alph*)]
    \item $1$.
    \item $2$.
    \item Greater than $2$, but finite.
    \item Infinite.
  \end{enumerate}

  \item Consider the following assertions about a commutative ring $R$ with identity and elements $a, b \in R$:
  \begin{enumerate}[label=(\Roman*)]
    \item There exist $p, q \in R$ such that $ap + bq = 1$.
    \item There exist $p, q \in R$ such that $a^2 p + b^2 q = 1$.
  \end{enumerate}
  Then:
  \begin{enumerate}[label=(\alph*)]
    \item (I) implies (II), and (II) implies (I).
    \item (I) implies (II), but (II) does not imply (I).
    \item (I) does not imply (II), but (II) implies (I).
    \item (I) does not imply (II), and (II) does not imply (I).
  \end{enumerate}

  \item The number of elements of finite order in the group
  \[
    \left\{ \begin{pmatrix} 1 & a & b \\ 0 & 1 & c \\ 0 & 0 & 1 \end{pmatrix} \;\middle|\; a, b, c \in \mathbb{R} \right\}
  \]
  is
  \begin{enumerate}[label=(\alph*)]
    \item $1$.
    \item Finite, but not $1$.
    \item Countably infinite.
    \item Uncountably infinite.
  \end{enumerate}

  \item The value of
  \[
    \max_{k \ge 1} \{x_1 x_2 \cdots x_k \mid x_1, \dots, x_k \in \mathbb{N}, \text{ and } x_1 + \dots + x_k = 100\}
  \]
  equals
  \begin{enumerate}[label=(\alph*)]
    \item $4 \times 3^{32}$.
    \item $2^{50}$.
    \item $2^{26} \times 3^{16}$.
    \item None of the remaining three options.
  \end{enumerate}

  \item Choose the option that completes the sentence correctly: There exists a $10 \times 10$ real symmetric matrix $A$, all of whose entries are nonnegative and all of whose diagonal entries are positive, such that $A^{10}$ has
  \begin{enumerate}[label=(\alph*)]
    \item exactly $67$ positive entries.
    \item exactly $68$ positive entries.
    \item exactly $69$ positive entries.
    \item exactly $70$ positive entries.
  \end{enumerate}

  \item The number of (nondegenerate Euclidean) triangles with sides of integer length and perimeter $8$, up to congruence, is
  \begin{enumerate}[label=(\alph*)]
    \item $1$.
    \item $2$.
    \item $3$.
    \item $4$.
  \end{enumerate}

  \item Let
  \[
    A = \{(\alpha, \beta) \in \mathbb{Z}^2 \mid \text{the roots } r_1, r_2, r_3 \text{ of the polynomial } p(x) = x^3 - 2x^2 + \alpha x - \beta \text{ satisfy } r_1^3 + r_2^3 + r_3^3 = 0\}.
  \]
  Which of the following statements is correct?
  \begin{enumerate}[label=(\alph*)]
    \item $A$ is infinite.
    \item $A$ is empty.
    \item $A$ is singleton.
    \item $A$ is finite, but neither empty nor singleton.
  \end{enumerate}
\end{enumerate}

\subsection*{Part B (True/False)}
\begin{enumerate}
  \item Let $\alpha$ be a positive real number, and let $f: (0,1) \to \mathbb{R}$ be a function such that $|f(x) - f(y)| \le |x - y|^\alpha$ for all $x, y \in (0,1)$. Then $f$ can be extended to a continuous function $[0,1] \to \mathbb{R}$.
  \item Suppose $f, g: \mathbb{R} \to \mathbb{R}$ are continuous functions such that $f^2 + g^2$ is uniformly continuous. Then at least one of the two functions $f$ and $g$ is uniformly continuous.
  \item Let $\{f_n\}_n$ be a sequence of (not necessarily continuous) functions from $[0,1]$ to $\mathbb{R}$. Let $f: [0,1] \to \mathbb{R}$ be such that for any $x \in [0,1]$ and any sequence $\{x_n\}_n$ consisting of elements from $[0,1]$, if $\lim_{n \to \infty} x_n = x$, then $\lim_{n \to \infty} f_n(x_n) = f(x)$. Then $f$ is continuous.
  \item Let $A, B \in M_2(\mathbb{Z}/2\mathbb{Z})$ be such that $\operatorname{tr}(A) = \operatorname{tr}(B)$ and $\operatorname{tr}(A^2) = \operatorname{tr}(B^2)$. Then $A$ and $B$ have the same eigenvalues.
  \item Let $v_1, v_2, w_1, w_2$ be nonzero vectors in $\mathbb{R}^2$. Then there exists a $2 \times 2$ real matrix $A$ such that $Av_1 = v_2$ and $Aw_1 = w_2$.
  \item Let $A = (a_{ij}) \in M_n(\mathbb{R})$ be such that $a_{ij} \ge 0$ for all $1 \le i, j \le n$. Assume that $\lim_{m \to \infty} A^m$ exists, and denote it by $B = (b_{ij})$. Then, for all $1 \le i, j \le n$, we have $b_{ij} \in \{0, 1\}$.
  \item Given any monic polynomial $f(x) \in \mathbb{R}[x]$ of degree $n$, there exists a matrix $A \in M_n(\mathbb{R})$ such that its characteristic polynomial equals $f$.
  \item If $A \in M_4(\mathbb{Q})$ is such that its characteristic polynomial equals $x^4 + 1$, then $A$ is diagonalizable in $M_4(\mathbb{C})$.
  \item If $A \in M_n(\mathbb{R})$ is such that $AB = BA$ for all invertible matrices $B \in M_n(\mathbb{R})$, then $A = \lambda \cdot \operatorname{Id}$ for some $\lambda \in \mathbb{R}$.
  \item There exists a homeomorphism $f: \mathbb{R} \to \mathbb{R}$ such that $f(2x) = 3f(x)$ for all $x \in \mathbb{R}$.
  \item There exists a continuous bijection from $[0,1] \times [0,1]$ to $\{(x, y) \in \mathbb{R}^2 \mid x^2 + y^2 \le 1\}$, which is not a homeomorphism.
  \item Let $f \in \mathbb{C}[z_1, \dots, z_n]$ be a nonzero polynomial ($n \ge 1$), and let $X = \{z \in \mathbb{C}^n \mid f(z) = 0\}$. Then $\mathbb{C}^n \setminus X$ is path connected.
  \item A connected metric space with at least two points is uncountable.
  \item If $A$ and $B$ are disjoint subsets of a metric space $(X, d)$, then $\inf\{d(x, y) \mid x \in A, \, y \in B\} \ne 0$.
  \item A countably infinite complete metric space has infinitely many isolated points (an element $x$ of a metric space $X$ is said to be an isolated point if $\{x\}$ is an open subset of $X$).
  \item Suppose $G$ and $H$ are two countably infinite abelian groups such that every nontrivial element of $G \times H$ has order $7$. Then $G$ is isomorphic to $H$.
  \item There exists a nonabelian group $G$ of order $26$ such that every proper subgroup of $G$ is abelian.
  \item Let $G$ be a group generated by two elements $x$ and $y$, each of order $2$. Then $G$ is finite.
  \item $\mathbb{R}[x]/(x^4 + x^2 + 2023)$ is an integral domain.
  \item Every finite group is isomorphic to a subgroup of a finite group generated by two elements.
\end{enumerate}

\newpage

% ==============================================================================
\section{GS-2024 (Mathematics)}
% ==============================================================================

\subsection*{Notation and Conventions}
\begin{itemize}
  \item $\mathbb{N}$ denotes the set of natural numbers $\{0, 1, \dots\}$, $\mathbb{Z}$ the set of integers, $\mathbb{Q}$ the set of rational numbers, $\mathbb{R}$ the set of real numbers, and $\mathbb{C}$ the set of complex numbers. These sets are assumed to carry the usual algebraic and metric structures.
  \item $\mathbb{R}^n$ denotes the Euclidean space of dimension $n$. Subsets of $\mathbb{R}^n$ are viewed as metric spaces using the standard Euclidean distance on $\mathbb{R}^n$.
  \item All rings are associative, with a multiplicative identity.
  \item For any ring $R$, $M_n(R)$ denotes the ring of $n \times n$ matrices with entries in $R$. The identity matrix in $M_n(R)$ will be denoted by $\operatorname{Id}$ or by $\operatorname{Id}_n$.
  \item $M_n(\mathbb{R})$ will also be viewed as a real vector space, and $M_n(\mathbb{C})$ as a complex vector space.
  \item For a ring $R$, $R[x_1, \dots, x_n]$ denotes the polynomial ring in $n$ variables $x_1, \dots, x_n$ over $R$.
  \item If $A$ is a set, $\#A$ stands for the cardinality of $A$, and equals $\infty$ if $A$ is infinite.
  \item If $B$ is a subset of a set $A$, we write $A \setminus B$ for the set $\{a \in A \mid a \notin B\}$.
  \item Let $G$ be a finite group, and let $S \subset G$. We say that $S$ generates $G$ if no proper subgroup of $G$ contains $S$.
\end{itemize}

\subsection*{Part A (Multiple Choice Questions)}
\begin{enumerate}
  \item What is the number of even positive integers $n$ such that every group of order $n$ is abelian?
  \begin{enumerate}[label=(\alph*)]
    \item $1$
    \item $2$
    \item Greater than $2$, but finite
    \item Infinite
  \end{enumerate}

  \item Let $n$ be a positive integer, and let $S = \{g \in \mathbb{R}[x] \mid g \text{ is a polynomial of degree at most } n\}$. For $g \in S$ let $A_g = \{x \in \mathbb{R} \mid e^x = g(x)\} \subset \mathbb{R}$. Let $m = \min\{\#A_g \mid g \in S\}$, and $M = \max\{\#A_g \mid g \in S\}$. Then
  \begin{enumerate}[label=(\alph*)]
    \item $m = 0$, $M = n$
    \item $m = 0$, $M = n + 1$
    \item $m = 1$, $M = n$
    \item $m = 1$, $M = n + 1$
  \end{enumerate}

  \item Let $V, W$ be nonzero finite dimensional vector spaces over $\mathbb{C}$. Let $m$ be the dimension of the space of $\mathbb{C}$-linear transformations $V \to W$ viewed as a real vector space. Let $n$ be the dimension of the space of $\mathbb{R}$-linear transformations $V \to W$ viewed as a real vector space. Then
  \begin{enumerate}[label=(\alph*)]
    \item $n = m$
    \item $2n = m$
    \item $n = 2m$
    \item $4n = m$
  \end{enumerate}

  \item Consider the real vector space of infinite sequences of real numbers
  \[
    S = \{(a_0, a_1, a_2, \dots) \mid a_k \in \mathbb{R}, \, k = 0, 1, 2, \dots\}.
  \]
  Let $W$ be the subspace of $S$ consisting of all sequences $(a_0, a_1, a_2, \dots)$ which satisfy the relation $a_{k+2} = 2a_{k+1} + a_k$, $k = 0, 1, 2, \dots$. What is the dimension of $W$?
  \begin{enumerate}[label=(\alph*)]
    \item $1$
    \item $2$
    \item $3$
    \item $\infty$
  \end{enumerate}

  \item Let $f: [0,\infty) \to \mathbb{R}$ be a continuous function. If $\lim_{n \to \infty} \int_0^1 f(x + n)\,dx = 2$, then which of the following statements about the limit $\lim_{n \to \infty} \int_0^1 f(nx)\,dx$ is correct?
  \begin{enumerate}[label=(\alph*)]
    \item The limit exists and equals $0$
    \item The limit exists and equals $\frac{1}{2}$
    \item The limit exists and equals $2$
    \item None of the remaining three options is correct
  \end{enumerate}

  \item Let $f: \mathbb{R} \to [0,\infty)$ be a function such that for any finite set $E \subset \mathbb{R}$ we have $\sum_{x \in E} f(x) \le 1$. Let $C_f = \{x \in \mathbb{R} \mid f(x) > 0\} \subset \mathbb{R}$. Then
  \begin{enumerate}[label=(\alph*)]
    \item $C_f$ is finite
    \item $C_f$ is a bounded subset of $\mathbb{R}$
    \item $C_f$ has at most one limit point
    \item $C_f$ is a countable set
  \end{enumerate}

  \item Let $p$ be a prime. Which of the following statements is true?
  \begin{enumerate}[label=(\alph*)]
    \item There exists a noncommutative ring with exactly $p$ elements
    \item There exists a noncommutative ring with exactly $p^2$ elements
    \item There exists a noncommutative ring with exactly $p^3$ elements
    \item None of the remaining three statements is correct
  \end{enumerate}

  \item Consider the sequence $\{a_n\}$ for $n \ge 1$ defined by $a_n = \lim_{N \to \infty} \sum_{k=n}^N \frac{1}{k^2}$. Which of the following statements about this sequence is true?
  \begin{enumerate}[label=(\alph*)]
    \item $\lim_{n \to \infty} n a_n$ does not exist
    \item $\lim_{n \to \infty} n a_n$ exists and equals $2$
    \item $\lim_{n \to \infty} n a_n$ exists and equals $1$
    \item $\lim_{n \to \infty} n^2 a_n$ exists and equals $1$
  \end{enumerate}

  \item Let $f: \mathbb{R} \to \mathbb{R}$ be a differentiable function that is a solution to the ordinary differential equation
  \[
    f'(t) = \sin^2(f(t)) \quad (\forall t \in \mathbb{R}), \qquad f(0) = 1.
  \]
  Which of the following statements is true?
  \begin{enumerate}[label=(\alph*)]
    \item $f$ is neither bounded nor periodic
    \item $f$ is bounded and periodic
    \item $f$ is bounded, but not periodic
    \item None of the remaining three statements is correct
  \end{enumerate}

  \item Let $B$ denote the set of invertible upper triangular $2 \times 2$ matrices with entries in $\mathbb{C}$, viewed as a group under matrix multiplication. Which of the following subgroups of $B$ is the normalizer of itself in $B$?
  \begin{enumerate}[label=(\alph*)]
    \item $\left\{ \begin{pmatrix} a & 0 \\ 0 & b \end{pmatrix} \;\middle|\; a, b \in \mathbb{C} \setminus \{0\} \right\}$
    \item $\left\{ \begin{pmatrix} a & 0 \\ 0 & a \end{pmatrix} \;\middle|\; a \in \mathbb{C} \setminus \{0\} \right\}$
    \item $\left\{ \begin{pmatrix} 1 & c \\ 0 & 1 \end{pmatrix} \;\middle|\; c \in \mathbb{C} \right\}$
    \item $\left\{ \begin{pmatrix} a & c \\ 0 & a \end{pmatrix} \;\middle|\; a \in \mathbb{C} \setminus \{0\}, \, c \in \mathbb{C} \right\}$
  \end{enumerate}

  \item What is the least positive integer $n > 1$ such that $x^n$ and $x$ are conjugate, for every $x \in S_{11}$? Here, $S_{11}$ denotes the symmetric group on $11$ letters.
  \begin{enumerate}[label=(\alph*)]
    \item $10$
    \item $11$
    \item $12$
    \item $13$
  \end{enumerate}

  \item Consider the following statements:
  \begin{enumerate}[label=(\Alph*)]
    \item Let $G$ be a group and let $H \subset G$ be a subgroup of index $2$. Then $[G, G] \subseteq H$.
    \item Let $G$ be a group and let $H \subset G$ be a subgroup that contains the commutator subgroup $[G, G]$ of $G$. Then $H$ is a normal subgroup of $G$.
  \end{enumerate}
  Which of the following statements is correct?
  \begin{enumerate}[label=(\alph*)]
    \item (A) and (B) are both true
    \item (A) and (B) are both false
    \item (A) is true and (B) is false
    \item (A) is false and (B) is true
  \end{enumerate}

  \item For any symmetric real matrix $A$, let $\lambda(A)$ denote the largest eigenvalue of $A$. Let $S$ be the set of positive definite symmetric $3 \times 3$ real matrices. Which of the following assertions is correct?
  \begin{enumerate}[label=(\alph*)]
    \item There exist $A, B \in S$ such that $\lambda(A + B) < \max(\lambda(A), \lambda(B))$
    \item For all $A, B \in S$, $\lambda(A + B) > \max(\lambda(A), \lambda(B))$
    \item There exist $A, B \in S$ such that $\lambda(A + B) = \max(\lambda(A), \lambda(B))$
    \item None of the remaining three assertions is correct
  \end{enumerate}

  \item Let $\theta \in (0, \pi/2)$. Let $T: \mathbb{R}^2 \to \mathbb{R}^2$ be the linear map which sends a vector $v$ to its reflection with respect to the line through $(0,0)$ and $(\cos\theta, \sin\theta)$. Then the matrix of $T$ with respect to the standard basis of $\mathbb{R}^2$ is given by
  \begin{enumerate}[label=(\alph*)]
    \item $\begin{pmatrix} \cos 2\theta & \sin 2\theta \\ \sin 2\theta & -\cos 2\theta \end{pmatrix}$
    \item $\begin{pmatrix} \cos 2\theta & \sin 2\theta \\ -\sin 2\theta & \cos 2\theta \end{pmatrix}$
    \item $\begin{pmatrix} \cos\theta & \sin\theta \\ \sin\theta & -\cos\theta \end{pmatrix}$
    \item $\begin{pmatrix} \cos\theta & \sin\theta \\ -\sin\theta & \cos\theta \end{pmatrix}$
  \end{enumerate}

  \item For a polynomial $f(x, y) \in \mathbb{R}[x, y]$ let $X_f = \{(a, b) \in \mathbb{R}^2 \mid f(a, b) = 1\} \subset \mathbb{R}^2$. Which of the following statements is correct?
  \begin{enumerate}[label=(\alph*)]
    \item If $f(x, y) = x^2 + 4xy + 3y^2$, then $X_f$ is compact
    \item If $f(x, y) = x^2 - 3xy + 3y^2$, then $X_f$ is compact
    \item If $f(x, y) = x^2 - 4xy - y^2$, then $X_f$ is compact
    \item None of the remaining three statements is correct
  \end{enumerate}

  \item What is the number of distinct subfields of $\mathbb{C}$ isomorphic to $\mathbb{Q}[\sqrt[3]{2}]$?
  \begin{enumerate}[label=(\alph*)]
    \item $1$
    \item $2$
    \item $3$
    \item Infinite
  \end{enumerate}

  \item Let $\mathbb{F}_3$ denote the finite field with $3$ elements. What is the number of one dimensional vector subspaces of the vector space $\mathbb{F}_3^5$ over $\mathbb{F}_3$?
  \begin{enumerate}[label=(\alph*)]
    \item $5$
    \item $121$
    \item $81$
    \item None of the remaining three options
  \end{enumerate}

  \item For a positive integer $n$, let $a_n, b_n, c_n, d_n$ be the real numbers such that
  \[
    \begin{pmatrix} 1 & 1 \\ 1 & 0 \end{pmatrix}^n = \begin{pmatrix} a_n & b_n \\ c_n & d_n \end{pmatrix}.
  \]
  Which of the following numbers equals $\lim_{n \to \infty} a_n / b_n$?
  \begin{enumerate}[label=(\alph*)]
    \item $1$
    \item $e$
    \item $3/2$
    \item None of the remaining three options
  \end{enumerate}

  \item Consider the complex vector space
  \[
    V = \{f \in \mathbb{C}[x] \mid f \text{ has degree at most } 50, \text{ and } f(ix) = -f(x) \text{ for all } x \in \mathbb{C}\}.
  \]
  Then the dimension of $V$ equals
  \begin{enumerate}[label=(\alph*)]
    \item $50$
    \item $25$
    \item $13$
    \item $47$
  \end{enumerate}

  \item Let $S$ denote the set of sequences $a = (a_1, a_2, \dots)$ of real numbers such that $a_k$ equals $0$ or $1$ for each $k$. Then the function $f: S \to \mathbb{R}$ defined by
  \[
    f((a_1, a_2, \dots)) = \frac{a_1}{10} + \frac{a_2}{10^2} + \dots
  \]
  is
  \begin{enumerate}[label=(\alph*)]
    \item injective but not surjective
    \item surjective but not injective
    \item bijective
    \item neither injective nor surjective
  \end{enumerate}
\end{enumerate}

\subsection*{Part B (True/False)}
\begin{enumerate}
  \item If $G$ is a group of order $361$, then $G$ has a normal subgroup $H$ such that $H \cong G/H$.
  \item There exists a metric space $X$ such that the number of open subsets of $X$ is exactly $2024$.
  \item The function $d: \mathbb{R} \times \mathbb{R} \to \mathbb{R}$ given by $d(x, y) = |e^x - e^y|$ defines a metric on $\mathbb{R}$, and $(\mathbb{R}, d)$ is a complete metric space.
  \item Let $n$ be a positive integer, and $A$ an $n \times n$ matrix over $\mathbb{R}$ such that $A^3 = \operatorname{Id}$. Then $A$ is diagonalizable in $M_n(\mathbb{R})$, i.e., there exists $P \in M_n(\mathbb{R})$ such that $P$ is invertible and $PAP^{-1}$ is a diagonal matrix.
  \item If $A \in M_n(\mathbb{Q})$ is such that the characteristic polynomial of $A$ is irreducible over $\mathbb{Q}$, then $A$ is diagonalizable in $M_n(\mathbb{C})$, i.e., there exists $P \in M_n(\mathbb{C})$ such that $P$ is invertible and $PAP^{-1}$ is a diagonal matrix.
  \item The complement of any countable union of lines in $\mathbb{R}^3$ is path connected.
  \item The subsets $\{(x, y) \in \mathbb{R}^2 \mid (y^2 - x)(y^2 - x - 1) = 0\}$ and $\{(x, y) \in \mathbb{R}^2 \mid y^2 - x^2 = 1\}$ of $\mathbb{R}^2$ (with the induced metric) are homeomorphic.
  \item $\mathbb{Q} \cap [0,1]$ is a compact subset of $\mathbb{Q}$.
  \item Suppose $f: X \to Y$ is a function between metric spaces, such that whenever a sequence $\{x_n\}$ converges to $x$ in $X$, the sequence $\{f(x_n)\}$ converges in $Y$ (but it is not given that the limit of $\{f(x_n)\}$ is $f(x)$). Then $f$ is continuous.
  \item Let $f: \mathbb{R} \to \mathbb{R}$ be differentiable, and assume that $|f'(x)| \ge 1$ for all $x \in \mathbb{R}$. Then for each compact set $C \subset \mathbb{R}$, the set $f^{-1}(C)$ is compact.
  \item There exists a function $f: [0,1] \to \mathbb{R}$ which is not Riemann integrable and satisfies
  \[
    \sum_{i=1}^n |f(t_i) - f(t_{i-1})|^2 < 1
  \]
  for every choice of a positive integer $n$ and of $0 \le t_0 < t_1 < t_2 < \dots < t_n \le 1$.
  \item Let $E \subset [0,1]$ be the subset consisting of numbers that have a decimal expansion which does not contain the digit $8$. Then $E$ is dense in $[0,1]$.
  \item Let $G$ be a proper subgroup of $(\mathbb{R}, +)$ which is closed as a subset of $\mathbb{R}$. Then $G$ is generated by a single element.
  \item There exists a unique function $f: \mathbb{R} \to \mathbb{R}$ such that $f$ is continuous at $x = 0$, and such that for all $x \in \mathbb{R}$, $f(x) + f\left(\frac{x}{2}\right) = x$.
  \item A map $f: V \to W$ between finite dimensional vector spaces over $\mathbb{R}$ is a linear transformation if and only if $f(x) = f(x - a) + f(x - b) - f(x - a - b)$ for all $x, a, b \in V$.
  \item Let $R$ be the ring $\mathbb{C}[x]/(x^2)$ obtained as the quotient of the polynomial ring $\mathbb{C}[x]$ by its ideal generated by $x^2$. Let $R^\times$ be the multiplicative group of units of this ring. Then there is an injective group homomorphism from $(\mathbb{Z}/2\mathbb{Z}) \times (\mathbb{Z}/2\mathbb{Z})$ into $R^\times$.
  \item Let $A \in M_2(\mathbb{Z})$ be such that $|A_{ij}(n)| \le 50$ for all $1 \le n \le 10^{50}$ and all $1 \le i, j \le 2$, where $A_{ij}(n)$ denotes the $(i, j)$-th entry of the $2 \times 2$ matrix $A^n$. Then $|A_{ij}(n)| \le 50$ for all positive integers $n$.
  \item Let $A, B$ be subsets of $\{0, \dots, 9\}$. It is given that, on choosing elements $a \in A$ and $b \in B$ at random, $a + b$ takes each of the values $0, \dots, 9$ with equal probability. Then one of $A$ or $B$ is singleton.
  \item If $f: \mathbb{R} \to \mathbb{R}$ is uniformly continuous, then there exists $M > 0$ such that for all $x \in \mathbb{R} \setminus [-M, M]$, we have $f(x) < x^{100}$.
  \item If a sequence $\{f_n\}$ of continuous functions from $[0,1]$ to $\mathbb{R}$ converges uniformly on $(0,1)$ to a continuous function $f: [0,1] \to \mathbb{R}$, then $\{f_n\}$ converges uniformly on $[0,1]$ to $f$.
\end{enumerate}

\newpage

% ==============================================================================
\section{GS-2025 (Mathematics PhD at CAM)}
% ==============================================================================

\subsection*{Notations and Conventions}
\begin{itemize}
  \item $\mathbb{N}$ denotes the set of natural numbers $\{0, 1, \dots\}$, $\mathbb{Z}$ the set of integers, $\mathbb{Q}$ the set of rational numbers, $\mathbb{R}$ the set of real numbers, and $\mathbb{C}$ the set of complex numbers. These sets are assumed to carry the usual algebraic and metric structures.
  \item For sets $A$ and $B$, $A \setminus B = \{x \in A : x \notin B\}$.
  \item $i = \sqrt{-1}$.
  \item $\mathbb{R}^n$, $n \ge 1$, is equipped with the Euclidean topology.
  \item For a subset $A$ of a topological space, $\overline{A}$ denotes the closure of $A$.
  \item For a real normed linear space $X$, the dual of $X$ is $X^* = \{\varphi: X \to \mathbb{R} : \varphi \text{ is linear and continuous}\}$, and $\langle Tx, y \rangle := (Tx)(y)$ for $T: X \to X^*$ and $x, y \in X$.
  \item $C([0,1])$ denotes the set of continuous real-valued functions on the interval $[0,1]$.
  \item For $k \in \mathbb{N}$, $C^k([0,1])$ denotes the set of real-valued functions having $k$ continuous derivatives on $(0,1)$, and $f^{(k)} \in C([0,1])$.
  \item For a set $S \subseteq \mathbb{R}$ of positive Lebesgue measure and $1 \le p < \infty$, $L^p(S) = \{f: S \to \mathbb{R} : f \text{ is measurable and } \int_S |f(x)|^p\,dx < \infty\}$.
  \item For an $n \times n$ matrix $A$, $\det(A)$ denotes the determinant of $A$.
\end{itemize}

\subsection*{Questions}
\begin{enumerate}
  \item For $0 < a < 1$, let $\lambda = \frac{1}{a} + ai$ and
  \[
    z_n = \int_{\mathbb{R}} e^{inx} e^{-\lambda x^2}\,dx.
  \]
  What is $\lim_{n \to \infty} z_n$?
  \begin{enumerate}[label=(\alph*)]
    \item $a$
    \item $\infty$
    \item $0$
    \item $1$
  \end{enumerate}

  \item Let $\mathbb{D} = \{(x, y) \in \mathbb{R}^2 : x^2 + y^2 < 1\}$. For $p \in \mathbb{R}$, define
  \[
    C_p = \int_{\mathbb{D}} \frac{\sin^2(x^2 + y^2)}{(x^2 + y^2)^p}\,dx\,dy.
  \]
  Choose the correct option.
  \begin{enumerate}[label=(\alph*)]
    \item $C_p < \infty$ for all $p > 0$.
    \item $C_p < \infty$ for all $p > 1$.
    \item $C_p < \infty$ for all $p < 3$.
    \item $C_p < \infty$ for all $p < 4$.
  \end{enumerate}

  \item Let $p(x) = x^3 - x^2 - x + 1$. Consider the following statements.
  \begin{enumerate}[label=(\Roman*)]
    \item If $T: \mathbb{R}^3 \to \mathbb{R}^3$ is a linear transformation and $p(T) = 0$, then $T$ is invertible.
    \item If $T: \mathbb{R}^3 \to \mathbb{R}^3$ is a linear transformation and $p(T) = 0$, then $T$ is diagonalizable.
  \end{enumerate}
  Choose the correct option.
  \begin{enumerate}[label=(\alph*)]
    \item (I) and (II) are both true.
    \item (I) and (II) are both false.
    \item (I) is true and (II) is false.
    \item (I) is false and (II) is true.
  \end{enumerate}

  \item What is the cardinality of $\{A \subseteq \mathbb{R} : \mathbb{R} \setminus \overline{A} = \overline{\mathbb{R} \setminus A}\}$?
  \begin{enumerate}[label=(\alph*)]
    \item $0$
    \item $1$
    \item $2$
    \item infinite
  \end{enumerate}

  \item Let $y \in C^2([0,1])$ be a non identically zero solution to
  \[
    \frac{d^2 y}{dt^2}(t) + (e^{t^2} \sin t) y(t) = 0, \qquad t \in (0,1).
  \]
  Choose the correct option.
  \begin{enumerate}[label=(\alph*)]
    \item $y$ can have infinitely many zeroes.
    \item Either $y$ has no zeroes or finitely many zeroes.
    \item The third derivative of $y$ does not exist at $1/2$.
    \item The equation has a unique non identically zero solution.
  \end{enumerate}

  \item Let $f$ be an entire function which satisfies $|f(z)| \ge |z|^2$ for all $z \in \mathbb{C}$ and assume moreover that $f(i) = -i$. Then,
  \begin{enumerate}[label=(\alph*)]
    \item $f(0) \ne 0$
    \item $f(1) = -3$
    \item $f(2)$ cannot be uniquely determined from the given information.
    \item $f(3) = 9i$
  \end{enumerate}

  \item Let $\mathbb{H} = \{x + iy : x \in \mathbb{R}, \, y > 0\}$ and $f: \mathbb{H} \to \mathbb{H}$ be a bijective holomorphic function such that $\lim_{|z| \to \infty} |f(z)| = \infty$. Choose the correct option.
  \begin{enumerate}[label=(\alph*)]
    \item $f$ is a linear function.
    \item $f$ is a rational function that is not a polynomial.
    \item $f(\{it : t > 0\})$ is a subset of a circle of finite radius.
    \item $f(\{it : t > 0\})$ is a subset of a horizontal line.
  \end{enumerate}

  \item Let $f: \mathbb{R} \to \mathbb{R}$ be defined by $f(y) = y \sin(1/y)$ for $y \ne 0$, and $f(0) = 0$. Then, the initial value problem
  \[
    \begin{cases}
      \frac{dy}{dt}(t) = f(y(t)), & t \in \mathbb{R}, \\
      y(0) = 0,
    \end{cases}
  \]
  has
  \begin{enumerate}[label=(\alph*)]
    \item at least one strictly monotonic solution.
    \item a unique solution; which is $y \equiv 0$.
    \item infinitely many solutions.
    \item at least one oscillatory solution (i.e., the solution is not identically zero and has infinitely many zeroes).
  \end{enumerate}

  \item Choose the correct option. For any real Banach space $X$ and any linear operator $T: X \to X^*$, if $\langle Tx, y \rangle = \langle Ty, x \rangle$ for all $x, y \in X$, then
  \begin{enumerate}[label=(\alph*)]
    \item $T$ is invertible.
    \item $T$ is surjective.
    \item $T$ is bounded.
    \item $T = 0$.
  \end{enumerate}

  \item Choose the correct option.
  \begin{enumerate}[label=(\alph*)]
    \item If $f: \mathbb{R} \to \mathbb{R}$ is a bijective differentiable map then its inverse is also a bijective differentiable map.
    \item If $f: \mathbb{R} \to \mathbb{R}$ is a bijective differentiable map then for all $x \in \mathbb{R}$ there exists an $\epsilon > 0$ and a differentiable map $g: (x - \epsilon, x + \epsilon) \to \mathbb{R}$ such that $f \circ g(t) = t$ for all $t \in (x - \epsilon, x + \epsilon)$.
    \item If $f: \mathbb{R} \to \mathbb{R}$ is a bijective differentiable map then $f$ is either monotonically increasing or monotonically decreasing.
    \item If $f: \mathbb{R} \to \mathbb{R}$ is a bijective differentiable map then $f'(x) > 0$ for all $x \in \mathbb{R}$ or $f'(x) < 0$ for all $x \in \mathbb{R}$.
  \end{enumerate}

  \item Let $f: \mathbb{R} \to \mathbb{R}$ be a continuous function satisfying the following condition: for each $x_0$ and for each $r > 0$,
  \[
    f(x_0) = \frac{1}{r} \int_{x_0 - r}^{x_0 + r} f(y)\,dy.
  \]
  Choose the correct statement.
  \begin{enumerate}[label=(\alph*)]
    \item There exists a non constant function $f \in L^2(\mathbb{R})$ satisfying the above property.
    \item There exists at most one function in $L^2(\mathbb{R})$ that satisfies the above property.
    \item The space of functions $f \in L^2(\mathbb{R})$ and satisfying the above property is one dimensional.
    \item The space of functions $f \in L^2(\mathbb{R})$ and satisfying the above property is infinite dimensional.
  \end{enumerate}

  \item Consider the $11 \times 11$ matrix:
  \[
    A = \begin{pmatrix}
      -2 & \frac{1}{1^2} & -\frac{1}{2^2} & \frac{1}{3^2} & \cdots & -\frac{1}{10^2} \\[4pt]
      \frac{1}{11^2} & -2 & -\frac{1}{12^2} & \frac{1}{13^2} & \cdots & -\frac{1}{20^2} \\[4pt]
      \frac{1}{21^2} & -\frac{1}{22^2} & -2 & \frac{1}{23^2} & \cdots & -\frac{1}{30^2} \\[4pt]
      \vdots & \vdots & \vdots & \vdots & \ddots & \vdots \\[4pt]
      \frac{1}{101^2} & \cdots & \cdots & \cdots & -\frac{1}{110^2} & -2
    \end{pmatrix}.
  \]
  Consider the following statements.
  \begin{enumerate}[label=(\Roman*)]
    \item The eigenvalues of $A$ are all non-negative.
    \item $A$ is invertible.
  \end{enumerate}
  Choose the correct option.
  \begin{enumerate}[label=(\alph*)]
    \item (I) and (II) are both true.
    \item (I) and (II) are both false.
    \item (I) is true and (II) is false.
    \item (I) is false and (II) is true.
  \end{enumerate}

  \item For sets $A, B \subseteq \mathbb{R}^n$, define $A + B = \{a + b : a \in A \text{ and } b \in B\}$. Choose the correct option.
  \begin{enumerate}[label=(\alph*)]
    \item If $A$ is compact and $B$ is closed, then $A + B$ is compact.
    \item If $A$ is closed and $B$ is closed, then $A + B$ is closed.
    \item If $A$ is open and $B$ is compact, then $A + B$ is open.
    \item If $A$ is open and $B$ is compact, then $A + B$ is compact.
  \end{enumerate}

  \item Choose the correct option. If $g: \mathbb{R} \to \mathbb{R}$ is a positive continuous function such that $\int_0^\infty e^{\alpha x} g(x)\,dx < \infty$ for some $\alpha > 0$, then
  \begin{enumerate}[label=(\alph*)]
    \item $\int_0^\infty \frac{g(x)}{x^n}\,dx < \infty$ for all $n \in \mathbb{N}$.
    \item $\lim_{x \to \infty} g(x)$ exists.
    \item $\int_0^\infty \frac{1}{g^\beta(x)}\,dx < \infty$ for some $\beta > 0$.
    \item there exist $C > 0$ and $\beta > 0$ such that $\int_t^\infty g(x)\,dx \le C e^{-\beta t}$, for all $t > 0$.
  \end{enumerate}

  \item Let $f_n: \mathbb{R} \to \mathbb{R}$ be a sequence of differentiable functions such that
  \[
    \int_0^1 f_n(x)\,dx = 0 \quad \text{and} \quad |f_n'(x)| \le x^{-1/2} \quad \text{for all } x \in (0,1].
  \]
  Consider the following statements.
  \begin{enumerate}[label=(\Roman*)]
    \item $\{f_n\}$ is uniformly bounded on $[0,1]$.
    \item $\{f_n\}$ has a subsequence that converges uniformly on $[0,1]$.
  \end{enumerate}
  Choose the correct option.
  \begin{enumerate}[label=(\alph*)]
    \item (I) and (II) are both true.
    \item (I) and (II) are both false.
    \item (I) is true and (II) is false.
    \item (I) is false and (II) is true.
  \end{enumerate}

  \item Let $n \ge 2$, and $A$ be an $n \times n$ real symmetric and positive definite matrix satisfying $\operatorname{trace}(A) = \det(A)$. Choose the correct statement that holds for every such $A$.
  \begin{enumerate}[label=(\alph*)]
    \item $\operatorname{trace}(A) < n^{\frac{n}{n-1}}$.
    \item $\operatorname{trace}(A) \le n$.
    \item $\det(A) \ge n^{\frac{n}{n-1}}$.
    \item all of the eigenvalues of $A$ are $\le \frac{1}{2}$.
  \end{enumerate}

  \item A box contains $2024$ balls in which $200$ are blue and the rest are red. Balls are chosen from the box one by one at random, and discarded. Consider the following statements.
  \begin{enumerate}[label=(\Roman*)]
    \item The probability of picking a red ball at the $k$-th pick is the same for all $k = 1, \dots, 2024$.
    \item The probability of picking a blue ball at the $200$-th pick is the same as the probability of picking a red ball at the $1824$-th pick.
    \item The probability of picking all the blue balls in the first $200$ draws is the same as the probability of picking all the red balls in the first $1824$ draws.
  \end{enumerate}
  Choose the correct option.
  \begin{enumerate}[label=(\alph*)]
    \item (I) and (II) are true, (III) is false.
    \item (I) and (III) are true, (II) is false.
    \item (II) and (III) are true, (I) is false.
    \item (I), (II), and (III) are all true.
  \end{enumerate}

  \item Let $X$ denote the inner product space $(C([0,1]), \langle \cdot, \cdot \rangle)$ where $\langle f, g \rangle = \int_0^1 f(x) g(x)\,dx$. Define $T: X \to X$ by $(Tf)(x) = \int_x^1 f(t)\,dt$ for $f \in X$ and $x \in [0,1]$. Choose the correct statement.
  \begin{enumerate}[label=(\alph*)]
    \item $T$ is surjective.
    \item There exists a bounded linear operator $U: X \to X$ such that $U \circ T$ is the identity map.
    \item $T(X)$ is dense in $X$.
    \item $T(X)$ is a closed subset of $X$.
  \end{enumerate}

  \item Let $T: L^2(0,1) \to L^2(0,1)$ be defined by $(Tf)(x) = f(\sqrt{x})$. We denote the operator norm of $T$ by $\|T\|$. Choose the correct option.
  \begin{enumerate}[label=(\alph*)]
    \item $\|T\| = \sqrt{2}$.
    \item $\|T\| = 2$.
    \item $\|T\| = \frac{1}{\sqrt{2}}$.
    \item $\|T\| = \frac{1}{2}$.
  \end{enumerate}

  \item Let $f: \mathbb{R} \to \mathbb{R}$ be defined by
  \[
    f(x) = \begin{cases}
      x^{-1/2}, & \text{if } 0 < x < 1, \text{ and} \\
      0, & \text{otherwise.}
    \end{cases}
  \]
  Let $r_n$ be an enumeration of rational numbers and correspondingly define
  \[
    h(x) = \sum_{n=1}^\infty 2^{-n} f(x - r_n), \qquad \text{for } x \in \mathbb{R}.
  \]
  Choose the correct statement.
  \begin{enumerate}[label=(\alph*)]
    \item $h \in L^1(\mathbb{R})$ and $h \in L^2(\mathbb{R})$.
    \item $h \in L^1(\mathbb{R})$ and $h \notin L^2(\mathbb{R})$.
    \item $h \notin L^1(\mathbb{R})$ and $h \in L^2(\mathbb{R})$.
    \item $h \notin L^1(\mathbb{R})$ and $h \notin L^2(\mathbb{R})$.
  \end{enumerate}
\end{enumerate}

\newpage

% ==============================================================================
\section{GS-2025 (Mathematics)}
% ==============================================================================

\subsection*{Notation and Conventions}
\begin{itemize}
  \item $\mathbb{N}$ denotes the set of natural numbers $\{0, 1, \dots\}$, $\mathbb{Z}$ the set of integers, $\mathbb{Q}$ the set of rational numbers, $\mathbb{R}$ the set of real numbers, and $\mathbb{C}$ the set of complex numbers. These sets are assumed to carry the usual algebraic and metric structures.
  \item $\mathbb{R}^n$ denotes the Euclidean space of dimension $n$. Subsets of $\mathbb{R}^n$ are viewed as metric spaces using the standard Euclidean distance on $\mathbb{R}^n$. For $x \in \mathbb{R}^n$, $\|x\|$ denotes the standard Euclidean norm of $x$, i.e., the distance from $x$ to $0$.
  \item All rings are associative, with a multiplicative identity. Ring homomorphisms are assumed to respect multiplicative identities.
  \item For any ring $R$, $M_n(R)$ denotes the ring of $n \times n$ matrices with entries in $R$. The identity matrix in $M_n(R)$ will be denoted by $\operatorname{Id}$ or by $\operatorname{Id}_n$. For a matrix $A \in M_n(R)$, $A^t$ will stand for its transpose, $\operatorname{trace}(A)$ for its trace, and $\det(A)$ for its determinant.
  \item $M_n(\mathbb{R})$ will also be viewed as a real vector space, and $M_n(\mathbb{C})$ as a complex vector space. $M_n(\mathbb{R})$ is given the topology such that any $\mathbb{R}$-linear isomorphism $M_n(\mathbb{R}) \to \mathbb{R}^{n^2}$ is a homeomorphism. Subsets of $M_n(\mathbb{R})$ are given the subspace topology.
  \item For a ring $R$, $R[x_1, \dots, x_n]$ denotes the polynomial ring in $n$ variables $x_1, \dots, x_n$ over $R$.
  \item If $A$ is a set, $\#A$ stands for the cardinality of $A$, and equals $\infty$ if $A$ is infinite.
  \item If $B$ is a subset of a set $A$, we write $A \setminus B$ for the set $\{a \in A \mid a \notin B\}$.
  \item Let $G$ be a group, and let $S \subset G$. The subgroup of $G$ generated by $S$ is defined to be the smallest subgroup of $G$ that contains $S$.
\end{itemize}

\subsection*{Part A (Multiple Choice Questions)}
\begin{enumerate}
  \item For a positive integer $m$, let $d(m)$ denote the number of divisors of $m$, including $1$ and $m$. Define a sequence $\{a_n\}_{n=1}^\infty$ by
  \[
    a_n = \#\{1 \le m \le n \mid d(m) \text{ is odd}\}.
  \]
  Find $\lim_{n \to \infty} \frac{a_n}{n}$.
  \begin{enumerate}[label=(\alph*)]
    \item $0$
    \item $1$
    \item $1/2$
    \item None of the remaining three options
  \end{enumerate}

  \item Let $S_3$ be the symmetric group on $3$ letters. Consider the real vector space
  \[
    V = \{f: S_3 \to \mathbb{R} \mid f(g) = f(g^3), \, \forall g \in S_3\}.
  \]
  What is the dimension of this space?
  \begin{enumerate}[label=(\alph*)]
    \item $1$
    \item $2$
    \item $3$
    \item Greater than $3$
  \end{enumerate}

  \item Consider the set $S = \{(x, y) \in \mathbb{R}^2 \mid (x^3 + 1)^3 = (y^5 + 1)^5\}$. How many connected components does $S$ have?
  \begin{enumerate}[label=(\alph*)]
    \item $1$
    \item $3$
    \item Greater than $3$, but finite
    \item Infinite
  \end{enumerate}

  \item Define a set $S$ of real polynomials as follows:
  \[
    S = \{p \in \mathbb{R}[x] \mid p \text{ is not constant, } |p(0)| < 1\}.
  \]
  For $p$ in $S$, define $U_p = \{x \in \mathbb{R} \mid |p(x)| < 1\}$. Which of the following statements is correct?
  \begin{enumerate}[label=(\alph*)]
    \item There exists $p \in S$ such that $U_p$ is a union of infinitely many disjoint open intervals.
    \item For each $p \in S$, $U_p$ is a union of infinitely many disjoint open intervals.
    \item For each $p \in S$, $U_p$ is a union of finitely many open intervals.
    \item There exists $p \in S$ such that $U_p$ is an unbounded set.
  \end{enumerate}

  \item What is the number of functions $f: \mathbb{R} \setminus \{0, 1\} \to \mathbb{R} \setminus \{0, 1\}$ such that $|f(x) - f(y)| = |x - y|$ for all $x, y \in \mathbb{R} \setminus \{0, 1\}$?
  \begin{enumerate}[label=(\alph*)]
    \item $1$
    \item $2$
    \item Greater than $2$, but finite
    \item Infinite
  \end{enumerate}

  \item Let $x_n = \left(1 - \frac{1}{\sqrt{n}}\right)^n \exp(n^{1/4})$. Which of the following statements is correct?
  \begin{enumerate}[label=(\alph*)]
    \item $\lim_{n \to \infty} x_n = 0$ and $\sum_{n=1}^\infty x_n$ converges.
    \item $\lim_{n \to \infty} x_n = 0$ but $\sum_{n=1}^\infty x_n$ diverges.
    \item $\lim_{n \to \infty} x_n$ exists but is non-zero.
    \item $\lim_{n \to \infty} x_n$ does not exist.
  \end{enumerate}

  \item Let $f: [0,1] \to \mathbb{R}$ be such that
  \[
    f(x) = \begin{cases}
      \frac{1}{2^n}, & x = \frac{p}{2^n} \text{ where } p \text{ is an odd integer, and} \\
      0, & \text{otherwise.}
    \end{cases}
  \]
  Which of the following statements is correct?
  \begin{enumerate}[label=(\alph*)]
    \item $f$ is continuous at $x \in [0,1]$ if and only if $x$ is rational.
    \item $f$ is continuous at $x \in [0,1]$ if and only if $x$ is irrational.
    \item $f$ is not continuous at $x$, for any $x \in [0,1]$.
    \item None of the remaining three statements is true.
  \end{enumerate}

  \item Let $\{q_n\}_{n=1}^\infty$ be an enumeration of the rationals. In other words, $\mathbb{Q} = \{q_n \mid n \in \mathbb{N} \setminus \{0\}\}$ and $q_m \ne q_n$ if $m \ne n$. Set
  \[
    X = \bigcup_{n=1}^\infty \left(q_n - \frac{1}{n^2}, \, q_n + \frac{1}{n^2}\right).
  \]
  Which of the following statements is correct?
  \begin{enumerate}[label=(\alph*)]
    \item $X$ is dense in $\mathbb{R}$, but not equal to $\mathbb{R}$.
    \item $X = \mathbb{R}$.
    \item $X$ contains an open interval of length $10$.
    \item None of the remaining three statements is correct.
  \end{enumerate}

  \item Consider the linear map $T: \mathbb{R}[x] \to \mathbb{R}[x]$ defined by $T(p(x)) = \frac{d}{dx}(x p(x))$. Which of the following statements is correct?
  \begin{enumerate}[label=(\alph*)]
    \item $T$ is injective but not surjective.
    \item $T$ is surjective but not injective.
    \item $T$ is bijective.
    \item $T$ is neither injective nor surjective.
  \end{enumerate}

  \item Let
  \[
    A = \begin{pmatrix} 1 & 2 & 1 \\ 2 & 1 & 2 \\ 1 & 2 & 1 \end{pmatrix}.
  \]
  Consider the following statements:
  \begin{enumerate}[label=(\roman*)]
    \item $A$ has a positive eigenvalue.
    \item $A$ has a negative eigenvalue.
  \end{enumerate}
  Which of the following statements is correct?
  \begin{enumerate}[label=(\alph*)]
    \item (i) and (ii) are true.
    \item (i) and (ii) are false.
    \item (i) is true and (ii) is false.
    \item (ii) is true and (i) is false.
  \end{enumerate}

  \item Let $A$ be a nonzero $n \times n$ matrix with real entries, where $n > 1$. Consider the following statements.
  \begin{enumerate}[label=(\roman*)]
    \item If $A^2 = 0$ then $\operatorname{rank}(A) \le \lfloor \frac{n}{2} \rfloor$.
    \item If $\operatorname{rank}(A) \le \lfloor \frac{n}{2} \rfloor$ then $A^2 = 0$.
  \end{enumerate}
  Which of the following statements is correct?
  \begin{enumerate}[label=(\alph*)]
    \item (i) and (ii) are true.
    \item (i) and (ii) are false.
    \item (i) is true and (ii) is false.
    \item (ii) is true and (i) is false.
  \end{enumerate}

  \item Consider the following subset of $\mathbb{R}^2$:
  \[
    S = \left\{\left(x, \sin\frac{1}{x}\right) \;\middle|\; x > 0\right\} \cup \left\{\left(0, \pm\frac{1}{n}\right) \;\middle|\; n \ge 1\right\}.
  \]
  Which of the following statements is correct?
  \begin{enumerate}[label=(\alph*)]
    \item $S$ is path connected.
    \item $S$ is connected, but not path connected.
    \item $S$ is not connected, but has finitely many connected components.
    \item $S$ has infinitely many connected components.
  \end{enumerate}

  \item Consider the following statements:
  \begin{enumerate}[label=(\roman*)]
    \item Any linear map $\mathbb{R}^3 \to \mathbb{R}^3$ has a one-dimensional invariant subspace.
    \item Any linear map $\mathbb{R}^3 \to \mathbb{R}^3$ has a two-dimensional invariant subspace.
  \end{enumerate}
  Which of the following statements is correct?
  \begin{enumerate}[label=(\alph*)]
    \item (i) and (ii) are both true.
    \item (i) and (ii) are both false.
    \item (i) is true and (ii) is false.
    \item (i) is false and (ii) is true.
  \end{enumerate}

  \item Consider the following statements:
  \begin{enumerate}[label=(\roman*)]
    \item For any linear map $T: \mathbb{R}^3 \to \mathbb{R}^3$ with a two-dimensional kernel, the trace of $T$ is an eigenvalue of $T$.
    \item For any linear map $T: \mathbb{R}^3 \to \mathbb{R}^3$ with a two-dimensional kernel, $X^2$ divides the characteristic polynomial of $T$.
  \end{enumerate}
  Which of the following statements is correct?
  \begin{enumerate}[label=(\alph*)]
    \item (i) and (ii) are both true.
    \item (i) and (ii) are both false.
    \item (i) is true and (ii) is false.
    \item (i) is false and (ii) is true.
  \end{enumerate}

  \item For a positive integer $n$, define the function $f_n: \mathbb{R} \to \mathbb{R}$ by $f_n(x) = \sin\left(\frac{x}{n}\right)$. Consider the sequences $\{f_n\}_{n=1}^\infty$ and $\{f_n'\}_{n=1}^\infty$ of functions from $\mathbb{R}$ to $\mathbb{R}$. Which of the following statements is correct?
  \begin{enumerate}[label=(\alph*)]
    \item The sequences $\{f_n\}_{n=1}^\infty$ and $\{f_n'\}_{n=1}^\infty$ both converge uniformly on $\mathbb{R}$.
    \item Neither of the sequences $\{f_n\}_{n=1}^\infty$ and $\{f_n'\}_{n=1}^\infty$ converges uniformly on $\mathbb{R}$.
    \item The sequence $\{f_n\}_{n=1}^\infty$ converges uniformly on $\mathbb{R}$, but the sequence $\{f_n'\}_{n=1}^\infty$ does not.
    \item The sequence $\{f_n'\}_{n=1}^\infty$ converges uniformly on $\mathbb{R}$, but the sequence $\{f_n\}_{n=1}^\infty$ does not.
  \end{enumerate}

  \item Let $n > 1$ be a positive integer. Let $A = [a_{ij}]_{1 \le i,j \le n}$ be an $n \times n$ matrix, such that $a_{i, i+1} = c_i$ for $i = 1, \dots, (n - 1)$ and $a_{n, 1} = c_n$ where $c_1, \dots, c_n$ are real numbers and all the other entries are $0$. Which of the following sentences is correct regardless of $n$?
  \begin{enumerate}[label=(\alph*)]
    \item $\det(\operatorname{Id} - A^n) = 0$ if $\prod_{i=1}^n c_i = 1$.
    \item $\det(\operatorname{Id} - A^n) > 0$ if $\prod_{i=1}^n c_i > 0$.
    \item $\det(\operatorname{Id} - A^n) < 0$ if $\prod_{i=1}^n c_i > 1$.
    \item $\det(\operatorname{Id} - A^n) > 0$ if $\prod_{i=1}^n c_i > 1$.
  \end{enumerate}

  \item Let $n \ge 1$. Consider non-zero $n \times n$ real matrices $A$ such that $\operatorname{trace}(AX) = 0$ for every real matrix $X$ with $\operatorname{trace}(X) = 0$. Consider the following assertions.
  \begin{enumerate}[label=(\roman*)]
    \item For every such matrix $A$, $\operatorname{trace}(A)^n = n^n \cdot \det(A)$.
    \item For every such matrix $A$, the rank of $A$ is $n$.
  \end{enumerate}
  Which of the following sentences is correct?
  \begin{enumerate}[label=(\alph*)]
    \item (i) is correct and (ii) is incorrect.
    \item (ii) is correct and (i) is incorrect.
    \item (i) and (ii) are correct.
    \item (i) and (ii) are incorrect.
  \end{enumerate}

  \item Let $S$ be the set of functions $f: (0,1) \to \mathbb{R}$ with the property that there is a sequence $\{f_n\}_{n=1}^\infty$ of functions that converges uniformly to $f$, where each $f_n: (0,1) \to \mathbb{R}$ is a twice continuously differentiable function (i.e., $f_n''$ exists and is continuous). Which of the following statements is correct?
  \begin{enumerate}[label=(\alph*)]
    \item Any $f \in S$ is continuous, but there exists $f \in S$ such that $f$ is not differentiable.
    \item Any $f \in S$ is differentiable, but there exists $f \in S$ such that $f'$ is not continuous.
    \item Any $f \in S$ is continuously differentiable, but there exists $f \in S$ such that $f''$ does not exist.
    \item Any $f \in S$ is twice continuously differentiable.
  \end{enumerate}

  \item For a real valued function $f: [0,1] \to \mathbb{R}$, let $\omega_f: [0,1] \to \mathbb{R}$ be the oscillation of $f$ defined by
  \[
    \omega_f(t) := \sup_{|x - y| \le t} |f(x) - f(y)|.
  \]
  Which one of the following statements is correct for every $f: [0,1] \to \mathbb{R}$?
  \begin{enumerate}[label=(\alph*)]
    \item $f$ is continuous if and only if $\lim_{t \to 0+} \omega_f(t) = 0$.
    \item If $f$ is continuous then $\lim_{t \to 0+} \omega_f(t) = 0$, but not conversely.
    \item If $\lim_{t \to 0+} \omega_f(t) = 0$ then $f$ is continuous, but not conversely.
    \item None of the remaining three statements is correct.
  \end{enumerate}

  \item The number of real cubic polynomials of the form $x^3 + ax + b$, each of whose complex roots lies on the unit circle $S^1 = \{z \in \mathbb{C} \mid z \overline{z} = 1\}$, equals
  \begin{enumerate}[label=(\alph*)]
    \item $0$
    \item $2$
    \item $\infty$
    \item None of the remaining three options
  \end{enumerate}
\end{enumerate}

\subsection*{Part B (True/False)}
\begin{enumerate}
  \item There exists a linear polynomial $p(x) \in \mathbb{R}[x]$ for which there exist exactly two subsets $S \subset \mathbb{R}$ such that $p(S) = S$ and such that $\#S = 2$.
  \item There are exactly $10$ abelian groups of order $2025$, up to isomorphism.
  \item If $A \subset \mathbb{R}$ is an uncountable subset, then there exist uncountably many elements $x \in A$ such that $x$ is a limit point of $A \setminus \{x\}$.
  \item There exists an open subset of $\mathbb{R}$ with uncountably many connected components.
  \item Let $(X, d)$ be a nonempty complete metric space, and let $T: X \to X$ be a continuous map such that $T^2$ is a contraction, i.e., there exists $a < 1$ such that for all distinct $x, y \in X$, we have $d(T^2(x), T^2(y)) < a d(x, y)$. Then there exists a unique element $x \in X$ such that $T(x) = x$.
  \item Let $S^2 = \{(x, y, z) \in \mathbb{R}^3 \mid x^2 + y^2 + z^2 = 1\}$, and let $d: S^2 \times S^2 \to \mathbb{R}$ be the restriction of the Euclidean metric on $\mathbb{R}^3$. If $f: S^2 \to S^2$ is a map such that $d(f(x), f(y)) = d(x, y)$ for all $x, y \in S^2$, then $f$ is surjective.
  \item Let $C([0,1], \mathbb{R})$ be the set of continuous functions from $[0,1]$ to $\mathbb{R}$, equipped with the metric $d$ given by $d(f_1, f_2) = \sup_{x \in [0,1]} |f_1(x) - f_2(x)|$. Let $X \subset C([0,1], \mathbb{R})$ be defined by
  \[
    X := \left\{f \in C([0,1], \mathbb{R}) \;\middle|\; \int_0^1 f(t)\,dt \ne 0\right\}.
  \]
  Then, with the induced metric, $X$ is connected.
  \item Let $SO_4 = \{A \in M_4(\mathbb{R}) \mid AA^t = \operatorname{Id} \text{ and } \det(A) > 0\}$. Then every matrix in $SO_4$ has $1$ or $-1$ as an eigenvalue.
  \item Let $T$ be an endomorphism of a finite dimensional real vector space $V$, and let $W$ denote the image of $T$. Let $T'$ denote the restriction of $T$ to $W$. Then $T$ and $T'$ have the same trace.
  \item Every subgroup of $S^1 = \{z \in \mathbb{C} \mid z \overline{z} = 1\}$, viewed as a group under multiplication, is finite.
  \item If $f: F \to F$ is a homomorphism of fields, then $f$ is surjective.
  \item Given any linear transformation $T: \mathbb{R}^2 \to \mathbb{R}^2$, there exist linear transformations $T_1, T_2: \mathbb{R}^2 \to \mathbb{R}^2$ such that $(1, 0)$ belongs to the kernel of $T_1$, $(0, 1)$ belongs to the image of $T_2$, and $T = T_1 + T_2$.
  \item For any positive integer $d$, the set $\{10^a - 10^b \mid a, b \text{ are distinct positive integers}\}$ contains a multiple of $d$.
  \item There exists $\theta \in (0, \pi/2)$ such that $\sin\theta + \sec\theta + \cot\theta = 3$.
  \item Let $m$ be a positive integer, and $S_m$ the symmetric group on $m$ letters. Let $n$ be a positive integer such that for every group $G$ of order $n$, there exists an injective group homomorphism $G \hookrightarrow S_m$. Then $m \ge n$.
  \item Let $l_1$ be the line in $\mathbb{R}^2$ joining $(0,0)$ and $\left(\frac{1}{2}, \frac{\sqrt{3}}{2}\right)$ and $l_2$ the line in $\mathbb{R}^2$ joining $(0,0)$ and $\left(\frac{\sqrt{3}}{2}, \frac{1}{2}\right)$. Consider the group of bijections $\mathbb{R}^2 \to \mathbb{R}^2$ under composition, and its subgroup $G$ generated by the reflections about $l_1$ and $l_2$. Then $G$ has exactly $12$ elements.
  \item There are only finitely many subsets $I \subset \mathbb{N}$ with the property that
  \[
    \left\{f(x) = \sum_{i=0}^n a_i x^i \in \mathbb{R}[x] \;\middle|\; n \in \mathbb{N}, \, a_i \in \mathbb{Q} \text{ for all } i \in I\right\}
  \]
  is a subring of $\mathbb{R}[x]$.
  \item There are uncountably many continuous functions $\mathbb{Q} \to \{0, 1\}$.
  \item Let $f: \mathbb{R} \to \mathbb{R}$ be a continuously differentiable function. Suppose there exists $C \in \mathbb{R}$ such that $f$ is injective on $(C, \infty)$. Then there exists $D \in \mathbb{R}$ such that $f'$ is injective on $(D, \infty)$.
  \item If $f, g: [0,1] \to \mathbb{R}$ are continuous and $f \cdot g = 0$, then one of $f$ or $g$ is zero on an open subset of $[0,1]$.
\end{enumerate}


\newpage
\section{GS-2026 (Mathematics PhD at CAM)}
\subsection*{Notations and Conventions}
\begin{itemize}
\item All sets considered are assumed to be nonempty.
\item $\mathbb{N}=\{1,2,\ldots\}$, $\mathbb{Z}$, $\mathbb{Q}$, $\mathbb{R}$, $\mathbb{C}$ carry the usual algebraic and metric structures.
\item $A\setminus B=\{x\in A: x\notin B\}$. \ $\Re z$, $\Im z$ are the real and imaginary parts of $z\in\mathbb{C}$. \ $\mathbb{D}=\{z\in\mathbb{C}:|z|<1\}$.
\item $\mathbb{R}^n$ ($n\ge1$) has the Euclidean topology. \ $\overline A$ denotes the closure of $A$.
\item $C([0,1])$ is the set of continuous real-valued functions on $[0,1]$. $C^1([0,1])$ is the set of real-valued functions on $[0,1]$ that are differentiable on $[0,1]$ (one-sided derivatives at $0$ and $1$) with $f'\in C([0,1])$.
\item For measurable $S\subseteq\mathbb{R}$ and $1\le p\le\infty$, $L^p(S)$ is the standard Lebesgue space with norm $\|\cdot\|_p$.
\item For a subset $S$ of a Hilbert space $H$ with inner product $\langle\cdot,\cdot\rangle$, $S^\perp=\{h\in H:\langle h,s\rangle=0 \text{ for all } s\in S\}$.
\end{itemize}

\subsection*{Part A}
\begin{enumerate}
\item Consider the following subsets of $\mathbb{R}^2$.
\begin{itemize}
\item $X=\{(t,\sin(1/t)):t\in(0,1)\}$
\item $Y=\{(t,\cos(1/t)):t\in(0,1)\}$
\item $Z=\{(t,\exp(1/t)):t\in(0,1)\}$
\end{itemize}
Choose the correct option.
\begin{enumerate}[label=(\alph*)]
\item Only one of the sets $X,Y,Z$ is connected.
\item Exactly two of the sets $X,Y,Z$ are connected.
\item All of the three sets $X,Y,Z$ are connected.
\item None of the sets $X,Y,Z$ are connected.
\end{enumerate}

\item \textbf{Theorem:} Let $(X,d)$ be a metric space. Suppose $E$ and $F$ are disjoint subsets of $X$, and that $E$ is unbounded. If \underline{\hspace{3cm}}, then there exists a $\delta>0$ such that $d(e,f)\ge\delta$ for all $e\in E$ and $f\in F$.

Choose the statement below that when used in the blank above makes the theorem true.
\begin{enumerate}[label=(\alph*)]
\item $E$ is open and $F$ is open.
\item $E$ is open and $F$ is compact.
\item $E$ is closed and $F$ is closed.
\item $E$ is closed and $F$ is compact.
\end{enumerate}

\item Let $X=(0,1)$ and $f_n,f:X\to\mathbb{R}$ be continuous functions. Consider the following statements.
\begin{enumerate}[label=(\roman*)]
\item $f_n$ converges uniformly to $f$ on $X$.
\item For any convergent sequence $(x_n)$ in $X$, that is $\lim x_n\in X$, we have $\lim f_n(x_n)=f(\lim x_n)$.
\end{enumerate}
Choose the correct option.
\begin{enumerate}[label=(\alph*)]
\item (i) $\Rightarrow$ (ii) and (ii) $\Rightarrow$ (i).
\item (i) $\Rightarrow$ (ii) but (ii) $\not\Rightarrow$ (i).
\item (i) $\not\Rightarrow$ (ii) but (ii) $\Rightarrow$ (i).
\item (i) $\not\Rightarrow$ (ii) and (ii) $\not\Rightarrow$ (i).
\end{enumerate}

\item For the two dimensional Lebesgue measure on the set $X=\{(t^2,t^3)\in\mathbb{R}^2:t\in\mathbb{R}\}$, consider the following statements.
\begin{enumerate}[label=(\roman*)]
\item If $Y\subseteq X$, then $Y$ is measurable.
\item $\mathbb{R}^2\setminus X$ is measurable.
\item $X$ has infinite measure.
\end{enumerate}
Choose the correct option.
\begin{enumerate}[label=(\alph*)]
\item (i), (ii), and (iii) are true.
\item (i) and (ii) are true, but (iii) is false.
\item (ii) and (iii) are true, but (i) is false.
\item (i), (ii), and (iii) are false.
\end{enumerate}

\item Let $f:\mathbb{R}\to\mathbb{R}$ be continuous such that $f(0)=1$ and $f(x)=0$ for $|x|>1$. Define $f_n(x)=f(x+n)$ for $n\in\mathbb{N}$ and $x\in\mathbb{R}$. Consider the following statements.
\begin{enumerate}[label=(\roman*)]
\item $(f_n)$ converges in $L^\infty(\mathbb{R})$.
\item $(f_n)$ converges in $L^p(\mathbb{R})$ for some $1\le p<\infty$.
\item $(f_n)$ converges in $L^p(\mathbb{R})$ for all $1\le p<\infty$.
\end{enumerate}
Choose the correct option.
\begin{enumerate}[label=(\alph*)]
\item (i), (ii), and (iii) are true.
\item (i) and (ii) are true, but (iii) is false.
\item (i) and (iii) are true, but (ii) is false.
\item (i), (ii), and (iii) are false.
\end{enumerate}

\item Let $M=\left\{f\in L^2(-1,1):\int_{-1}^1 f(x)\,\mathrm{d} x=0\right\}$, and $g(x)=x^2$. What is $\inf\{\|g-f\|_2: f\in M\}$?
\begin{enumerate}[label=(\alph*)]
\item $\dfrac{\sqrt2}{3}$. \item $\dfrac13$. \item $\dfrac23$. \item $1$.
\end{enumerate}

\item Let $T:\bigl(C^1([0,1]),\|\cdot\|_\infty\bigr)\to\bigl(C([0,1]),\|\cdot\|_\infty\bigr)$ be defined by $T(f)=f'$. Consider the following statements.
\begin{enumerate}[label=(\roman*)]
\item $T$ is a closed operator.
\item $\bigl(C^1([0,1]),\|\cdot\|_\infty\bigr)$ is a Banach space.
\end{enumerate}
Choose the correct option.
\begin{enumerate}[label=(\alph*)]
\item (i) is true but (ii) is false.
\item (i) and (ii) are both true.
\item (i) and (ii) are both false.
\item (i) is false but (ii) is true.
\end{enumerate}

\item Consider the holomorphic functions
\[ f(z)=\frac{\sin z}{z^3}\quad\text{and}\quad g(w)=f\!\left(\frac1w\right),\qquad z,w\in\mathbb{C}\setminus\{0\}. \]
Choose the correct statement.
\begin{enumerate}[label=(\alph*)]
\item $f(z)$ has a pole of order $3$ at $z=0$ and $g(w)$ has a pole of order $3$ at $w=0$.
\item $f$ has a pole of order $3$ at $z=0$ and $g(w)$ has an essential singularity at $w=0$.
\item $f$ has a pole of order $2$ at $z=0$ and $g(w)$ has a pole of order $3$ at $w=0$.
\item $f$ has a pole of order $2$ at $z=0$ and $g(w)$ has an essential singularity at $w=0$.
\end{enumerate}

\item Consider the following domains in $\mathbb{C}$.
\begin{enumerate}[label=(\roman*)]
\item $\{z\in\mathbb{C}:\Re z>0\text{ and }\Im z>0\}$.
\item A bounded domain.
\item $\mathbb{C}\setminus\{0\}$.
\item $\mathbb{C}\setminus\overline{U}$ where $U$ is an open subset of $\mathbb{C}$ with $U\neq\mathbb{C}$.
\end{enumerate}
From which of the above domains does there exist an injective holomorphic function into the unit disc $\mathbb{D}$?
\begin{enumerate}[label=(\alph*)]
\item (i), (ii), and (iii).
\item (i), (ii), and (iv).
\item (i), (iii), and (iv).
\item (i) and (ii).
\end{enumerate}

\item Let $X$ be a Hausdorff topological space such that for all closed sets $A\subseteq X$ and continuous functions $f:A\to[0,1]$ there exists a continuous function $\tilde f:X\to[0,1]$ such that $\tilde f|_A=f$. Which of the following statements is necessarily true?
\begin{enumerate}[label=(\alph*)]
\item If $C_1$ and $C_2$ are disjoint closed sets, then there are disjoint open sets $O_1$ and $O_2$ such that $C_1\subseteq O_1$ and $C_2\subseteq O_2$.
\item $X$ is compact.
\item $X$ is connected.
\item All singleton subsets of $X$ are open.
\end{enumerate}
\end{enumerate}

\subsection*{Part B}
\begin{enumerate}[resume]
\item Let $X$ and $Y$ be metric spaces, and $f:X\to Y$ be a bijective continuous function. Choose the correct statement.
\begin{enumerate}[label=(\alph*)]
\item If $X$ is compact, then $f$ is a homeomorphism.
\item If $Y$ is compact, then $X$ is compact.
\item If $Y$ is connected, then $X$ is connected.
\item If $X=\mathbb{R}$ and $Y\subseteq\mathbb{R}^2$, then $Y$ is a closed set.
\end{enumerate}

\item Let $E$ be a closed convex subset of a Hilbert space $H$. For each $x\in H$, there exists a unique $y_0\in E$ minimizing $\{\|x-y\|:y\in E\}$. Which of the following does this $y_0$ satisfy?
\begin{enumerate}[label=(\alph*)]
\item $\langle x-y_0,z-y_0\rangle=0,\ \forall z\in E$.
\item $\langle x-y_0,z-y_0\rangle\le0,\ \forall z\in E$.
\item $\langle x-y_0,z-y_0\rangle\ge0,\ \forall z\in E$.
\item $y_0\in x+E^\perp$.
\end{enumerate}

\item Consider the following functional equation for $f$ in $C([0,1])$:
\[ 2f(x)-\int_0^1\cos(f(y))\sqrt{1+(x-y)^2}\,\mathrm{d} y+\sin(e^x)=0,\qquad x\in[0,1]. \]
Choose the correct statement.
\begin{enumerate}[label=(\alph*)]
\item There exists no solution in $C([0,1])$.
\item There exists a unique solution $f\in C([0,1])$.
\item There exist exactly two solutions in $C([0,1])$.
\item There exist infinitely many solutions in $C([0,1])$.
\end{enumerate}

\item Consider the following functions on $\mathbb{R}\setminus\{0\}$.
\begin{enumerate}[label=(\roman*)]
\item $f(t)=t+2t^2\sin(1/t)$.
\item $g(t)=t+2t^3\sin(1/t)$.
\item $h(t)=t^3\cos(t)$.
\end{enumerate}
Which among the above functions has an inverse in $(-\epsilon,\epsilon)\setminus\{0\}$ for some $\epsilon>0$?
\begin{enumerate}[label=(\alph*)]
\item Only (iii).
\item Only (ii) and (iii).
\item Only (ii).
\item All the functions (i), (ii), and (iii).
\end{enumerate}

\item Let $f_n$ and $f$ be real-valued Lebesgue measurable functions on $[0,1]$. Suppose that $(f_n)$ converges to $f$ pointwise. Consider the following statements.
\begin{enumerate}[label=(\roman*)]
\item For all $\epsilon>0$ there exists a compact set $K\subseteq[0,1]$ such that $(f_n)$ converges uniformly to $f$ on $K$ and the Lebesgue measure of $K$ is greater than $1-\epsilon$.
\item There exists a dense set $U\subseteq[0,1]$ such that $(f_n)$ converges uniformly to $f$ on $U$.
\end{enumerate}
Which of the above statements are necessarily true?
\begin{enumerate}[label=(\alph*)]
\item (i) and (ii) are both true.
\item (i) is true but (ii) is false.
\item (ii) is true but (i) is false.
\item (i) and (ii) are both false.
\end{enumerate}

\item Choose the correct option.
\begin{enumerate}[label=(\alph*)]
\item There exists a function $f\in C(\mathbb{R})$ such that $\displaystyle\lim_{\epsilon\to0}\int_0^1\bigl(f(x+\epsilon)-f(x)\bigr)\,\mathrm{d} x\ne0$.
\item There exists a function $f\in L^\infty(\mathbb{R})$ such that $\displaystyle\lim_{\epsilon\to0}\int_0^1\bigl(f(x+\epsilon)-f(x)\bigr)\,\mathrm{d} x\ne0$.
\item There exists a function $f\in L^2(\mathbb{R})$ such that $\displaystyle\lim_{\epsilon\to0}\int_0^1\bigl(f(x+\epsilon)-f(x)\bigr)\,\mathrm{d} x\ne0$.
\item For all functions $f\in L^1(\mathbb{R})$, $\displaystyle\lim_{\epsilon\to0}\int_0^1\bigl(f(x+\epsilon)-f(x)\bigr)\,\mathrm{d} x=0$.
\end{enumerate}

\item Let $f_n(x)=\sin(n\pi x)$. Choose the correct statement.
\begin{enumerate}[label=(\alph*)]
\item $(f_n)$ converges in $L^2(0,2)$.
\item $(f_n)$ does not have any convergent subsequence in $L^2(0,2)$.
\item There exists $E\subseteq(0,2)$ such that $E$ has Lebesgue measure $2$ and $f_n|_E$ converges pointwise.
\item There exists an interval $F\subset(0,2)$ such that $f_n|_F$ converges in $L^2(F)$.
\end{enumerate}

\item Let $X$ denote $\mathbb{R}^{2025}$ equipped with the norm $\|(x_1,\ldots,x_{2025})\|=\sum_{i=1}^{2025}|x_i|$ and $T:X\to X$ be defined by
\[ T(u)=\begin{cases}u,&\text{if }\|u\|\le1,\\ u/\|u\|,&\text{if }\|u\|>1.\end{cases} \]
Consider the following statements.
\begin{enumerate}[label=(\roman*)]
\item $\|T(u)-T(v)\|\le\|u-v\|$ for all $u,v\in X$.
\item $\|T(u)-T(v)\|\le2\|u-v\|$ for all $u,v\in X$.
\item $\|T(u)-T(v)\|\le3\|u-v\|$ for all $u,v\in X$.
\end{enumerate}
Choose the correct option.
\begin{enumerate}[label=(\alph*)]
\item (i), (ii) and (iii) are true.
\item Only (ii) and (iii) are true.
\item Only (iii) is true.
\item (i), (ii) and (iii) are false.
\end{enumerate}

\item Let $\mathcal F=\{f:\mathbb{D}\to\mathbb{D}:\ f \text{ is holomorphic and } f(\tfrac12)=0\}$ and $\alpha=\sup\{\Im f(0):f\in\mathcal F\}$. Choose the correct option.
\begin{enumerate}[label=(\alph*)]
\item $\alpha=1$.
\item $\alpha=\dfrac13$.
\item $\alpha=\dfrac12$.
\item $\alpha=\dfrac23$.
\end{enumerate}

\item Let $f:\mathbb{D}\to\mathbb{C}$ be holomorphic such that $|f(z)|\ge\sqrt{|z|}$ for all $z\in\mathbb{D}$. Consider the following statements.
\begin{enumerate}[label=(\roman*)]
\item There exists $z_0\in\mathbb{D}$ such that $|f(z_0)|<1$.
\item There exists $z_0\in\mathbb{D}$ such that $|f(z_0)|\ge1$.
\item The function $f$ is bounded.
\end{enumerate}
Choose the correct option.
\begin{enumerate}[label=(\alph*)]
\item Only (i) is true.
\item Only (ii) is true.
\item Only (i) and (ii) are true.
\item Only (ii) and (iii) are true.
\end{enumerate}
\end{enumerate}

\newpage
%=====================================================================
\section{GS-2026 (Mathematics: Stage I Questions)}
\subsection*{Notation and Conventions}
\begin{itemize}
\item $\mathbb{N}=\{0,1,\ldots\}$, $\mathbb{Z}$, $\mathbb{Q}$, $\mathbb{R}$, $\mathbb{C}$ carry the usual algebraic and metric structures. For $x>0$, $\log(x)$ is the natural logarithm.
\item For a metric space $(X,d)$, $x\in X$, $r>0$: $B(x,r)=\{y\in X\mid d(x,y)<r\}$.
\item $\mathbb{R}^n$ is Euclidean space; $\|x\|$ is the Euclidean norm.
\item All rings are associative with identity; ring homomorphisms preserve identity.
\item $M_n(R)$ is the ring of $n\times n$ matrices over $R$; $\mathrm{Id}$ or $\mathrm{Id}_n$ is the identity; $A^{\mathsf J}$ is the transpose, $\operatorname{trace}(A)$ the trace, $\det(A)$ the determinant. $GL_n(R)\subset M_n(R)$ is the group of invertible matrices. $M_n(\mathbb{R})$, $M_n(\mathbb{C})$ carry the Euclidean topology.
\item $R[x_1,\ldots,x_n]$ is the polynomial ring over $R$.
\item For a compact metric space $X$, $C(X,\mathbb{R})$ is the real vector space of continuous real-valued functions with metric $d(f,g)=\sup_{x\in X}|f(x)-g(x)|$.
\item $\#A$ is the cardinality of $A$ ($=\infty$ if infinite). \ $A\setminus B=\{a\in A\mid a\notin B\}$.
\item For a group $G$ and $S\subseteq G$, the subgroup generated by $S$ is the smallest subgroup of $G$ containing $S$.
\end{itemize}

\subsection*{Part A --- MCQ}
\begin{enumerate}
\item For $\alpha>0$, and for all positive integers $n$, let $a_n=\log(1+n^{-\alpha})$. Then,
\begin{enumerate}[label=(\alph*)]
\item $\sum_{n=1}^\infty a_n$ is convergent for all $\alpha>0$.
\item $\sum_{n=1}^\infty a_n$ is convergent for all $\alpha>1$.
\item $\sum_{n=1}^\infty a_n$ is convergent for all $0<\alpha\le1$.
\item $\sum_{n=1}^\infty a_n$ is divergent for all $\alpha>1$.
\end{enumerate}

\item What is the value of $\displaystyle\lim_{n\to\infty}\left(\prod_{k=1}^n\frac{k}{k+n}\right)^{1/n}$?
\begin{enumerate}[label=(\alph*)]
\item $1/e$ \item $1/2$ \item $1/\log(4)$ \item $1/4$
\end{enumerate}

\item For $x\in[0,1]$, let $f(x)=\displaystyle\sum_{n=1}^\infty\frac{x^n}{n^2}$. Then
\begin{enumerate}[label=(\alph*)]
\item $f$ is continuous but nowhere differentiable.
\item $f$ is differentiable on $(0,1)$ and its derivative extends continuously to $[0,1]$.
\item $f$ is differentiable on $(0,1)$ but the derivative diverges as $x\to1^-$.
\item $f$ can be extended to a real analytic function on $(0,1+\varepsilon)$ for some $\varepsilon>0$.
\end{enumerate}

\item For each positive integer $n$ and each $x\in[0,1]$, define
\[ f_n(x)=\sin(n)\exp\!\Bigl(-n^4\bigl(x-\tfrac1n\bigr)^2\Bigr). \]
Which of the following is true about the sequence of functions $(f_n)$ on $[0,1]$?
\begin{enumerate}[label=(\alph*)]
\item $(f_n)$ converges to $0$ uniformly on $[0,1]$.
\item $(f_n)$ converges to $0$ pointwise, but not uniformly, on $[0,1]$.
\item $(f_n)$ does not converge pointwise on $[0,1]$.
\item $(f_n)$ converges uniformly to a nonzero function on $[0,1]$.
\end{enumerate}

\item Let $f:[0,\infty)\to\mathbb{R}$ be a continuous function. Consider the statements:
\begin{enumerate}[label=(\roman*)]
\item If $\lim_{x\to\infty}f(x)=\infty$, then $f$ is not uniformly continuous.
\item If $f$ is bounded, then $f$ is uniformly continuous.
\end{enumerate}
Which of the following is correct?
\begin{enumerate}[label=(\alph*)]
\item Both (i) and (ii) are true.
\item (i) is true, (ii) is false.
\item (i) is false, (ii) is true.
\item Both (i) and (ii) are false.
\end{enumerate}

\item Consider the set $\mathcal A=\{A\in M_{10}(\mathbb{R})\mid A^{10}=0\}$. For each $A\in M_{10}(\mathbb{R})$, define the linear transformation $T_A:M_{10}(\mathbb{R})\to M_{10}(\mathbb{R})$ by $T_A(X)=AX$. Then which of the following is true?
\begin{enumerate}[label=(\alph*)]
\item For all $A\in\mathcal A$, we have $\det(T_A)=0$ and $\operatorname{trace}(T_A)=0$.
\item There exists $A\in\mathcal A$ for which $\det(T_A)\neq0$ whereas $\operatorname{trace}(T_A)=0$.
\item There exists $A\in\mathcal A$ for which $\det(T_A)=0$ whereas $\operatorname{trace}(T_A)\neq0$.
\item There exists $A\in\mathcal A$ for which $\det(T_A)\neq0$ and $\operatorname{trace}(T_A)\neq0$.
\end{enumerate}

\item Let $(X,d)$ be a complete metric space and $U$ be a proper dense subset of $X$. Let $f:U\to\mathbb{R}$ be a function. Consider the following two statements:
\begin{enumerate}[label=(\roman*)]
\item If $f$ is continuous, then it can be extended to a continuous function $g:X\to\mathbb{R}$.
\item If $f$ is uniformly continuous, then it can be uniquely extended to a continuous function $g:X\to\mathbb{R}$.
\end{enumerate}
Then which of the following is correct?
\begin{enumerate}[label=(\alph*)]
\item Both (i) and (ii) are true.
\item Both (i) and (ii) are false.
\item (i) is true but (ii) is false.
\item (ii) is true but (i) is false.
\end{enumerate}

\item Let $(X,d)$ be a metric space with the additional property that $d(x,y)\le\max\{d(x,z),d(y,z)\}$ for all $x,y,z\in X$. Consider the following statements:
\begin{enumerate}[label=(\roman*)]
\item For all $x,y,z\in X$, at least two of the numbers $d(x,y),d(y,z),d(x,z)$ are equal.
\item For each positive real number $r$, and $x,y\in X$, the sets $B(x,r)$ and $B(y,r)$ are either equal or disjoint.
\end{enumerate}
Which of the following is true?
\begin{enumerate}[label=(\alph*)]
\item (i) and (ii) are true.
\item (i) and (ii) are false.
\item (i) is true and (ii) is false.
\item (ii) is true and (i) is false.
\end{enumerate}

\item Let $T:\mathbb{R}^n\to\mathbb{R}^n$ be a linear map such that $T\neq0$, and $\mathbb{R}^n=\mathrm{range}(T)\oplus\mathrm{null\ space}(T)$. Consider the following statements:
\begin{enumerate}[label=(\roman*)]
\item $T$ is diagonalisable.
\item There does not exist any integer $k>0$ such that $T^k=0$.
\end{enumerate}
Which of the following must be true?
\begin{enumerate}[label=(\alph*)]
\item (i) is false and (ii) is true.
\item (i) is true and (ii) is false.
\item (i) and (ii) are true.
\item (i) and (ii) are false.
\end{enumerate}

\item For any differentiable function $g:\mathbb{R}\to\mathbb{R}$, let
\[ A_g=\{x\in[0,1]\mid g(x)=0\},\qquad B_g=\{x\in[0,1]\mid g'(x)=0\}. \]
Let $S=\{g:\mathbb{R}\to\mathbb{R}\mid g\text{ is differentiable and }A_g\cap B_g=\emptyset\}$. Which of the following statements is true?
\begin{enumerate}[label=(\alph*)]
\item For all $g\in S$, the set $B_g$ is finite.
\item For all $g\in S$, the set $A_g$ is finite.
\item There exists $g\in S$ such that both $A_g$ and $B_g$ are infinite.
\item None of the other statements is true.
\end{enumerate}

\item Let $g:[1,\infty)\to\mathbb{R}$ be a function such that
\[ g(1)=1\quad\text{and}\quad g'(x)=\frac{1}{x^2+g^2(x)}\ \text{ for all } x\ge1. \]
Then which of the following statements is true?
\begin{enumerate}[label=(\alph*)]
\item $\lim_{x\to\infty}g(x)$ exists, and belongs to $(2,\infty)$.
\item $\lim_{x\to\infty}g(x)$ exists, and belongs to $(-\infty,2]$.
\item None of the other statements is true.
\item $\lim_{x\to\infty}g(x)=\infty$.
\end{enumerate}

\item Let $A,B$ be linear endomorphisms of a finite dimensional vector space over a field. Consider the following statements:
\begin{enumerate}[label=(\roman*)]
\item Every eigenvalue of $AB$ is an eigenvalue of $BA$.
\item Every eigenvector of $AB$ is an eigenvector of $BA$.
\end{enumerate}
Which of the following is correct?
\begin{enumerate}[label=(\alph*)]
\item Both (i) and (ii) are true.
\item (i) is true, (ii) is false.
\item (i) is false, (ii) is true.
\item Both (i) and (ii) are false.
\end{enumerate}

\item For a real number $a$, let $\lfloor a\rfloor$ denote the greatest integer less than or equal to $a$, and let $\{a\}=a-\lfloor a\rfloor$. Let $n\ge4$ be an integer. Then the expression
\[ \left\lfloor\frac{1}{\{n!\,e\}}\right\rfloor \]
equals
\begin{enumerate}[label=(\alph*)]
\item $n$.
\item $n-1$.
\item $\lfloor(n-2)e\rfloor$.
\item None of the other options.
\end{enumerate}


\item Consider the space $X=\left\{A\in M_2(\mathbb{R})\ \Big|\ AA^{\mathsf J}=\begin{pmatrix}2&-2\\-2&3\end{pmatrix}\right\}$. Then the number of connected components of $X$ equals
\begin{enumerate}[label=(\alph*)]
\item $0$. \item $1$. \item $2$. \item $4$.
\end{enumerate}

\item Let $V$ be the real vector space of all polynomials in $\mathbb{R}[x]$ of degree less than or equal to $10$. Consider the linear operator $T:V\to V$ defined by $(Tp)(x)=p(2x+1)$ for all $p\in V$. What is the trace of $T$?
\begin{enumerate}[label=(\alph*)]
\item $11$. \item $2047$. \item $121$. \item None of the other options.
\end{enumerate}

\item Let $1\le p<q$ be integers, and let $f(x,y)=y^p-x^q$. Consider the complex vector space
\[ V=\frac{\mathbb{C}[x,y]}{\bigl(f,\ \frac{\partial f}{\partial x},\ \frac{\partial f}{\partial y}\bigr)}. \]
Then the dimension of $V$ over $\mathbb{C}$ is:
\begin{enumerate}[label=(\alph*)]
\item $pq$. \item $pq-p-q+1$. \item $p+q$. \item Infinite.
\end{enumerate}

\item Let $\Gamma\subset\mathbb{R}^2$ be the graph of $y=\sin(x)$ and let $\ell$ be a line in $\mathbb{R}^2$ such that $\#(\Gamma\cap\ell)=n\in\mathbb{N}$. Which of the following statements is necessarily true?
\begin{enumerate}[label=(\alph*)]
\item $n$ is even.
\item $n$ is odd.
\item $n\le10$.
\item None of the remaining three options.
\end{enumerate}

\item Suppose that $f:(-1,1)\to\mathbb{R}$ is infinitely differentiable, and $|f(x)-1-x-x^2|\le|x|^3$ for all $x\in(-\tfrac13,\tfrac13)$. Then,
\begin{enumerate}[label=(\alph*)]
\item $f(0)=f'(0)=f''(0)=1$.
\item $f(0)=f'(0)=1$ and $f'''(0)=6$.
\item $f(0)=f'(0)=1$ and $f''(0)=2$.
\item None of the other statements is true.
\end{enumerate}

\item Consider the following two assertions about a group $G$.
\begin{enumerate}[label=(\Roman*)]
\item $G$ is abelian.
\item The map $\phi:G\to G$ defined by $\phi(g)=g^2$ is a group homomorphism.
\end{enumerate}
Which of the following is true?
\begin{enumerate}[label=(\alph*)]
\item (I) implies (II), and (II) implies (I).
\item (I) implies (II), but (II) does not imply (I).
\item (I) does not imply (II), but (II) implies (I).
\item (I) does not imply (II), and (II) does not imply (I).
\end{enumerate}

\item Suppose that $f_n:[0,1]\to[0,1]$ is continuous and increasing for each $n\ge1$, and the sequence $(f_n(x))_{n\ge1}$ of real numbers is decreasing for all $x\in[0,1]$. Let $f:[0,1]\to[0,1]$ be the function defined as $f(x)=\lim_{n\to\infty}f_n(x)$. Then
\begin{enumerate}[label=(\alph*)]
\item $f$ is continuous.
\item $f$ is right-continuous but not necessarily left-continuous.
\item $f$ is left-continuous but not necessarily right-continuous.
\item None of the other assertions is true.
\end{enumerate}
\end{enumerate}

\subsection*{Part B --- True/False}
\begin{enumerate}[resume]
\item If $f:\mathbb{R}\to\mathbb{R}$ is a continuously differentiable function and $S\subset\mathbb{R}$ is a bounded subset, then there exists a constant $C_S>0$ such that $|f(x)-f(y)|\le C_S|x-y|$ for all $x,y\in S$.

\item There exists a continuous surjection from $\mathbb{Q}$ to $\{0,1,\ldots,2026\}$.

\item The subset $X=\{f\in C([0,1],\mathbb{R}):\ |f(x)-f(y)|\le|x-y|^{1/2}\ \text{for all } x,y\in[0,1]\}$ of $C([0,1],\mathbb{R})$ is compact.

\item There is no continuous bijection $f:\mathbb{R}\to[0,1]$.

\item Let $A,B$ be linear endomorphisms of a finite dimensional vector space over a field. If $A$ and $B$ are diagonalizable, then $AB$ is also diagonalizable.

\item If $A\in GL_5(\mathbb{C})$ and $A^2$ is diagonalizable, then $A$ is diagonalizable.

\item A function $f:\mathbb{R}\to\mathbb{R}$ is said to be periodic if there exists a positive real number $c$ such that $f(x+c)=f(x)$ for all $x\in\mathbb{R}$. Then the set of all periodic functions is a linear subspace of the real vector space of all functions from $\mathbb{R}$ to $\mathbb{R}$.

\item If $p(x)\in\mathbb{C}[x]$ is a polynomial which is not equal to $x$, then the polynomial $p(x)-x$ divides the polynomial $p(p(x))-x$.

\item Let $A$ be a square matrix with non-negative real entries such that the sum of entries in each row of $A$ is $1$. Then each eigenvalue of $A$ has absolute value less than or equal to $1$.

\item Let $X$ be a metric space having exactly $4$ points. Then, for some positive integer $n$, there exists a distance preserving map $f:X\to\mathbb{R}^n$, where $\mathbb{R}^n$ is equipped with its Euclidean metric.

\item Let $f:\mathbb{R}\to\mathbb{R}$ be a continuous function. For each positive integer $n$, let $f_n:\mathbb{R}\to\mathbb{R}$ be defined by $f_n(x)=f(x^n)$. If the sequence of functions $(f_n)$ converges uniformly on $\mathbb{R}$, then $f$ must be a constant function.

\item Let $X,Y$ be finite metric spaces. Then $X$ and $Y$ are homeomorphic if and only if the real vector spaces $C(X,\mathbb{R})$ and $C(Y,\mathbb{R})$ are isomorphic.

\item Suppose a finite group $G$ is the union of its Sylow subgroups. Then the order of $G$ is a power of some prime number.

\item For a $2\times2$ real matrix $A$, if the limit
\[ \sin(A)=\lim_{N\to\infty}\sum_{n=1}^N(-1)^{n-1}\frac{A^{2n-1}}{(2n-1)!} \]
exists in $M_2(\mathbb{R})$, then all eigenvalues of the limit matrix $\sin(A)$ have absolute value less than or equal to $1$.

\item Let $A$ be a finite abelian group of order $2025$. Then for every subgroup $H\le A$ we have $A\cong H\times A/H$ as groups.

\item Consider the map $L:\mathbb{R}[t]\to\mathbb{R}[t]$ given by $L(f)=3f+2f'$, where $f'$ denotes the derivative (with respect to $t$) of the polynomial $f$. Then $L$ is injective, but not surjective.

\item Let $\sigma_1,\ldots,\sigma_{80}$ be elements of the permutation group $S_{236}$ such that $\sigma_i^3$ is the identity permutation for each $i$. Then there exist $1\le i<j\le80$ such that $\sigma_i$ is conjugate to $\sigma_j$.

\item For a permutation $\sigma\in S_{10}$, let $\mathrm{Fix}(\sigma)\subseteq\{1,2,\ldots,10\}$ denote the set of fixed points of $\sigma$. Then the average $\dfrac1{10!}\displaystyle\sum_{\sigma\in S_{10}}\#\mathrm{Fix}(\sigma)$ is equal to $1$.

\item The set $X=\{(b,c)\in\mathbb{R}^2\mid\text{the polynomial }x^2+bx+c\text{ has distinct roots in }\mathbb{C}\}$ is connected.

\item The integer $1^{6n}+2^{6n}+3^{6n}+\cdots+2025^{6n}$ is divisible by $7$ for all integers $n\ge10$.
\end{enumerate}


\end{document}
